\documentclass[12pt,a4paper]{article}

% --- 1. Cấu hình Tiếng Việt và Font ---
\usepackage[utf8]{inputenc}
\usepackage[T5]{fontenc}           
\usepackage[vietnamese]{babel}
\usepackage{mathptmx} % Font Times New Roman tương đương

% --- 2. Xử lý màu sắc và bảng (Sửa lỗi Option Clash ở đây) ---
\usepackage[table,svgnames]{xcolor} 
\usepackage{booktabs}
\usepackage{multirow}
\usepackage{colortbl}
\usepackage{tabularx}
\usepackage{array}
\usepackage{float}

% --- 3. Đồ họa và Hình ảnh ---
\usepackage{graphicx}
\usepackage{tikz}
\usetikzlibrary{calc,positioning}
\usepackage{caption}
\captionsetup[figure]{labelformat=empty}

% --- 4. Định dạng trang ---
\usepackage{geometry}
\geometry{left=2.5cm,right=2.5cm,top=2.5cm,bottom=2.5cm}
\usepackage{parskip}           
\setlength{\parindent}{0pt}

% --- 5. Code Listings (Sửa lỗi khai báo trùng) ---
\usepackage{listings}
\lstdefinestyle{pythonStyle}{
    language=Python,
    backgroundcolor=\color{gray!10},
    basicstyle=\ttfamily\footnotesize,
    keywordstyle=\color{blue}\bfseries,
    stringstyle=\color{teal},
    commentstyle=\color{green!50!black}\itshape,
    numbers=left,
    numberstyle=\tiny\color{gray},
    breaklines=true,
    frame=single,
    showstringspaces=false,
    tabsize=4,
    captionpos=b
}
\lstset{style=pythonStyle}

% --- 6. Toán học và Link ---
\usepackage{amsmath,amssymb,amsfonts}
\usepackage{url}
\usepackage[
    colorlinks=true,
    linkcolor=black,
    citecolor=green,
    urlcolor=cyan
]{hyperref}

\usepackage{pdfpages}

\begin{document}

% Dưới comment nì là một cái bìa khác xíu mí anh coi thử nha
\begin{titlepage}
\begin{tikzpicture}[remember picture,overlay]
  \draw[line width=1.5pt]
    ($(current page.north west)+(1.5cm,-1.5cm)$) rectangle
    ($(current page.south east)+(-1.5cm,1.5cm)$);
\end{tikzpicture}

\begin{center}

\vspace*{0.5cm}

{\Large \textbf{ĐẠI HỌC QUỐC GIA THÀNH PHỐ HỒ CHÍ MINH}}\\[0.4cm]
{\large \textbf{TRƯỜNG ĐẠI HỌC CÔNG NGHỆ THÔNG TIN}}\\[0.3cm]
{\large \textbf{KHOA KHOA HỌC MÁY TÍNH}}

\vspace{0.8cm}
\rule{0.65\textwidth}{1pt}

\vspace{1.3cm}

\includegraphics[width=0.32\linewidth]{Logo ĐH Công Nghệ Thông Tin - UIT.png}

\vspace{1.3cm}

{\LARGE \textbf{BÁO CÁO ĐỒ ÁN CUỐI KỲ}}\\[0.4cm]
{\Large \textbf{Môn: Nhập môn Công nghệ Phần mềm}}

\vspace{1cm}

{\Large \textbf{\textcolor{blue}{
Phát triển Website tư vấn và kiểm tra\\
thành phần mỹ phẩm - SKINMATE
}}}

\vspace{1.4cm}

\begin{tabular}{rl}
\Large \textbf{Giảng viên hướng dẫn:} & \Large Đỗ Văn Tiến \\[0.5cm]

\Large \textbf{Sinh viên thực hiện:} & \Large 24520119 - Nguyễn Thị Vân Anh \\
& \Large 24521766 - Lê Đan Thảo Tiên \\
& \Large 24520749 - Nguyễn Khang \\
\end{tabular}

\vfill

{\Large TP. Hồ Chí Minh, 2026}

\end{center}
\end{titlepage}


\tableofcontents

\section{Tổng quan đề tài}

\subsection{Giới thiệu đề tài}
Trong bối cảnh thị trường làm đẹp phát triển mạnh mẽ, việc chăm sóc da (\textit{skincare}) đã trở thành một nhu cầu thiết yếu hàng ngày, đặc biệt là đối với giới trẻ. Cùng với đó, xu hướng tiêu dùng thông minh đang lên ngôi; người dùng không chỉ quan tâm đến công dụng quảng cáo mà còn chú trọng vào bảng thành phần (\textit{INCI}) cấu tạo nên sản phẩm. Tuy nhiên, rào cản lớn nhất hiện nay là các thành phần hóa học thường mang danh pháp phức tạp, gây khó khăn cho người tiêu dùng phổ thông trong việc tra cứu và ghi nhớ. Việc lựa chọn mỹ phẩm sai lệch, chứa thành phần kỵ với cơ địa da không chỉ gây lãng phí tài chính mà còn dẫn đến các rủi ro sức khỏe như kích ứng, viêm mụn hay tổn thương da liễu.

Hiện tại, quá trình tự tra cứu thủ công trên Internet mất rất nhiều thời gian và thông tin thường thiếu tính cá nhân hóa. Xuất phát từ thực tiễn đó, đề tài "\textit{Phát triển Website tư vấn và kiểm tra thành phần mỹ phẩm - SKINMATE}" được thực hiện nhằm mang đến một giải pháp công nghệ tiện lợi. Hệ thống cho phép người dùng cá nhân hóa hồ sơ làn da và tự động phân tích tính an toàn của các thành phần mỹ phẩm dựa trên tập luật logic (Rule-based), từ đó hỗ trợ người dùng đưa ra quyết định mua sắm an toàn và khoa học hơn.
\subsection{Mục tiêu phần mềm}
Phần mềm được xây dựng hướng tới các mục tiêu chính sau đây:

Mục tiêu về mặt ứng dụng: Xây dựng một nền tảng Web trực quan, thân thiện, cho phép người dùng tra cứu nhanh chóng tính chất của các thành phần mỹ phẩm bằng cách nhập chuỗi INCI hoặc trích xuất từ ảnh bảng thành phần bằng OCR. Hệ thống cung cấp các cảnh báo mức độ phù hợp (Tốt, Dễ kích ứng, Bình thường) thông qua hệ thống phân loại bằng màu sắc dễ hiểu.

Mục tiêu về mặt dữ liệu và logic: Xây dựng thành công cơ sở dữ liệu lưu trữ thông tin các chất hóa học phổ biến trong mỹ phẩm. Đồng thời, thiết lập được hệ thống đối chiếu dựa trên luật (Rule-based system) để tự động hóa việc ánh xạ giữa đặc tính thành phần và loại da cụ thể của từng người dùng.

Mục tiêu về mặt giá trị xã hội: Nâng cao nhận thức của người tiêu dùng về hóa mỹ phẩm, tạo dựng thói quen mua sắm an toàn, minh bạch và bảo vệ sức khỏe làn da cho cộng đồng.
\subsection{Phạm vi phần mềm}
Để đảm bảo tính khả thi và tập trung vào giải quyết vấn đề cốt lõi trong khuôn khổ thời gian của đồ án, phạm vi của phần mềm được giới hạn như sau:

Đối tượng phục vụ: Phần mềm hướng tới đối tượng là người tiêu dùng cá nhân, đặc biệt là học sinh, sinh viên có nhu cầu chăm sóc da cơ bản nhưng chưa có kiến thức chuyên sâu về da liễu.

Phạm vi chức năng:
\begin{itemize}
    \item Hệ thống tập trung vào việc quản lý tài khoản người dùng và thiết lập hồ sơ loại da (Da dầu, Da khô, Da nhạy cảm,...).
    \item Tính năng cốt lõi là phân tích chuỗi văn bản (text) chứa bảng thành phần mỹ phẩm do người dùng nhập vào, bóc tách và đối chiếu với cơ sở dữ liệu để đưa ra cảnh báo.
    \item Hỗ trợ Quản trị viên (Admin) cập nhật, quản lý danh sách thành phần và thiết lập các quy tắc phân tích (Rules).
    \item Hỗ trợ gợi ý sản phẩm trong danh mục do Admin quản lý dựa trên độ an toàn và sự tương thích với loại da người dùng.
    \item Hệ thống tích hợp AI (Gemini Flash, model \texttt{gemini-flash-latest}) đóng vai trò hỗ trợ phân tích các thành phần chưa có trong cơ sở dữ liệu và tự động cập nhật tri thức mới vào hệ thống (AI-powered Fallback).
\end{itemize}







Giới hạn của hệ thống (Out of scope): 
\begin{itemize}
    % \item Phần mềm sử dụng phương pháp đối chiếu logic theo luật (Rule-based matching) đã được định nghĩa sẵn trong cơ sở dữ liệu, không áp dụng các mô hình Trí tuệ nhân tạo (AI) phức tạp hay học máy (Machine Learning).
    \item Hệ thống hỗ trợ nhập chuỗi văn bản và trích xuất bảng thành phần từ ảnh bằng OCR; chưa hỗ trợ quét mã vạch sản phẩm.
\end{itemize}

\section{Xác định và thu thập yêu cầu}

\subsection{Kế hoạch và phương pháp khảo sát}

\subsubsection{Đối tượng và địa điểm khảo sát}

Các đối tượng chính trong khảo sát:
\begin{itemize}
    \item Học sinh, sinh viên và người trẻ (độ tuổi chủ yếu từ 16 - 25): Đây là nhóm khách hàng có nhu cầu chăm sóc da (skincare) ngày càng cao, nhạy bén với công nghệ và ứng dụng web. Tuy nhiên, nhóm này thường bị hạn chế về tài chính và chưa có nhiều kinh nghiệm chuyên sâu về da liễu, dễ mua sai sản phẩm nếu chỉ dựa vào quảng cáo.
    \item Người tiêu dùng từng gặp vấn đề về da (mụn, kích ứng, nhạy cảm): Nhóm đối tượng này có xu hướng mua sắm cẩn trọng, đặc biệt quan tâm đến bảng thành phần (INCI) của mỹ phẩm nhưng lại gặp khó khăn trong việc ghi nhớ các danh pháp hóa học phức tạp.
    \item Những người yêu thích làm đẹp: Những người thường xuyên có nhu cầu tra cứu, đối chiếu thành phần mỹ phẩm trước khi mua nhưng cảm thấy tốn quá nhiều thời gian với phương pháp tra cứu thủ công hiện tại.
\end{itemize}

\subsubsection{Các phương pháp khảo sát}

Để thu thập thông tin một cách chính xác, khách quan và đầy đủ nhất nhằm phục vụ cho việc xác định yêu cầu của hệ thống phần mềm SKINMATE, nhóm đã kết hợp sử dụng 4 phương pháp khảo sát thực tế sau đây:

\textbf{a. Phương pháp sử dụng bảng câu hỏi (Khảo sát trực tuyến - Survey):}

Cách thức thực hiện: Nhóm thiết kế một bảng câu hỏi trắc nghiệm ngắn gọn thông qua công cụ Google Forms và gửi đường dẫn khảo sát đến các hội nhóm sinh viên, diễn đàn mạng xã hội để thu thập ý kiến.

Mục đích: Thu thập số liệu định lượng trên diện rộng về thói quen mua sắm mỹ phẩm, tỷ lệ người dùng từng bị kích ứng da do dùng sai sản phẩm, và mức độ quan tâm đến bảng thành phần INCI.

Kết quả đạt được: Cung cấp các số liệu thống kê trực quan (biểu đồ) làm minh chứng thực tế cho tính cấp thiết của phần mềm. Đồng thời, kết quả khảo sát giúp nhóm xác định được tính năng cốt lõi (Core feature) mà người dùng mong muốn nhất là phân tích thành phần tự động và dịch ra tiếng Việt.

\textbf{b. Phương pháp phỏng vấn (Interview):}

Cách thức thực hiện: Chuẩn bị sẵn một bộ câu hỏi mở và tiến hành phỏng vấn ngắn trực tiếp một nhóm nhỏ (khoảng 5 - 7 người) là các sinh viên có thói quen sử dụng mỹ phẩm chăm sóc da thường xuyên.

Mục đích: Tìm hiểu sâu hơn về những khó khăn thực sự (pain points) của người dùng khi phải tự tra cứu thành phần hóa học một cách thủ công trên Internet.

Kết quả đạt được: Thu nhận được những đóng góp chi tiết về mặt giao diện và trải nghiệm. Người dùng cho biết họ cần một kết quả phân tích trực quan bằng màu sắc (Đỏ/Xanh) để dễ hiểu ngay lập tức thay vì phải đọc những đoạn giải thích học thuật dài dòng.

\textbf{c. Phương pháp quan sát thực tế (Observation):}

Cách thức thực hiện: Nhóm tiến hành quan sát hành vi của người dùng tại các cửa hàng mỹ phẩm và đặc biệt là trên các hội nhóm làm đẹp lớn (như Facebook "Ghiền Skincare", "Đẹp Chanh Sả").

Mục đích: Xem xét quy trình nhờ tư vấn thành phần hiện tại đang diễn ra như thế nào. Nhóm ghi nhận được việc người dùng thường xuyên chụp ảnh bảng thành phần và đăng lên mạng nhờ cộng đồng mạng phân tích giúp.

Kết quả đạt được: Khẳng định nhược điểm của quy trình hiện tại là tính thủ công, tốn thời gian chờ đợi phản hồi và thiếu độ tin cậy khoa học. Điều này làm nổi bật ưu điểm của hệ thống Rule-based (xử lý logic tức thời) mà nhóm đang xây dựng.

\textbf{d. Phương pháp nghiên cứu tài liệu và phân tích sản phẩm tương tự:}

Cách thức thực hiện: Trải nghiệm thực tế và phân tích các nền tảng tra cứu thành phần quốc tế đang có trên thị trường như CosDNA, INCIDecoder.

Mục đích: Học hỏi cách thức tổ chức cơ sở dữ liệu hóa học và nhận diện những điểm yếu của các hệ thống cũ.

Kết quả đạt được: Nhận thấy các website hiện tại có rào cản lớn về ngôn ngữ (chủ yếu là tiếng Anh), giao diện phức tạp và thiếu tính "Cá nhân hóa theo loại da". Từ đó, nhóm quyết định đưa tính năng thiết lập hồ sơ loại da vào SKINMATE làm thế mạnh cạnh tranh.

\subsection{Khảo sát hiện trạng}

\subsubsection{Hiện trạng tổ chức}
Hiện nay, quy trình tìm hiểu, tư vấn và lựa chọn mỹ phẩm chăm sóc da (skincare) của người tiêu dùng tại Việt Nam chưa được chuẩn hóa hay quản lý bởi một hệ thống tập trung. Vòng lặp mua sắm này đang diễn ra phân tán và phụ thuộc vào sự tương tác của 4 nhóm đối tượng chính sau đây:
\begin{itemize}
    \item Người tiêu dùng cá nhân (Đặc biệt là học sinh, sinh viên):
Đây là nhóm đối tượng đóng vai trò trung tâm, có nhu cầu thường xuyên trong việc tìm kiếm các sản phẩm làm đẹp phù hợp. Tuy nhiên, họ lại là nhóm thiếu hụt nhiều nhất về mặt kiến thức chuyên môn (da liễu, hóa học). Người tiêu dùng đang phải tự đóng vai trò là chuyên gia: tự tìm kiếm thông tin, đọc bảng thành phần và đưa ra quyết định mua hàng. Do đó, họ là người chịu rủi ro cao nhất khi thông tin bị sai lệch hoặc không đầy đủ.

    \item Nhân viên tư vấn tại các cửa hàng và đại lý phân phối mỹ phẩm:
Nhóm này đóng vai trò cầu nối cung cấp thông tin sản phẩm trực tiếp đến khách hàng. Tuy nhiên, hoạt động tư vấn chủ yếu dựa trên tài liệu tiếp thị của nhãn hàng, thiếu sự phân tích chuyên sâu, độc lập về từng thành phần hóa học cụ thể (INCI). Hơn nữa, việc tư vấn đôi khi mang tính thiên vị nhằm mục đích thúc đẩy doanh số, chưa thực sự đặt sự an toàn và đặc tính cơ địa da của từng khách hàng lên hàng đầu.

    \item Cộng đồng làm đẹp mạng xã hội (KOLs, KOCs, các hội nhóm):
Với sự phát triển của mạng xã hội, nhóm này nổi lên như những "tư vấn viên không chính thức". Họ chia sẻ trải nghiệm cá nhân và đánh giá (review) sản phẩm. Mặc dù giúp lan tỏa kiến thức, nhưng mô hình hoạt động này tự phát, thông tin mang tính chủ quan. Một sản phẩm tốt với KOL này không có nghĩa là an toàn với mọi người tiêu dùng do khác biệt về loại da.

    \item Các nền tảng tra cứu quốc tế và từ điển hóa học:
Đây là những tổ chức lưu trữ dữ liệu chuyên sâu về mỹ phẩm (ví dụ: CosDNA, INCIDecoder). Mặc dù cung cấp thông tin rất khách quan, nhưng các hệ thống này hoạt động độc lập, không hỗ trợ tiếng Việt và thiếu đi tính năng "cá nhân hóa" (không cho phép người dùng lưu trữ loại da của bản thân để hệ thống tự động lọc các chất kỵ/phù hợp).
\end{itemize}









=> Đánh giá chung: Mô hình "tổ chức" thông tin chăm sóc da hiện tại đang bị phân mảnh và thiếu tính liên kết. Người tiêu dùng bị mắc kẹt giữa vô vàn luồng thông tin nhưng lại thiếu đi một "trợ lý trung gian" (một hệ thống độc lập, khách quan) có khả năng tự động đối chiếu các đặc tính khoa học của mỹ phẩm với tình trạng da cá nhân của họ một cách nhanh chóng. Điều này là tiền đề cấp thiết để xây dựng phần mềm SKINMATE.

\subsubsection{Hiện trạng nghiệp vụ}
Trong bối cảnh hiện nay, khi người tiêu dùng có nhu cầu lựa chọn và sử dụng mỹ phẩm an toàn, quy trình thực hiện thường mang tính thủ công, phân tán và thiếu sự hỗ trợ từ các hệ thống chuyên biệt. Quy trình nghiệp vụ thực tế có thể được mô tả qua các bước sau:

\textbf{Bước 1: Tiếp cận sản phẩm}

Người dùng bắt gặp sản phẩm mỹ phẩm thông qua các kênh như:
\begin{itemize}
    \item Mạng xã hội (Facebook, TikTok, Instagram)
    \item Sàn thương mại điện tử
    \item Đề xuất từ bạn bè hoặc KOLs
\end{itemize}

Tại bước này, quyết định ban đầu thường dựa trên thương hiệu, quảng cáo hoặc đánh giá cảm tính, chưa có sự phân tích thành phần cụ thể.

\textbf{Bước 2: Kiểm tra bảng thành phần (INCI)}

Người dùng đọc bảng thành phần được in trên bao bì hoặc mô tả sản phẩm. Tuy nhiên:

Các thành phần được trình bày bằng thuật ngữ chuyên ngành (tiếng Anh/Latin)
Danh sách dài, khó hiểu đối với người không có kiến thức chuyên môn

Điều này dẫn đến việc người dùng không thể tự đánh giá mức độ an toàn hoặc phù hợp với loại da của mình.

\textbf{Bước 3: Tra cứu từng thành phần}

Người dùng tiến hành tra cứu thủ công bằng cách:
\begin{itemize}
    \item Tìm kiếm từng thành phần trên Google
    \item Truy cập các website từ điển mỹ phẩm quốc tế
    \item Đọc blog, bài viết phân tích
\end{itemize}

Quá trình này tồn tại nhiều bất cập:

\begin{itemize}
    \item Tốn thời gian do phải tra cứu từng thành phần riêng lẻ
    \item Thông tin phân tán, không đồng nhất giữa các nguồn
    \item Khó xác định thông tin nào là chính xác hoặc đáng tin cậy
\end{itemize}

\textbf{Bước 4: Tham khảo ý kiến từ cộng đồng}

Sau khi có thông tin sơ bộ, người dùng thường:
\begin{itemize}
    \item Hỏi ý kiến trong các group làm đẹp
    \item Xem review từ KOLs/beauty bloggers
    \item Nhận tư vấn từ người bán
\end{itemize}

Hạn chế của bước này:
\begin{itemize}
    \item Thông tin mang tính chủ quan
    \item Bị ảnh hưởng bởi quảng cáo
    \item Không bảo đảm phù hợp tình trạng da của mình
\end{itemize}


\textbf{Bước 5: Tổng hợp và ra quyết định}

Người dùng tự tổng hợp các thông tin đã thu thập để đưa ra quyết định mua hoặc không mua sản phẩm.

\begin{itemize}
    \item Không có công cụ hỗ trợ tổng hợp dữ liệu một cách hệ thống
    \item Dễ xảy ra sai sót do thiếu kiến thức chuyên môn
    \item Quyết định mang tính cảm tính hoặc không chính xác.
\end{itemize}



\subsubsection{Hiện trạng tin học}

\textbf{Thiết bị sử dụng của người dùng:}  
Đối tượng người dùng mục tiêu của hệ thống SKINMATE chủ yếu là học sinh, sinh viên – nhóm có mức độ tiếp cận công nghệ cao. Phần lớn người dùng sở hữu điện thoại thông minh (smartphone) và/hoặc máy tính xách tay (laptop). Trong đó, smartphone là thiết bị phổ biến nhất, được sử dụng thường xuyên để truy cập internet và mạng xã hội. Laptop hỗ trợ tốt hơn cho các thao tác học tập và tra cứu thông tin chuyên sâu. Do đó, việc phát triển hệ thống dưới dạng ứng dụng web là phù hợp, cho phép truy cập đa nền tảng mà không cần cài đặt.

\textbf{Hạ tầng mạng và kết nối Internet:}  
Hạ tầng mạng tại Việt Nam hiện nay đã phát triển mạnh mẽ với sự phổ biến của Wi-Fi, 4G và 5G. Người dùng có thể truy cập internet tại nhiều địa điểm như nhà riêng, trường học hoặc các không gian công cộng. Tốc độ kết nối ổn định, đáp ứng tốt nhu cầu sử dụng các hệ thống web. Điều này đảm bảo SKINMATE có thể hoạt động liên tục và dễ dàng tiếp cận người dùng.

\textbf{Phần mềm và trình duyệt:}  
Người dùng chủ yếu sử dụng các trình duyệt phổ biến như Google Chrome, Microsoft Edge và Safari. Các trình duyệt này đều hỗ trợ tốt các công nghệ web hiện đại như HTML, CSS, JavaScript và các thư viện giao diện như Tailwind CSS. Vì vậy, hệ thống có thể triển khai bằng công nghệ đơn giản nhưng vẫn đảm bảo tính tương thích cao.

\textbf{Thói quen sử dụng công nghệ:}  
Người dùng có thói quen tìm kiếm thông tin qua Google, sử dụng các website để tra cứu và dễ dàng thực hiện các thao tác như copy/paste nội dung. Đồng thời, người dùng ưu tiên các hệ thống có giao diện đơn giản, dễ sử dụng và phản hồi nhanh. Điều này hoàn toàn phù hợp với mô hình hoạt động của SKINMATE.

\textbf{Đánh giá chung:}  
Hiện trạng tin học cho thấy không có rào cản lớn về thiết bị hay hạ tầng đối với việc triển khai hệ thống. Người dùng có đầy đủ điều kiện để tiếp cận và sử dụng một ứng dụng web đơn giản, từ đó tạo nền tảng thuận lợi cho việc phát triển và triển khai SKINMATE.


\subsection{Đánh giá hiện trạng và những tính năng mong muốn}

\subsubsection{Ưu điểm}

\begin{itemize}
    \item Người dùng có khả năng tiếp cận Internet và thiết bị công nghệ cao.
    \item Nhiều nguồn thông tin về thành phần mỹ phẩm sẵn có trên Internet.
    \item Người dùng có ý thức ngày càng cao về việc lựa chọn mỹ phẩm an toàn.
    \item Thói quen sử dụng công nghệ (tìm kiếm, tra cứu, copy/paste) phù hợp với hệ thống đề xuất.
\end{itemize}

\subsubsection{Nhược điểm và khó khăn cần giải quyết}

\begin{itemize}
    \item Thông tin về thành phần mỹ phẩm phân tán, không tập trung tại một nền tảng duy nhất.
    \item Người dùng phải tra cứu thủ công từng thành phần, gây mất thời gian và công sức.
    \item Nội dung mang tính chuyên ngành, khó hiểu đối với người không có kiến thức nền.
    \item Thiếu công cụ hỗ trợ đánh giá mức độ phù hợp của thành phần theo từng loại da.
    \item Thông tin từ cộng đồng hoặc KOLs mang tính chủ quan, thiếu độ tin cậy.
    \item Không có hệ thống trực quan giúp tổng hợp và hiển thị kết quả một cách dễ hiểu.
\end{itemize}

\subsubsection{Giải pháp và tính năng mong muốn}

Dựa trên các hạn chế của hệ thống hiện tại, đề tài đề xuất xây dựng hệ thống SKINMATE với các giải pháp và tính năng sau:

\begin{itemize}
    \item Xây dựng hệ thống tra cứu thành phần mỹ phẩm tự động dựa trên cơ chế rule-based.
    \item Cho phép người dùng nhập danh sách thành phần bằng cách dán (paste) văn bản.
    \item Tự động tách và phân tích từng thành phần trong chuỗi nhập liệu.
    \item Đối chiếu với cơ sở dữ liệu để đánh giá mức độ phù hợp theo loại da.
    \item Hiển thị kết quả trực quan bằng màu sắc:
    \begin{itemize}
        \item Màu đỏ: Thành phần không phù hợp / dễ gây kích ứng
        \item Màu xanh: Thành phần tốt / phù hợp
        \item Màu xám: Thành phần trung tính
    \end{itemize}
    \item Cho phép người dùng lưu thông tin cá nhân cơ bản (loại da) để cá nhân hóa kết quả.
    \item Giao diện đơn giản, dễ sử dụng, tối ưu cho cả thiết bị di động và máy tính.
\end{itemize}

\textbf{Kết luận:}
Các tính năng đề xuất tập trung giải quyết trực tiếp các vấn đề tồn tại trong quy trình hiện tại, hướng đến việc xây dựng một hệ thống đơn giản, hiệu quả, dễ sử dụng và phù hợp với phạm vi của đồ án môn học.


\section{Đặc tả và phân tích yêu cầu chức năng phần mềm}

\subsection{Yêu cầu chức năng}

\subsubsection{Yêu cầu chức năng nghiệp vụ}

\begin{center}
\begin{tabular}{|c|p{3cm}|p{5cm}|p{2.5cm}|p{3cm}|}
\hline
\textbf{STT} & \textbf{Nghiệp vụ} & \textbf{Thao tác cụ thể} & \textbf{Biểu mẫu} & \textbf{Quy định / Ghi chú} \\
\hline

1 & Đăng ký tài khoản 
& Người dùng nhập username, password và chọn loại da 
& Form đăng ký 
& Username không trùng, password tối thiểu 6 ký tự \\
\hline

2 & Đăng nhập 
& Người dùng nhập username và password 
& Form đăng nhập 
& Kiểm tra đúng thông tin tài khoản và trạng thái hoạt động \\
\hline

3 & Chọn loại da 
& Người dùng chọn loại da (Bình thường, Dầu, Khô, Nhạy cảm, Hỗn hợp) 
& Form hồ sơ cá nhân 
& Mỗi user chỉ có 1 loại da tại một thời điểm \\
\hline

4 & Nhập danh sách thành phần 
& Người dùng dán (paste) chuỗi thành phần cách nhau bằng dấu phẩy 
& Form nhập liệu 
& Chuỗi nhập phải có ít nhất 1 thành phần \\
\hline

5 & Xử lý chuỗi thành phần 
& Hệ thống tách chuỗi thành các thành phần riêng lẻ 
& Không có 
& Tách theo dấu phẩy, loại bỏ khoảng trắng dư \\
\hline

6 & Tra cứu thành phần 
& Hệ thống đối chiếu từng thành phần với database 
& Không có 
& So khớp theo tên ingredient (không phân biệt hoa thường) \\
\hline

7 & Phân tích AI dự phòng (Fallback)
& Hệ thống gọi Gemini API khi thành phần không có trong CSDL
& Không có 
& Tự động lưu kết quả AI vào CSDL để làm giàu tri thức \\
\hline

8 & Đánh giá thành phần 
& Xác định mức độ phù hợp theo loại da của user 
& Không có 
& Dựa trên bảng Ingredient\_Rules hoặc kết quả từ AI \\
\hline

9 & Hiển thị kết quả 
& Hiển thị danh sách thành phần kèm màu sắc 
& Giao diện kết quả 
& Đỏ: không phù hợp, Xanh: tốt, Xám: trung tính \\
\hline

10 & Xem mô tả thành phần 
& Người dùng xem mô tả chi tiết của từng chất 
& Popup / modal 
& Lấy dữ liệu từ bảng Ingredients \\
\hline

11 & Đề xuất sản phẩm 
& Hiển thị sản phẩm an toàn dựa trên loại da
& Giao diện kết quả 
& Loại bỏ sản phẩm chứa bất kỳ chất "Đỏ" (BAD) nào \\
\hline

12 & Quản lý lịch sử 
& Người dùng xem lại hoặc xóa các lần phân tích cũ
& Trang Lịch sử 
& Tự động lưu sau mỗi lần phân tích thành công \\
\hline

13 & Quản lý tài khoản (Admin)
& Admin khóa/mở khóa hoặc xóa tài khoản người dùng
& Quản lý User 
& Tài khoản bị khóa (Inactive) không thể đăng nhập \\
\hline

14 & Báo cáo thống kê (Admin)
& Xem biểu đồ phân bố loại da và tổng lượng phân tích
& Dashboard Admin 
& Trực quan hóa dữ liệu bằng biểu đồ (Recharts) \\
\hline

15 & Quản lý dữ liệu Excel (Admin)
& Admin nhập/xuất Ingredients, Rules, Products bằng Excel
& Import/Export 
& Hỗ trợ định dạng .xlsx, xử lý lỗi dữ liệu khi Import \\
\hline


\end{tabular}
\end{center}



\begin{center}
\begin{tabular}{|c|p{3cm}|p{5cm}|p{2.5cm}|p{3cm}|}
\hline
\textbf{STT} & \textbf{Nghiệp vụ} & \textbf{Thao tác cụ thể} & \textbf{Biểu mẫu} & \textbf{Quy định / Ghi chú} \\
\hline
% Thêm vào cuối bảng nghiệp vụ (trước \end{tabular})
16 & Trích xuất INCI bằng ảnh (OCR) 
& Người dùng tải ảnh chụp bảng thành phần, hệ thống tự động nhận dạng ký tự và điền chuỗi INCI 
& Form phân tích (có vùng tải ảnh) 
& Ảnh định dạng JPG, PNG; xử lý qua OCR.space API \\
\hline
17 & Báo cáo phân loại sai thành phần 
& Người dùng đề xuất phân loại lại một thành phần (GOOD/BAD/NEUTRAL) kèm lý do và dẫn chứng 
& Form báo cáo (BM8) 
& Lý do $\ge$ 20 ký tự, link dẫn chứng tùy chọn; báo cáo ở trạng thái PENDING cho Admin duyệt \\
\hline
18 & Bình chọn báo cáo cộng đồng 
& Người dùng bình chọn (UP/DOWN) cho các báo cáo đang chờ duyệt 
& Không có biểu mẫu (nút bấm) 
& Mỗi báo cáo chỉ được bình chọn 1 lần, có thể thay đổi hoặc hủy bình chọn \\
\hline
19 & Xem thông báo 
& Người dùng xem danh sách thông báo, đánh dấu đã đọc từng thông báo hoặc tất cả 
& Không có biểu mẫu (dropdown) 
& Thông báo gồm kết quả duyệt báo cáo, tin nhắn từ Admin; tải lại khi mở danh sách và định kỳ mỗi 30 giây khi danh sách đang đóng \\
\hline
20 & Duyệt báo cáo cộng đồng (Admin) 
& Admin xét duyệt báo cáo (APPROVED/REJECTED), nếu duyệt → tự động cập nhật IngredientRule 
& Form duyệt (modal) 
& Kèm ghi chú Admin; gửi thông báo tới người báo cáo \\
\hline
21 & Gửi thông báo đến người dùng (Admin) 
& Admin gửi thông báo tùy chỉnh đến một người dùng cụ thể 
& API (có thể mở rộng giao diện) 
& Endpoint POST /api/v1/notifications/send \\
\hline


\end{tabular}
\end{center}
%====================================================

\subsubsection{Danh sách các Biểu mẫu và Quy định}

\textbf{Danh sách biểu mẫu (BM):}

\begin{itemize}
    \item \textbf{BM 1:} Form đăng ký tài khoản (Thông tin cơ bản + Loại da).
    \item \textbf{BM 2:} Form đăng nhập (Hệ thống cấp JWT sau khi xác thực thông tin đăng nhập).
    \item \textbf{BM 3:} Form cập nhật hồ sơ cá nhân (Quản lý loại da, tên hiển thị, đổi mật khẩu).
    \item \textbf{BM 4:} Form phân tích thành phần INCI (Có tích hợp tải ảnh OCR, hiển thị kết quả màu sắc và gợi ý sản phẩm).
    \item \textbf{BM 5:} Form quản trị dữ liệu tập trung dành cho Admin (Ingredients, Rules, Products) – hỗ trợ Thêm/Sửa/Xóa.
    \item \textbf{BM 6:} Bảng điều khiển người dùng dành cho Admin (Quản lý trạng thái tài khoản, khóa/mở khóa, xóa).
    \item \textbf{BM 7:} Giao diện Import / Export dữ liệu (Xử lý file Excel .xlsx, hiển thị báo cáo kết quả).
    \item \textbf{BM 8:} Form báo cáo phân loại sai thành phần (Hiển thị dưới dạng Modal, yêu cầu nhập lý do $\ge$ 20 ký tự, link dẫn chứng tùy chọn).
    \item \textbf{BM 9:} Modal duyệt báo cáo dành cho Admin (Cho phép nhập ghi chú và xác nhận APPROVE/REJECT).
\end{itemize}

\textbf{Quy định (QĐ):}

\begin{itemize}
    \item \textbf{Xác thực \& Phân quyền:} Người dùng phải đăng nhập để sử dụng tính năng Phân tích, Lịch sử, Báo cáo và Bình chọn. Chỉ Admin mới có quyền truy cập vào các route \texttt{/admin/*} thông qua \textit{adminMiddleware}.
    
    \item \textbf{Cá nhân hóa dữ liệu:} Loại da (\textit{skinType}) được sử dụng để chọn tập luật đối chiếu (\textit{IngredientRules}); biểu mẫu đăng ký yêu cầu lựa chọn loại da và CSDL đặt mặc định \texttt{NORMAL} khi không có giá trị.
    
    \item \textbf{Chuẩn hóa đầu vào INCI:} Danh sách thành phần phải phân tách bằng dấu phẩy (`,`). Hệ thống tự động loại bỏ khoảng trắng thừa và chuyển tên thành phần về chữ thường để đảm bảo khớp chính xác với cơ sở dữ liệu.
    
    \item \textbf{Cơ chế Hybrid Analysis (AI Fallback):}  
    - Ưu tiên tra cứu trong CSDL nội bộ.  
    - Nếu thành phần chưa có, hệ thống tự động gọi \textbf{Gemini Flash API} (model \texttt{gemini-flash-latest}) để phân tích dựa trên loại da của người dùng.  
    - Kết quả từ AI được tự động lưu (\textit{auto-caching}) vào bảng \texttt{Ingredient} và \texttt{IngredientRule} để tái sử dụng cho các lần sau.
    
    \item \textbf{An toàn sản phẩm (Safety-first):} Sản phẩm được gợi ý phải \textbf{tuyệt đối không chứa bất kỳ thành phần nào bị đánh giá là BAD} đối với loại da của người dùng. Trong nhóm sản phẩm an toàn, hệ thống chấm điểm theo số thành phần GOOD, lấy tối đa 6 sản phẩm điểm cao và xáo trộn để hiển thị tối đa 3 gợi ý.
    
    \item \textbf{Kiểm soát tài khoản:} Tài khoản có trường \textit{isActive = false} sẽ bị từ chối cấp JWT token ngay tại tầng xác thực (\textit{auth.service}) với thông báo “Tài khoản của bạn đã bị khóa bởi Quản trị viên”.
    
    \item \textbf{Quản lý dữ liệu quy mô lớn (Excel Import/Export):}  
    - Hỗ trợ tệp \texttt{.xlsx} và \texttt{.xls}.  
    - Cơ chế \textit{Upsert}: nếu bản ghi đã tồn tại (theo tên), hệ thống sẽ cập nhật; nếu chưa có thì tạo mới.  
    - Sau khi import, hệ thống hiển thị báo cáo chi tiết: số dòng tạo mới, số dòng cập nhật, số dòng lỗi kèm lý do.
    
    \item \textbf{Báo cáo phân loại sai từ cộng đồng:}  
    - Người dùng phải đăng nhập và có loại da trong hồ sơ.  
    - Lý do báo cáo phải có độ dài tối thiểu 20 ký tự.  
    - Báo cáo được lưu với trạng thái \texttt{PENDING}, chờ Admin xử lý.  
    - Mỗi báo cáo có thể được bình chọn UP/DOWN bởi nhiều người dùng, mỗi người chỉ được bình chọn một lần (có thể thay đổi hoặc hủy).
    
    \item \textbf{Duyệt báo cáo (Admin):}  
    - Khi Admin chọn \texttt{APPROVED}, hệ thống tự động tạo hoặc cập nhật \texttt{IngredientRule} tương ứng với \texttt{ingredientId}, \texttt{skinType} và \texttt{reportedEffect}.  
    - Sau khi duyệt (dù APPROVED hay REJECTED), hệ thống tự động gửi thông báo (\texttt{Notification}) đến người dùng đã gửi báo cáo; nội dung thông báo sử dụng ghi chú của Admin nếu có.  
    - Báo cáo chuyển sang trạng thái \texttt{APPROVED} hoặc \texttt{REJECTED} sẽ không còn hiển thị trong danh sách chờ.
    
    \item \textbf{Thông báo hệ thống:}  
    - Thông báo bao gồm các loại: \texttt{REPORT\_RESOLVED} (kết quả duyệt báo cáo), \texttt{ADMIN\_MESSAGE} (tin nhắn từ Admin).  
    - Thông báo được tải khi người dùng đăng nhập, khi mở dropdown và được polling mỗi 30 giây trong lúc dropdown đang đóng.  
    - Người dùng có thể đánh dấu đã đọc từng thông báo hoặc tất cả.
    
    \item \textbf{Rate Limit:} Người dùng chỉ được thực hiện tối đa \textbf{25 lần phân tích} trong vòng 24 giờ để tránh lạm dụng tài nguyên và API Gemini.
\end{itemize}





\subsubsection{Chi tiết các Biểu mẫu và Quy định}

Để làm rõ các yêu cầu nghiệp vụ, hệ thống quy định các Biểu mẫu (BM) nhập liệu và các Quy định (QĐ) ràng buộc đi kèm như sau:

\vspace{0.5cm}
\textbf{1. Biểu mẫu 1 và Quy định 1 (Đăng ký tài khoản)}

\textbf{* BM 1: Form đăng ký tài khoản}
\begin{center}
\begin{tabular}{|c|p{3.5cm}|p{3.5cm}|p{5cm}|}
\hline
\textbf{STT} & \textbf{Trường/Mục} & \textbf{Kiểu dữ liệu} & \textbf{Ghi chú} \\
\hline
1 & Username & Chuỗi (String) & Tên đăng nhập của người dùng. \\
\hline

2 & Họ và tên (Display name) & Chuỗi (String) & Tên hiển thị trên hệ thống (Tùy chọn). \\
\hline

3 & Password & Chuỗi (String) & Mật khẩu bảo mật. \\
\hline
4 & Loại da (Skin type) & Lựa chọn (Enum) & Gồm: Bình thường, Dầu, Khô, Nhạy cảm, Hỗn hợp. \\
\hline
\end{tabular}
\end{center}

\textbf{- QĐ 1: Ràng buộc đăng ký:}
\begin{itemize}
    \item Username là duy nhất, không được trùng lặp với tài khoản đã có trong hệ thống.
    \item Password phải có độ dài tối thiểu 6 ký tự để đảm bảo bảo mật cơ bản.
    \item Tất cả các trường (Username, Password, Loại da) đều bắt buộc, không được để trống.
    \item Display name: Có thể để trống (Nullable), nếu trống hệ thống sẽ mặc định hiển thị theo Username.
\end{itemize}

\vspace{0.5cm}
\textbf{2. Biểu mẫu 2 và Quy định 2 (Đăng nhập)}

\textbf{* BM 2: Form đăng nhập}
\begin{center}
\begin{tabular}{|c|p{3.5cm}|p{3.5cm}|p{5cm}|}
\hline
\textbf{STT} & \textbf{Trường/Mục} & \textbf{Kiểu dữ liệu} & \textbf{Ghi chú} \\
\hline
1 & Username & Chuỗi (String) & Tên đăng nhập đã đăng ký. \\
\hline
2 & Password & Chuỗi (String) & Mật khẩu tương ứng. \\
\hline
\end{tabular}
\end{center}

\textbf{- QĐ 2: Ràng buộc đăng nhập:}
\begin{itemize}
    \item Username phải tồn tại trong cơ sở dữ liệu.
    \item Password phải khớp với Username tương ứng. Hệ thống báo lỗi chung "Sai tên đăng nhập hoặc mật khẩu" nếu không khớp.
\end{itemize}

\vspace{0.5cm}
\textbf{3. Biểu mẫu 3 và Quy định 3 (Cập nhật hồ sơ loại da)}

\textbf{* BM 3: Form cập nhật loại da}
\begin{center}
\begin{tabular}{|c|p{3.5cm}|p{3.5cm}|p{5cm}|}
\hline
\textbf{STT} & \textbf{Trường/Mục} & \textbf{Kiểu dữ liệu} & \textbf{Ghi chú} \\
\hline
1 & Loại da (Skin type) & Lựa chọn (Enum) & Cho phép người dùng thay đổi tình trạng da hiện tại. \\
\hline
\end{tabular}
\end{center}

\textbf{- QĐ 3: Ràng buộc cập nhật:}
\begin{itemize}
    \item Mỗi tài khoản (User) chỉ được phép có 01 trạng thái loại da tại một thời điểm (để hệ thống áp dụng Rule ánh xạ chuẩn xác).
    \item Cập nhật loại da mới sẽ ghi đè lên loại da cũ trong CSDL.
\end{itemize}

\vspace{0.5cm}
\textbf{4. Biểu mẫu 4 và Quy định 4 (Nhập và Phân tích thành phần - Core Flow)}

\textbf{* BM 4: Form nhập danh sách thành phần (INCI)}
\begin{center}
\begin{tabular}{|c|p{3.5cm}|p{3.5cm}|p{5cm}|}
\hline
\textbf{STT} & \textbf{Trường/Mục} & \textbf{Kiểu dữ liệu} & \textbf{Ghi chú} \\
\hline
1 & Danh sách thành phần (inciString) & Chuỗi văn bản dài (Text) & Chuỗi chứa danh sách các danh pháp quốc tế INCI do người dùng dán (paste) vào. \\
\hline
\end{tabular}
\end{center}

\textbf{- QĐ 4: Ràng buộc phân tích và Gợi ý sản phẩm:}
\begin{itemize}
    \item \textbf{Điều kiện thực hiện:} Người dùng phải đăng nhập để phân tích. Loại da (\textit{skinType}) được dùng khi đối chiếu luật và gợi ý sản phẩm; endpoint gợi ý yêu cầu giá trị này có trong thông tin xác thực của người dùng.
    \item \textbf{Định dạng đầu vào:} Chuỗi nhập liệu (\textit{inciString}) không được để trống. Các thành phần bên trong phải phân tách bằng dấu phẩy (,). Hệ thống sử dụng phương thức \textit{.split(',').map(i => i.trim().toLowerCase())} để chuẩn hóa dữ liệu.
    \item \textbf{Cơ chế Phân tích Hybrid (Kết hợp Rule-based và AI Fallback):} 
    \begin{itemize}
        \item \textbf{Bước 1 (Local Lookup):} Hệ thống ưu tiên truy vấn các thành phần và quy tắc an toàn (\textit{IngredientRule}) có sẵn trong cơ sở dữ liệu PostgreSQL.
        \item \textbf{Bước 2 (AI Fallback):} Đối với các thành phần lạ (không tồn tại trong CSDL), hệ thống tự động gọi API \textbf{Gemini Flash} (model \texttt{gemini-flash-latest}) để phân tích dựa trên loại da của người dùng.
        \item \textbf{Bước 3 (Auto-caching):} Các kết quả từ AI sẽ được tự động lưu (\textit{upsert}) vào bảng \textit{Ingredient} và \textit{IngredientRule} để tái sử dụng cho các lần truy vấn sau, giúp tối ưu tốc độ và chi phí API.
        \item \textbf{Mặc định:} Chỉ khi cả CSDL và AI đều không phản hồi, thành phần mới được gán nhãn "Trung tính" (Màu xám).
    \end{itemize}
    \item \textbf{Cơ chế Gợi ý sản phẩm (Safety-first Filter):}
    \begin{itemize}
        \item Sau khi phân tích, hệ thống tự động kích hoạt bộ lọc sản phẩm. 
        \item \textbf{Quy tắc loại trừ:} Một sản phẩm chỉ được xét gợi ý nếu \textbf{không chứa bất kỳ thành phần nào} bị gắn nhãn \textbf{BAD} (Màu đỏ) đối với loại da hiện tại của người dùng.
        \item \textbf{Xếp hạng hiển thị:} Sản phẩm an toàn được chấm điểm theo số thành phần \textbf{GOOD}; hệ thống lấy nhóm tối đa 6 sản phẩm điểm cao nhất, xáo trộn và hiển thị tối đa 3 sản phẩm.
    \end{itemize}
    \item \textbf{Lưu trữ lịch sử:} Mỗi lần phân tích thành công, chuỗi \textit{rawInput} sẽ được tự động lưu vào bảng \textit{AnalysisHistory} gắn với \textit{userId} để người dùng có thể xem lại hoặc tái phân tích.
\end{itemize}




\vspace{0.5cm}
\textbf{5. Biểu mẫu 5 và Quy định 5 (Quản trị dữ liệu hệ thống - Admin CRUD)}

\textbf{* BM 5: Form quản lý Ingredients / Rules / Products}
\begin{center}
\begin{tabular}{|c|p{3.5cm}|p{3.5cm}|p{5cm}|}
\hline
\textbf{STT} & \textbf{Trường/Mục} & \textbf{Kiểu dữ liệu} & \textbf{Ghi chú} \\
\hline
1 & Tên thành phần / SP & Chuỗi (String) & Tên định danh duy nhất (Unique). \\
\hline
2 & Mô tả / Brand & Văn bản (Text) & Thông tin chi tiết hiển thị cho User. \\
\hline
3 & Tập luật (Effect) & Lựa chọn (Enum) & Gán nhãn: GOOD, BAD, NEUTRAL. \\
\hline
4 & INCI String & Văn bản (Text) & Danh sách thành phần của sản phẩm. \\
\hline
\end{tabular}
\end{center}

\textbf{- QĐ 5: Ràng buộc quản trị dữ liệu:}
\begin{itemize}
    \item \textbf{Phân quyền hệ thống:} Chỉ tài khoản có \textit{role: ADMIN} mới có quyền truy cập các đầu cuối (\textit{endpoints}) bắt đầu bằng \textit{/api/v1/admin/*}. Hệ thống kiểm tra qua \textit{adminMiddleware}.
    \item \textbf{Tính toàn vẹn dữ liệu:} Khi xóa một \textit{Ingredient}, hệ thống tự động xóa tất cả các \textit{IngredientRule} liên quan (Cơ chế \textit{Cascade Delete} trong Prisma).
    \item \textbf{Chuẩn hóa dữ liệu:} Tên thành phần khi lưu vào CSDL phải được tự động chuyển về chữ thường (\textit{lowercase}) để đảm bảo tính chính xác khi so khớp ở luồng Phân tích.
\end{itemize}

\vspace{0.5cm}
\textbf{6. Biểu mẫu 6 và Quy định 6 (Quản lý người dùng và Bảo mật)}

\textbf{* BM 6: Bảng điều khiển người dùng (User Management)}
\begin{center}
\begin{tabular}{|c|p{3.5cm}|p{3.5cm}|p{5cm}|}
\hline
\textbf{STT} & \textbf{Trường/Mục} & \textbf{Kiểu dữ liệu} & \textbf{Ghi chú} \\
\hline
1 & Trạng thái (isActive) & Boolean & Cho phép Khóa (Lock) hoặc Mở khóa. \\
\hline
2 & Quyền hạn (Role) & Enum & User hoặc Admin. \\
\hline
\end{tabular}
\end{center}

\textbf{- QĐ 6: Quy định về trạng thái tài khoản:}
\begin{itemize}
    \item \textbf{Cơ chế khóa tài khoản:} Admin có quyền chuyển trạng thái người dùng thành \textit{isActive: false}. 
    \item \textbf{Ràng buộc đăng nhập:} Người dùng có tài khoản bị khóa sẽ bị hệ thống từ chối cấp JWT Token tại tầng \textit{auth.service} với thông báo: "Tài khoản của bạn đã bị khóa bởi Quản trị viên."
    \item \textbf{Bảo vệ thao tác trên giao diện:} Màn hình quản lý không hiển thị nút Khóa/Mở khóa/Xóa đối với tài khoản có vai trò \texttt{ADMIN}; các request quản trị được bảo vệ bởi \textit{adminMiddleware}.
\end{itemize}

\vspace{0.5cm}
\textbf{7. Biểu mẫu 7 và Quy định 7 (Xử lý dữ liệu quy mô lớn - Excel Service)}

\textbf{* BM 7: Giao diện Import / Export Excel}
\begin{center}
\begin{tabular}{|c|p{3.5cm}|p{3.5cm}|p{5cm}|}
\hline
\textbf{STT} & \textbf{Trường/Mục} & \textbf{Kiểu dữ liệu} & \textbf{Ghi chú} \\
\hline
1 & Tệp dữ liệu & File (.xlsx) & File Excel chứa danh sách hàng loạt. \\
\hline
\end{tabular}
\end{center}

\textbf{- QĐ 7: Ràng buộc nhập/xuất dữ liệu:}
\begin{itemize}
    \item \textbf{Định dạng file:} Hệ thống chỉ chấp nhận tệp có định dạng \textit{.xlsx} hoặc \textit{.xls}.
    \item \textbf{Cơ chế Upsert:} Khi Import, nếu tên thành phần đã tồn tại, hệ thống sẽ thực hiện cập nhật (\textit{Update}) mô tả; nếu chưa có, hệ thống sẽ tạo mới (\textit{Create}).
    \item \textbf{Báo cáo lỗi:} Sau khi xử lý file, hệ thống phải trả về kết quả chi tiết: Số dòng thành công, số dòng bị lỗi và lý do lỗi cụ thể (ví dụ: Sai định dạng SkinType).
\end{itemize}

\vspace{0.5cm}
\textbf{8. Biểu mẫu 8 và Quy định 8 (Báo cáo phân loại sai thành phần)}

\textbf{* BM 8: Form báo cáo phân loại}
\begin{center}
\begin{tabular}{|c|p{3.5cm}|p{3.5cm}|p{5cm}|}
\hline
\textbf{STT} & \textbf{Trường/Mục} & \textbf{Kiểu dữ liệu} & \textbf{Ghi chú} \\
\hline
1 & Thành phần & Chuỗi (String) & Tên INCI đã được chuẩn hóa từ kết quả phân tích. \\
\hline
2 & Phân loại đề xuất & Lựa chọn (Enum) & GOOD, BAD, NEUTRAL – đánh giá mà người dùng cho là đúng. \\
\hline
3 & Lý do & Văn bản (Text) & Mô tả chi tiết lý do đề xuất thay đổi (tối thiểu 20 ký tự). \\
\hline
4 & Link dẫn chứng & URL (Tùy chọn) & Đường dẫn đến nghiên cứu, bài báo khoa học... \\
\hline
\end{tabular}
\end{center}

\textbf{- QĐ 8: Ràng buộc báo cáo:}
\begin{itemize}
    \item Người dùng phải đăng nhập và đã có loại da trong hồ sơ.
    \item Lý do báo cáo phải có ít nhất 20 ký tự để đảm bảo chất lượng đóng góp.
    \item Mỗi báo cáo được gắn với \texttt{skinType} của người dùng tại thời điểm báo cáo, vì ảnh hưởng đến luật phân tích tương ứng.
    \item Báo cáo được lưu với trạng thái \texttt{PENDING}, chờ Admin duyệt.
\end{itemize}

% \vspace{0.5cm}
% \textbf{9. Biểu mẫu 9 và Quy định 9 (Gửi thông báo từ Admin – Tùy chọn)}

% \textit{(Chức năng hiện có qua API, giao diện có thể bổ sung sau)}

% \textbf{- QĐ 9:} Admin phải cung cấp userId, tiêu đề và nội dung thông báo. Thông báo được lưu vào bảng \texttt{Notification} và hiển thị trong chuông thông báo của người dùng.


\vspace{0.5cm}
\textbf{* BM 9: Form duyệt báo cáo (Dành cho Quản trị viên)}
\begin{center}
\begin{tabular}{|c|p{3.5cm}|p{3.5cm}|p{5cm}|}
\hline
\textbf{STT} & \textbf{Trường/Mục} & \textbf{Kiểu dữ liệu} & \textbf{Ghi chú} \\
\hline
1 & Báo cáo cần duyệt & Chỉ đọc (Read-only) & Hiển thị thông tin: Thành phần, Người báo cáo, Loại da, Đề xuất, Lý do, Link dẫn chứng (nếu có). \\
\hline
2 & Quyết định & Lựa chọn (Radio/Button) & Gồm 2 lựa chọn: \texttt{APPROVED} (Duyệt) hoặc \texttt{REJECTED} (Từ chối). \\
\hline
3 & Ghi chú của Admin & Văn bản (Text), Tùy chọn & Nội dung giải thích lý do duyệt hoặc từ chối, sẽ dùng làm nội dung thông báo nếu được nhập. \\
\hline
\end{tabular}
\end{center}

\textbf{- QĐ 9: Ràng buộc khi duyệt báo cáo:}
\begin{itemize}
    \item Chỉ tài khoản có vai trò \texttt{ADMIN} mới được truy cập chức năng duyệt báo cáo (thông qua \textit{adminMiddleware}).
    \item Khi Admin chọn \texttt{APPROVED}, hệ thống tự động cập nhật (upsert) \texttt{IngredientRule} dựa trên cặp (\texttt{ingredientId}, \texttt{skinType}) với giá trị \texttt{reportedEffect} từ báo cáo.
    \item Sau khi duyệt (bất kỳ trạng thái nào), hệ thống bắt buộc tạo một bản ghi \texttt{Notification} gửi đến người dùng đã tạo báo cáo; trường nội dung sử dụng ghi chú của Admin nếu có, ngược lại dùng thông báo mặc định.
    \item Báo cáo sau khi được duyệt sẽ chuyển từ trạng thái \texttt{PENDING} sang \texttt{APPROVED} hoặc \texttt{REJECTED}, đồng thời cập nhật thời gian giải quyết (\texttt{resolvedAt}) và người giải quyết (\texttt{resolvedBy}).
    \item Báo cáo đã được duyệt sẽ không còn xuất hiện trong danh sách “chờ duyệt” trên giao diện của Admin.
\end{itemize}



\subsubsection{Bảng trách nhiệm yêu cầu nghiệp vụ}

Bảng dưới đây mô tả sự tương tác và phân chia trách nhiệm rõ ràng giữa Người dùng và Hệ thống phần mềm đối với các nghiệp vụ cốt lõi của SKINMATE:

\begin{center}
\begin{tabular}{|c|p{3.5cm}|p{4cm}|p{4.5cm}|p{2cm}|}
\hline
\textbf{STT} & \textbf{Tên yêu cầu} & \textbf{Người dùng} & \textbf{Phần mềm} & \textbf{Ghi chú} \\
\hline
1 & Quản lý tài khoản (Đăng ký / Đăng nhập) & Cung cấp thông tin tài khoản (Username, Password) theo BM 1 và BM 2. & Kiểm tra quy định (QĐ 1, QĐ 2). Nếu hợp lệ: Lưu trữ hoặc cấp quyền truy cập. Nếu vi phạm: Báo lỗi chi tiết. & Quá trình bắt buộc. \\
\hline
2 & Thiết lập hồ sơ loại da & Lựa chọn và cập nhật loại da cá nhân (Dầu, Khô,...) theo BM 3. & Ghi nhận thay đổi và lấy đây làm tiêu chí (điều kiện) để kích hoạt các luật phân tích. & \\
\hline
3 & Nhập liệu thành phần mỹ phẩm & Sao chép và dán (paste) chuỗi INCI cần kiểm tra theo BM 4. & Tiếp nhận văn bản, tự động tách chuỗi theo dấu phẩy và loại bỏ các khoảng trắng thừa. & Đầu vào là Text. \\
\hline
4 & Phân tích và đối chiếu dữ liệu & Yêu cầu hệ thống tiến hành kiểm tra. & Thực hiện truy vấn cơ sở dữ liệu, đối chiếu từng thành phần với tập luật (Ingredient\_Rules) khớp với loại da của user. & Xử lý Rule-based. \\
\hline
5 & Trả kết quả cảnh báo trực quan & Xem và tham khảo kết quả để ra quyết định mua sắm. & Trực quan hóa kết quả phân loại bằng màu sắc (Đỏ: Kích ứng, Xanh: Tốt, Xám: Trung tính) kèm mô tả ngắn gọn. & Tối ưu UI/UX. \\
\hline
6 & Gợi ý sản phẩm phù hợp & Xem danh sách sản phẩm an toàn được đề xuất để có thêm lựa chọn mua sắm. & Tự động truy vấn bảng Products dựa trên loại da, loại sản phẩm có luật \texttt{BAD}, tính điểm \texttt{GOOD}, xáo trộn trong nhóm điểm cao và hiển thị tối đa 3 sản phẩm. & Nối tiếp luồng phân tích. \\
\hline
7 & Trích xuất thành phần từ ảnh (OCR) & Tải ảnh chụp bảng thành phần & Gọi API OCR.space, xử lý kết quả và tự động điền vào ô INCI & Hỗ trợ nhập liệu nhanh hơn \\
\hline
8 & Báo cáo sai phân loại & Gửi đề xuất kèm lý do, dẫn chứng & Lưu báo cáo (PENDING), cập nhật danh sách báo cáo đang chờ & Đóng góp tri thức từ cộng đồng \\
\hline
9 & Bình chọn báo cáo & Bấm nút UP/DOWN trên báo cáo & Cập nhật số phiếu, hiển thị lại tổng điểm & Giúp Admin ưu tiên xử lý báo cáo được đồng thuận cao \\
\hline
\end{tabular}
\end{center}

\begin{center}
\renewcommand{\arraystretch}{1.3}
\begin{tabular}{|c|p{3.5cm}|p{4cm}|p{4.5cm}|p{2cm}|}
\hline
10 & Xem thông báo & Mở chuông thông báo & Lấy danh sách thông báo, cho phép đánh dấu đã đọc / xem chi tiết & Tăng tương tác và tính minh bạch của hệ thống \\
\hline
11 & Duyệt báo cáo (Admin) & Chọn chấp nhận hoặc từ chối & Nếu APPROVED: tự động tạo/cập nhật IngredientRule; gửi thông báo cho người báo cáo & Quản lý chất lượng dữ liệu \\
\hline
12 & Gửi thông báo (Admin) & Gửi dữ liệu thông báo qua API quản trị & Lưu vào Notification để người dùng xem trên chuông thông báo; hiện chưa có màn hình soạn tin riêng & Chức năng Backend \\
\hline
\end{tabular}
\end{center}





\subsubsection{Yêu cầu hiệu quả}
Yêu cầu hiệu quả liên quan đến tốc độ xử lý, thời gian phản hồi và khả năng tối ưu hóa tài nguyên. Với hệ thống SKINMATE, tính năng phân tích kết hợp AI đòi hỏi sự cân bằng giữa độ chính xác và tốc độ phản hồi để đảm bảo trải nghiệm người dùng không bị gián đoạn.

\textbf{- Bảng yêu cầu hiệu quả:}

\begin{center}
\renewcommand{\arraystretch}{1.3}
\begin{tabular}{|c|>{\raggedright\arraybackslash}p{3.5cm}|>{\raggedright\arraybackslash}p{4.5cm}|>{\centering\arraybackslash}p{2.5cm}|>{\centering\arraybackslash}p{2.5cm}|}
\hline
\textbf{STT} & \textbf{Nghiệp vụ} & \textbf{Tốc độ xử lý} & \textbf{Dung lượng lưu trữ} & \textbf{Ghi chú} \\ 
\hline

1 & Quản lý tài khoản (Đăng ký/Đăng nhập) 
& Phản hồi kết quả xác thực dưới 1.5 giây. 
& Không đáng kể. 
& Mã hóa Bcrypt cần thời gian xử lý. \\ 
\hline

2 & Phân tích chuỗi thành phần (Core flow) 
& Dưới 2 giây (với dữ liệu nội bộ) hoặc 5-8 giây (nếu gọi AI Gemini). 
& Tối ưu hóa Cache dữ liệu AI. 
& Phụ thuộc vào độ trễ của API Gemini. \\ 
\hline

3 & Gợi ý sản phẩm an toàn 
& Truy vấn và lọc loại trừ chất BAD dưới 2 giây. 
& Tối ưu hóa câu lệnh Prisma. 
& Thực hiện song song với phân tích thành phần. \\ 
\hline

4 & Báo cáo thống kê (Admin) 
& Tổng hợp dữ liệu và vẽ biểu đồ dưới 3 giây. 
& Tính toán Aggregation trên Database. 
& Xử lý dữ liệu lớn trên Recharts. \\ 
\hline

5 & Quản lý dữ liệu Excel (Import/Export) 
& Xử lý tệp tin 1000 dòng dưới 10 giây. 
& Tối ưu hóa bộ nhớ đệm (Memory Storage). 
& Phụ thuộc vào số lượng thành phần mới cần tạo. \\ 
\hline
\end{tabular}
\end{center}

\textbf{- Bảng trách nhiệm yêu cầu hiệu quả:}

\begin{center}
\renewcommand{\arraystretch}{1.3}
\begin{tabular}{|c|>{\raggedright\arraybackslash}p{3.5cm}|>{\raggedright\arraybackslash}p{3.5cm}|>{\raggedright\arraybackslash}p{4.5cm}|>{\raggedright\arraybackslash}p{2cm}|}
\hline
\textbf{STT} 
& \multicolumn{1}{c|}{\textbf{Nghiệp vụ}} 
& \multicolumn{1}{c|}{\textbf{Người dùng}} 
& \multicolumn{1}{c|}{\textbf{Phần mềm}} 
& \multicolumn{1}{c|}{\textbf{Ghi chú}} \\
\hline

1 & Phân tích thành phần Hybrid (CSDL + AI) 
& Cung cấp danh sách INCI đúng định dạng dấu phẩy. 
& Tự động nhận diện chất lạ, gọi AI và thực hiện ghi đệm (caching) vào DB để tăng tốc cho lần sau. 
& Xử lý thông minh. \\
\hline

2 & Đề xuất sản phẩm an toàn 
& Không cần thao tác thêm. 
& Sử dụng thuật toán lọc loại trừ chính xác, đảm bảo không có sản phẩm vi phạm được hiển thị. 
& Safety-first logic. \\
\hline

3 & Xử lý tệp Excel (Admin) 
& Cung cấp file đúng định dạng mẫu (.xlsx). 
& Thư viện \textit{xlsx} bóc tách dữ liệu nhanh, cơ chế \textit{upsert} giảm thiểu xung đột dữ liệu. 
& Xử lý hàng loạt. \\
\hline

4 & Báo cáo Dashboard 
& Truy cập trang Reports. 
& Truy vấn gộp nhóm (\textit{groupBy}) hiệu quả, hiển thị trực quan thông qua Recharts. 
& Trực quan hóa dữ liệu. \\
\hline
\end{tabular}
\end{center}

\subsubsection{Yêu cầu tiện dụng}
Yêu cầu tiện dụng đảm bảo phần mềm có giao diện trực quan, sang trọng (Luxury Design System), phù hợp với đối tượng trẻ - những người ưu tiên trải nghiệm nhanh và thẩm mỹ. Hệ thống tối ưu hóa UX thông qua cơ chế phân loại màu sắc Rose-Sage và các thao tác "một chạm".

\textbf{- Bảng yêu cầu tiện dụng:}

\begin{center}
\renewcommand{\arraystretch}{1.3}
\begin{tabular}{|c|p{3.5cm}|p{3cm}|p{4.5cm}|p{2.5cm}|}
\hline
\textbf{STT} & \textbf{Nghiệp vụ} & \textbf{Mức độ dễ học} & \textbf{Mức độ dễ sử dụng} & \textbf{Ghi chú} \\
\hline
1 & Quản lý tài khoản & Không cần hướng dẫn. & Biểu mẫu Glassmorphism hiện đại, thông báo lỗi (Validation) hiển thị tức thì. & Giao diện Luxury. \\
\hline
2 & Phân tích thành phần & 1 phút làm quen. & Chỉ cần dán (Paste) văn bản. Kết quả màu sắc trực quan: Rose (Bad), Sage (Good), Gray (Neutral). & Tối ưu hóa thị giác. \\
\hline
3 & Quản lý lịch sử & Không cần hướng dẫn. & Xem lại và tái phân tích (Re-analyze) chỉ với một nút bấm. & Tiết kiệm thời gian. \\
\hline
4 & Đề xuất sản phẩm & Không cần hướng dẫn. & Sản phẩm an toàn tự động hiển thị dưới dạng thẻ (Cards) sinh động, dễ tương tác. & UX mượt mà. \\
\hline
5 & Quản trị hệ thống (Admin) & 5 phút làm quen. & Dashboard có biểu đồ trực quan, hỗ trợ quản lý hàng loạt qua Excel thay vì nhập tay. & Hiệu quả quản lý. \\
\hline
\end{tabular}
\end{center}

\textbf{- Bảng trách nhiệm yêu cầu tiện dụng:}

\begin{center}
\renewcommand{\arraystretch}{1.3}
\begin{tabular}{|c|p{3.5cm}|p{3.5cm}|p{4.5cm}|p{2cm}|}
\hline
\textbf{STT} & \textbf{Nghiệp vụ} & \textbf{Người dùng} & \textbf{Phần mềm} & \textbf{Ghi chú} \\
\hline
1 & Trải nghiệm người dùng tổng thể & Sử dụng trên đa thiết bị (Mobile/Laptop). & Sử dụng \textbf{Tailwind CSS v4} tạo giao diện Responsive. Áp dụng font Playfair Display và Inter cho cảm giác sang trọng. & Thẩm mỹ cao. \\
\hline
2 & Phân tích & Cung cấp chuỗi INCI. & Tự động xử lý logic phức tạp phía Backend, trả về kết quả đơn giản, dễ hiểu nhất cho User. & Giảm tải nhận thức. \\
\hline
3 & Quản trị dữ liệu & Admin cung cấp file Excel mẫu. & Xử lý tự động việc đối chiếu và lưu trữ, cung cấp báo cáo biểu đồ qua Recharts để Admin dễ nắm bắt. & Tối ưu quản trị. \\
\hline
\end{tabular}
\end{center}


\subsubsection{Yêu cầu tương thích}
Yêu cầu tương thích mô tả khả năng hoạt động của phần mềm trên các nền tảng, môi trường phần mềm hoặc thiết bị phần cứng khác nhau. Do hệ thống SKINMATE hướng tới đối tượng người dùng trẻ thường xuyên sử dụng điện thoại di động, tính tương thích đa thiết bị thông qua nền tảng Web là yếu tố được ưu tiên hàng đầu.

\textbf{- Bảng yêu cầu tương thích:}

\begin{center}
\begin{tabular}{|c|p{4cm}|p{5cm}|p{3.5cm}|}
\hline
\textbf{STT} & \textbf{Nghiệp vụ / Chức năng} & \textbf{Đối tượng liên quan (Môi trường)} & \textbf{Ghi chú} \\
\hline
1 & Truy cập và sử dụng toàn bộ hệ thống & Các trình duyệt Web hiện đại (Google Chrome, Safari, Microsoft Edge, Firefox). & Tương thích tốt nhất trên các phiên bản trình duyệt mới. \\
\hline
2 & Hiển thị giao diện và thao tác & Thiết bị di động (Smartphone), Máy tính bảng (Tablet) và Máy tính cá nhân (Laptop/PC). & Giao diện hiển thị co giãn (Responsive) linh hoạt nhờ thư viện Tailwind CSS v4. \\
\hline
\end{tabular}
\end{center}


\subsubsection{Yêu cầu tiến hoá}
Yêu cầu tiến hóa mô tả khả năng thích nghi và thay đổi các quy định, tham số nghiệp vụ của phần mềm trong tương lai mà không cần can thiệp vào mã nguồn. Đối với SKINMATE, yêu cầu này đảm bảo hệ thống không bị lạc hậu trước sự ra đời liên tục của các hoạt chất mỹ phẩm mới và sự thay đổi trong các nghiên cứu da liễu.

\textbf{- Bảng yêu cầu tiến hoá:}

\begin{center}
\renewcommand{\arraystretch}{1.3}
\begin{tabular}{|c|p{3.5cm}|p{4.5cm}|p{4cm}|}
\hline
\rowcolor{blue!10}
\textbf{STT} & \textbf{Nghiệp vụ} & \textbf{Tham số cần thay đổi} & \textbf{Miền giá trị cần thay đổi} \\
\hline
1 & Thay đổi luật phân tích thành phần & Đánh giá an toàn (\textit{Effect}) & Bảng \texttt{Ingredient\_Rule} \\
\hline
2 & Cập nhật tri thức hóa mỹ phẩm & Tên hoạt chất (\textit{INCI}) và mô tả khoa học & Bảng \texttt{Ingredient} \\
\hline
3 & Mở rộng danh mục gợi ý & Thông tin sản phẩm, thương hiệu và bảng thành phần & Bảng \texttt{Product} \\
\hline
4 & Tự động hóa tri thức mới & Kết quả phân tích hoạt chất lạ từ AI & Tích hợp \textbf{Gemini AI API} \\
\hline
\end{tabular}
\end{center}

\textbf{- Bảng trách nhiệm yêu cầu tiến hoá:}

\begin{center}
\renewcommand{\arraystretch}{1.3}
\begin{tabular}{|c|p{3cm}|p{3.5cm}|p{4cm}|p{2.5cm}|}
\hline
\rowcolor{blue!10}
\textbf{STT} & \textbf{Nghiệp vụ} & \textbf{Người dùng (Admin)} & \textbf{Phần mềm} & \textbf{Ghi chú} \\
\hline
1 & Thay đổi luật phân tích & Cập nhật đánh giá (Tốt/Kích ứng) dựa trên các nghiên cứu mới. & Tự động áp dụng luật mới cho toàn bộ các yêu cầu phân tích ngay lập tức. & Không cần sửa Code. \\
\hline
2 & Cập nhật tri thức hóa học & Thêm hóa chất mới hoặc tinh chỉnh mô tả chuyên môn. & Ghi nhận hoạt chất mới vào bộ từ điển nội bộ để nhận diện ở luồng phân tích. & Linh hoạt dữ liệu. \\
\hline
3 & Mở rộng dữ liệu sản phẩm & Bổ sung các sản phẩm đang thịnh hành trên thị trường. & Tự động đưa sản phẩm mới vào thuật toán lọc an toàn (\textit{Safety-filter}). & Đa dạng lựa chọn. \\
\hline
4 & Tự tiến hóa tri thức & Không cần thao tác (hoặc phê duyệt lại kết quả AI). & \textbf{AI Fallback:} Tự động gọi Gemini AI khi gặp chất lạ và \textbf{Auto-caching} ngược vào DB. &  \\
\hline
\end{tabular}
\end{center}

\textbf{Đánh giá khả năng tiến hóa:} 
Hệ thống được thiết kế theo kiến trúc hướng dữ liệu (\textit{Data-driven}). Mọi logic về độ an toàn đều nằm ở các bảng liên kết (\textit{Relation tables}), giúp đội ngũ vận hành có thể thay đổi toàn bộ "bộ não" của ứng dụng chỉ thông qua giao diện quản trị hoặc cơ chế tự học từ AI mà không cần lập trình viên can thiệp vào Logic Backend.

\subsection{Yêu cầu hệ thống}

\subsubsection{Yêu cầu bảo mật}
Yêu cầu bảo mật xác định các tiêu chuẩn và cơ chế để bảo vệ an toàn dữ liệu của người dùng và hệ thống (thông tin tài khoản, dữ liệu thành phần mỹ phẩm, tập luật) khỏi các truy cập trái phép hoặc các cuộc tấn công phá hoại.

\textbf{- Bảng yêu cầu bảo mật:}

\begin{center}
\begin{tabular}{|c|p{3cm}|p{5.5cm}|p{3.5cm}|}
\hline
\textbf{STT} & \textbf{Phạm vi} & \textbf{Yêu cầu bảo mật} & \textbf{Ghi chú} \\
\hline
1 & Bảo mật mật khẩu & Mật khẩu phải được mã hóa một chiều bằng thuật toán Salted Hashing trước khi lưu vào cơ sở dữ liệu PostgreSQL. & Sử dụng thư viện \textit{bcryptjs} với độ phức tạp (salt rounds) phù hợp. \\
\hline
2 & Xác thực và Phân quyền & Sử dụng cơ chế JSON Web Token (JWT) để duy trì phiên đăng nhập. Phân quyền rõ ràng qua Middleware: \textit{authMiddleware} cho người dùng và kết hợp thêm \textit{adminMiddleware} cho quản trị viên. & Token có thời hạn hết hạn (expiresIn) để giảm thiểu rủi ro. \\
\hline
3 & Trạng thái tài khoản & Hệ thống phải kiểm tra trường \textit{isActive} của người dùng. Các tài khoản bị khóa bởi Admin sẽ bị từ chối cấp Token ngay tại tầng Service. & Ngăn chặn truy cập từ người dùng vi phạm. \\
\hline
4 & Bảo vệ truy vấn dữ liệu & Sử dụng ORM (Prisma) để thực hiện các truy vấn tham số hóa, ngăn chặn triệt để tấn công SQL Injection. & Tận dụng khả năng bảo mật của Prisma Client. \\
\hline
5 & Giới hạn truy cập (Rate Limit) & Giới hạn endpoint phân tích theo \texttt{userId} của tài khoản đã xác thực trong cửa sổ 24 giờ; tài khoản Admin được bỏ qua giới hạn. & Sử dụng \textit{express-rate-limit} để hạn chế lạm dụng tài nguyên phân tích. \\
\hline
\end{tabular}
\end{center}

\textbf{- Bảng trách nhiệm yêu cầu bảo mật:}

\begin{center}
\begin{tabular}{|c|p{3.5cm}|p{3.5cm}|p{4.5cm}|p{2cm}|}
\hline
\textbf{STT} & \textbf{Nghiệp vụ} & \textbf{Người dùng} & \textbf{Phần mềm} & \textbf{Ghi chú} \\
\hline
1 & Đăng ký và Đăng nhập & Tự bảo mật mật khẩu cá nhân, không chia sẻ JWT Token. & Thực hiện mã hóa mật khẩu, cấp và xác thực Token theo tiêu chuẩn bảo mật. & \\
\hline
2 & Truy cập trang quản trị & Không cố tình truy cập các đường dẫn trái phép (\textit{/admin/*}). & Sử dụng Middleware chặn các yêu cầu không có quyền Admin, trả về mã lỗi 403 Forbidden. & \\
\hline
3 & Gửi yêu cầu phân tích & Cung cấp chuỗi INCI hợp lệ qua giao diện. & Thực hiện làm sạch dữ liệu (Sanitization), giới hạn tần suất gửi tin để bảo vệ tài nguyên Server. & \\
\hline
\end{tabular}
\end{center}


\subsubsection{Yêu cầu an toàn}
Yêu cầu an toàn đảm bảo hệ thống hoạt động ổn định, có khả năng phòng vệ trước các lỗi logic, sự cố hạ tầng hoặc thao tác sai từ người dùng mà không làm mất mát dữ liệu.

\textbf{- Bảng yêu cầu an toàn:}

\begin{center}
\begin{tabular}{|c|p{3.5cm}|p{5.3cm}|p{3.5cm}|}
\hline
\textbf{STT} & \textbf{Phạm vi} & \textbf{Yêu cầu an toàn} & \textbf{Ghi chú} \\
\hline
1 & Quản lý hạ tầng & Mã nguồn hiện tại chưa kèm cấu hình Docker/container. & Có thể bổ sung Dockerfile và cấu hình triển khai ở giai đoạn đóng gói hệ thống. \\
\hline
2 & Xử lý ngoại lệ (Exception Handling) & Toàn bộ các hàm xử lý nghiệp vụ phải nằm trong khối \textit{try-catch}. Các lỗi phát sinh phải được Log lại và trả về thông báo thân thiện cho người dùng thay vì làm sập tiến trình Node.js. & Sử dụng cơ chế tập trung xử lý lỗi. \\
\hline
3 & Toàn vẹn dữ liệu & Sử dụng quan hệ khóa ngoại (Foreign Keys) và cơ chế \textit{onDelete: Cascade} trong Prisma để đảm bảo không có dữ liệu "mồ côi" khi xóa tài khoản hoặc thành phần. & Ràng buộc chặt chẽ trong CSDL. \\
\hline
4 & Dự phòng AI API & Trong trường hợp Gemini API không phản hồi hoặc hết hạn mức, hệ thống phải tự động chuyển sang chế độ phân tích mặc định (Neutral) để không làm gián đoạn trải nghiệm người dùng. & Cơ chế Fallback an toàn. \\
\hline
\end{tabular}
\end{center}

\textbf{- Bảng trách nhiệm yêu cầu an toàn:}

\begin{center}
\begin{tabular}{|c|p{3cm}|p{3.5cm}|p{4.5cm}|p{2cm}|}
\hline
\textbf{STT} & \textbf{Nghiệp vụ} & \textbf{Người dùng} & \textbf{Phần mềm} & \textbf{Ghi chú} \\
\hline
1 & Lưu trữ dữ liệu & Không có thao tác trực tiếp. & Tự động ghi nhật ký lỗi (Logger) và duy trì các ràng buộc dữ liệu thông qua Prisma ORM. & \\
\hline
2 & Phân tích thành phần & Đảm bảo nhập liệu văn bản không chứa các ký tự điều khiển hệ thống. & Sử dụng cơ chế xử lý chuỗi (String manipulation) an toàn trong TypeScript để bóc tách thành phần. & \\
\hline
\end{tabular}
\end{center}

\subsection{Yêu cầu công nghệ}
Yêu cầu công nghệ xác định các nền tảng, ngôn ngữ lập trình, hệ quản trị cơ sở dữ liệu và công cụ hỗ trợ hiện đại được sử dụng để phát triển hệ thống SKINMATE. Việc lựa chọn công nghệ dựa trên tiêu chí tối ưu hiệu suất, tính an toàn dữ liệu cao và khả năng tích hợp trí tuệ nhân tạo (AI) linh hoạt.

\textbf{- Bảng yêu cầu công nghệ:}

\begin{center}
\begin{tabular}{|c|p{3.5cm}|p{5.5cm}|p{3.5cm}|}
\hline
\textbf{STT} & \textbf{Thành phần} & \textbf{Công nghệ / Công cụ sử dụng} & \textbf{Mục đích / Ghi chú} \\
\hline
1 & Giao diện người dùng (Frontend) & Next.js 16 (App Router), React 19, TypeScript, Tailwind CSS v4, Recharts. & Xây dựng giao diện Luxury, tối ưu Render phía Server (SSR) và Client (CSR). \\
\hline
2 & Xử lý nghiệp vụ (Backend) & Node.js (TypeScript), Express.js framework. & Xử lý logic nghiệp vụ, quản lý phiên bằng JWT và mã hóa mật khẩu bằng bcryptjs. \\
\hline
3 & Trí tuệ nhân tạo (AI Fallback) & Google Gemini Flash API (\texttt{gemini-flash-latest}). & Phân tích thành phần chưa có trong CSDL và tự động hóa việc gán nhãn an toàn. \\
\hline
4 & Cơ sở dữ liệu \& ORM & Hệ quản trị CSDL PostgreSQL 15 và Prisma ORM v5. & Quản lý quan hệ dữ liệu phức tạp, đảm bảo tính nhất quán qua Prisma Schema. \\
\hline
5 & Hạ tầng \& Công cụ & Node.js/npm, Visual Studio Code, Git/GitHub, Prisma Studio. & Vận hành môi trường phát triển và quản lý mã nguồn. \\
\hline
\end{tabular}
\end{center}


\textbf{- Bảng trách nhiệm yêu cầu công nghệ (Hệ thống):}

\begin{center}
\renewcommand{\arraystretch}{1.2}
\begin{tabular}{|c|p{3.5cm}|p{4cm}|p{4.5cm}|p{2cm}|}
\hline
\textbf{STT} & \textbf{Thành phần} & \textbf{Người dùng} & \textbf{Phần mềm / Hệ thống} & \textbf{Ghi chú} \\
\hline
1 & Giao diện (Frontend) & Truy cập qua trình duyệt hiện đại. Thao tác mượt mà nhờ React 19. & Đảm bảo tính Responsive toàn diện, hiển thị biểu đồ thống kê bằng Recharts. & Tối ưu hóa UI/UX. \\
\hline
2 & Xử lý (Backend) & Gửi yêu cầu (Request) qua RESTful API. & Kiểm soát lỗi bằng TypeScript, thực hiện xác thực người dùng và gọi AI API khi cần. & Xử lý Type-safe. \\
\hline
3 & Lưu trữ (Database \& ORM) & Nhập liệu thông tin cá nhân và bảng thành phần (INCI). & Prisma ORM tự động thực thi các truy vấn SQL tối ưu, quản lý quan hệ khóa ngoại giữa User và History. & PostgreSQL 15. \\
\hline
\end{tabular}
\end{center}

\subsection{Phân tích yêu cầu}

Dựa trên các yêu cầu chức năng và phi chức năng đã thu thập, nhóm tiến hành phân tích tính khả thi của dự án SKINMATE trên 3 phương diện cốt lõi:

\begin{itemize}
    \item \textbf{Tính khả thi kỹ thuật:} Hoàn toàn khả thi. Hệ thống ứng dụng kiến trúc Web hiện đại với các công nghệ mạnh mẽ như Next.js 16, TypeScript và Prisma ORM kết hợp cơ sở dữ liệu PostgreSQL. Mặc dù logic cốt lõi là hệ thống đối chiếu tập luật (Rule-based), nhưng dự án đã tích hợp \textbf{Google Gemini Flash API} (model \texttt{gemini-flash-latest}) để giải quyết bài toán các thành phần nằm ngoài cơ sở dữ liệu. Việc sử dụng AI dưới dạng API giúp hệ thống thông minh hơn mà không đòi hỏi hạ tầng máy chủ phức tạp để huấn luyện mô hình. TypeScript đảm bảo tính an toàn về kiểu dữ liệu, giúp giảm thiểu rủi ro trong việc xử lý các chuỗi văn bản (INCI) phức tạp từ phía người dùng.

    \item \textbf{Tính khả thi nghiệp vụ:} Hoàn toàn khả thi. Quy trình nghiệp vụ được chuẩn hóa từ thực tế: Tiếp nhận bảng thành phần $\rightarrow$ Bóc tách dữ liệu $\rightarrow$ Đối chiếu tập luật/AI $\rightarrow$ Cá nhân hóa theo loại da. Việc tin học hóa quy trình này không chỉ giúp tiết kiệm thời gian mà còn mang lại độ chính xác cao nhờ khả năng lưu trữ tri thức chuyên môn vào CSDL tập trung. Tính năng gợi ý sản phẩm và theo dõi lịch sử phân tích giúp khép kín vòng đời trải nghiệm người dùng, tạo ra giá trị nghiệp vụ thực tiễn cho cộng đồng làm đẹp.

    \item \textbf{Tính khả thi về con người:} Đối với người dùng cuối (Gen Z, học sinh, sinh viên), hệ thống rất dễ tiếp cận nhờ giao diện Luxury theo xu hướng "Clean Beauty", tối ưu hóa thao tác Copy-Paste. Về phía quản trị, rủi ro về khối lượng nhập liệu chuyên môn đã được giảm bớt đáng kể nhờ cơ chế \textbf{AI Auto-caching}: AI sẽ gợi ý mô tả và mức độ an toàn của chất mới, Admin chỉ đóng vai trò kiểm duyệt và tinh chỉnh luật. Điều này giúp hệ thống có thể vận hành ổn định ngay cả khi đội ngũ quản trị không có quá nhiều nhân sự chuyên sâu về hóa học.
\end{itemize}

\section{Mô hình hoá yêu cầu chức năng}
\subsection{Tổng quan về mô hình hoá yêu cầu}

Mô hình hoá yêu cầu chức năng là bước quan trọng trong quá trình phân tích và thiết kế hệ thống, nhằm xác định rõ các chức năng mà hệ thống cần cung cấp cũng như cách thức người dùng tương tác với hệ thống. Việc mô hình hoá giúp chuyển đổi các yêu cầu nghiệp vụ ban đầu thành các biểu diễn trực quan, dễ hiểu, hỗ trợ cho việc thiết kế và triển khai hệ thống sau này.

Trong đồ án này, hệ thống được xây dựng với mục tiêu hỗ trợ người dùng phân tích bảng thành phần mỹ phẩm (INCI) dựa trên loại da, từ đó đưa ra đánh giá mức độ phù hợp của từng thành phần và gợi ý sản phẩm phù hợp. Do đó, việc mô hình hoá yêu cầu tập trung vào các chức năng chính như: quản lý tài khoản người dùng, phân tích thành phần mỹ phẩm, quản lý dữ liệu thành phần và luật phân tích, cũng như đề xuất sản phẩm.

Quá trình mô hình hoá được thực hiện thông qua các công cụ và kỹ thuật như Use Case Diagram, Use Case Specification và các biểu đồ luồng dữ liệu. Các Use Case được xây dựng nhằm mô tả chi tiết các tương tác giữa tác nhân (người dùng, quản trị viên) và hệ thống, từ đó xác định rõ phạm vi chức năng của từng thành phần trong hệ thống.

Cụ thể, hệ thống bao gồm hai nhóm chức năng chính:
\begin{itemize}
\item Nhóm chức năng xử lý dữ liệu (Lưu trữ): bao gồm các chức năng như đăng ký tài khoản, cập nhật hồ sơ người dùng, quản lý thành phần, quản lý luật phân tích và quản lý sản phẩm.
\item Nhóm chức năng khai thác dữ liệu (Tra cứu): bao gồm các chức năng như đăng nhập, phân tích thành phần mỹ phẩm, xem chi tiết thành phần và xem gợi ý sản phẩm.
\end{itemize}

Việc phân chia rõ ràng các nhóm chức năng giúp hệ thống dễ dàng tổ chức, thiết kế và kiểm soát luồng dữ liệu. Đồng thời, đây cũng là cơ sở để xây dựng các sơ đồ chi tiết như sơ đồ luồng dữ liệu (DFD) và sơ đồ lớp (Class Diagram) trong các phần tiếp theo của báo cáo.

Nhờ vào quá trình mô hình hoá yêu cầu, hệ thống được định hình một cách rõ ràng về mặt chức năng, đảm bảo đáp ứng đúng nhu cầu của người dùng cũng như tạo tiền đề cho việc phát triển và mở rộng trong tương lai.

\subsection{Mô hình hoá theo Chức năng (DFD - Data Flow Diagram}

\subsubsection {Nhóm Use Case dùng chung}
\textbf {- Đăng nhập}
\begin{figure}[H]
    \centering
    \includegraphics[width=0.75\linewidth]{image.png}
\end{figure}

Mô tả luồng dữ liệu:
\begin{itemize}
    \item D1: Username, Password.
    \item D2: Không có.
    \item D3: Thông tin tài khoản trong CSDL.
    \item D4: Không có.
    \item D5: Kết quả xác thực.
    \item D6: D5.
\end{itemize}

Thuật toán:
\begin{enumerate}
    \item Nhận D1 từ người dùng.
    \item Kết nối CSDL.
    \item Đọc D3 từ bảng Users.
    \item So sánh Username và Password.
    \item Nếu sai $\rightarrow$ trả D6 (Thất bại) $\rightarrow$ Kết thúc.
    \item Nếu đúng $\rightarrow$ trả D6 (Thành công).
    \item Đóng kết nối.
\end{enumerate}

\textbf{- Cập nhật hồ sơ cá nhân}

\begin{figure}[H]
    \centering
    \includegraphics[width=0.75\linewidth]{Screenshot 2026-05-06 201725.png} % Thay tên file ảnh của bạn
    \caption{Sơ đồ luồng dữ liệu chức năng Cập nhật hồ sơ cá nhân}
\end{figure}

\vspace{0.5cm}

Mô tả luồng dữ liệu:
\begin{itemize}
    \item D1: Thông tin cập nhật (Loại da mới, Tên hiển thị mới).
    \item D2: Thao tác nhập liệu từ bàn phím và chuột.
    \item D3: Thông tin User hiện tại từ CSDL (để đối chiếu).
    \item D4: Dữ liệu hồ sơ mới cần ghi đè (Update).
    \item D5: Tín hiệu thay đổi giao diện Navbar (cập nhật tên/loại da trên màn hình).
    \item D6: Thông báo kết quả "Cập nhật thành công" (hoặc báo lỗi).
\end{itemize}

Thuật toán:
\begin{enumerate}
    \item Nhận D1 (Thông tin cập nhật) từ người dùng và D2 từ thiết bị nhập.
    \item Kết nối CSDL.
    \item Đọc D3 (Thông tin User hiện tại) từ bảng Users dựa trên Token xác thực.
    \item Kiểm tra tính hợp lệ của D1 (Loại da phải nằm trong danh mục cho phép).
    \item Nếu không hợp lệ $\rightarrow$ trả D6 (Lỗi) $\rightarrow$ Kết thúc.
    \item Cập nhật D4 (Dữ liệu mới) vào bảng Users trong CSDL.
    \item Trả D5 để Frontend cập nhật lại giao diện Navbar (Reactive UI).
    \item Trả D6 (Thông báo thành công) cho người dùng.
    \item Đóng kết nối.
\end{enumerate}

\textbf{- Thay đổi mật khẩu}

\begin{figure}[H]
    \centering
    \includegraphics[width=0.75\linewidth]{Screenshot 2026-05-06 201746.png} % Thay tên file ảnh của bạn
    \caption{Sơ đồ luồng dữ liệu chức năng Thay đổi mật khẩu}
\end{figure}

\vspace{0.5cm}

Mô tả luồng dữ liệu:
\begin{itemize}
    \item D1: Mật khẩu cũ, Mật khẩu mới.
    \item D2: Thao tác nhập liệu từ bàn phím.
    \item D3: Mật khẩu băm (Password Hash) hiện tại trong CSDL.
    \item D4: Mật khẩu băm mới (đã qua mã hóa Bcrypt) cần lưu trữ.
    \item D5: Tín hiệu đóng form và làm sạch các ô nhập liệu.
    \item D6: Thông báo kết quả (Thành công/Lỗi).
\end{itemize}

Thuật toán:
\begin{enumerate}
    \item Nhận D1 (Mật khẩu cũ, Mật khẩu mới) từ người dùng và D2 từ thiết bị nhập.
    \item Kết nối CSDL.
    \item Lấy D3 (Mật khẩu băm hiện tại) của User từ bảng Users.
    \item Dùng thuật toán Bcrypt so sánh "Mật khẩu cũ" (trong D1) với D3.
    \item Nếu sai $\rightarrow$ trả D6 (Lỗi: Sai mật khẩu hiện tại) $\rightarrow$ Kết thúc.
    \item Nếu mật khẩu mới < 6 ký tự $\rightarrow$ trả D6 (Lỗi định dạng) $\rightarrow$ Kết thúc.
    \item Mã hóa (Hash) "Mật khẩu mới" thành D4.
    \item Lưu D4 vào CSDL (Ghi đè mật khẩu cũ).
    \item Trả D5 để thu gọn form nhập liệu trên màn hình.
    \item Trả D6 (Thông báo đổi mật khẩu thành công) cho người dùng.
    \item Đóng kết nối.
\end{enumerate}

\subsubsection{Nhóm Use Case Dành Cho Người Dùng (User)}
\textbf{- Đăng ký tài khoản}
\begin{figure}[H]
    \centering
    \includegraphics[width=0.75\linewidth]{Screenshot 2026-04-22 120437.png}
\end{figure}

Mô tả luồng dữ liệu:
\begin{itemize}
    \item D1: Thông tin đăng ký (Username, Password, Loại da).
    \item D2: Không có.
    \item D3: Danh sách Username hiện có.
    \item D4: Thông tin user hợp lệ cần lưu trữ.
    \item D5: Không có.
    \item D6: Thông báo kết quả (Đăng ký thành công/Lỗi).
\end{itemize}

Thuật toán:
\begin{enumerate}
    \item Nhận D1 từ người dùng.
    \item Kết nối cơ sở dữ liệu.
    \item Đọc D3 từ bảng Users.
    \item Kiểm tra Username đã tồn tại chưa.
    \item Nếu trùng $\rightarrow$ trả D6 (Lỗi) $\rightarrow$ Kết thúc.
    \item Lưu D4 vào CSDL.
    \item Trả D6 (Thành công).
    \item Đóng kết nối.
\end{enumerate}




\textbf{- Phân tích thành phần mỹ phẩm (Core Flow)}

\begin{figure}[H]
    \centering
    \includegraphics[width=0.75\linewidth]{UC04.png}
\end{figure}

\vspace{0.5cm}
Mô tả luồng dữ liệu:
\begin{itemize}
    \item D1: Chuỗi thành phần mỹ phẩm (INCI).
    \item D2: Không có.
    \item D3: Dữ liệu Users, Ingredients, Ingredient\_Rules.
    \item D4: Không có.
    \item D5: Kết quả phân tích (Đỏ/Xanh/Xám).
    \item D6: D5.
\end{itemize}

Thuật toán:
\begin{enumerate}
    \item Nhận D1 từ người dùng.
    \item Kết nối CSDL.
    \item Lấy loại da user và dữ liệu thành phần, luật.
    \item Tách danh sách thành phần từ chuỗi INCI.
    \item Với mỗi thành phần:
    \begin{itemize}
        \item So khớp với luật.
        \item Gán nhãn Good/Bad/Neutral.
    \end{itemize}
    \item Tổng hợp kết quả D5.
    \item Trả D6 cho người dùng.
    \item Đóng kết nối.
\end{enumerate}

\textbf{ - Xem chi tiết mô tả thành phần}

\begin{figure}[H]
    \centering
    \includegraphics[width=0.75\linewidth]{UC05.png}
\end{figure}
\vspace{0.5cm}

Mô tả luồng dữ liệu:
\begin{itemize}
    \item D1: Yêu cầu xem chi tiết thành phần.
    \item D2: Không có.
    \item D3: Mô tả thành phần từ CSDL.
    \item D4: Không có.
    \item D5: Nội dung chi tiết.
    \item D6: D5.
\end{itemize}

Thuật toán:
\begin{enumerate}
    \item Nhận D1.
    \item Kết nối CSDL.
    \item Truy vấn bảng Ingredients.
    \item Lấy thông tin D3.
    \item Trả D6.
    \item Đóng kết nối.
\end{enumerate}

\textbf{- Xem gợi ý sản phẩm}


\begin{figure}[H]
    \centering
    \includegraphics[width=0.75\linewidth]{UC08.png}
\end{figure}

\vspace{0.5cm}
Mô tả luồng dữ liệu:
\begin{itemize}
    \item D1: Tiêu chí lọc (loại da).
    \item D2: Không có.
    \item D3: Danh sách sản phẩm và thành phần.
    \item D4: Không có.
    \item D5: Danh sách sản phẩm phù hợp.
    \item D6: D5.
\end{itemize}

Thuật toán:
\begin{enumerate}
    \item Nhận D1.
    \item Kết nối CSDL.
    \item Lấy danh sách sản phẩm.
    \item Lọc sản phẩm theo loại da và thành phần.
    \item Sắp xếp theo mức độ phù hợp.
    \item Trả D6.
    \item Đóng kết nối.
\end{enumerate}


\textbf{- Quản lý lịch sử phân tích}

\begin{figure}[H]
    \centering
    \includegraphics[width=0.75\linewidth]{Screenshot 2026-05-06 201754.png} % Thay tên file ảnh của bạn
    \caption{Sơ đồ luồng dữ liệu chức năng Quản lý lịch sử phân tích}
\end{figure}

\vspace{0.5cm}

Mô tả luồng dữ liệu:
\begin{itemize}
    \item D1: Yêu cầu thao tác (Xem danh sách / Xóa 1 bản ghi / Xóa tất cả).
    \item D2: Thao tác click chuột từ người dùng.
    \item D3: Danh sách các lần phân tích cũ từ CSDL.
    \item D4: Lệnh xóa (ID bản ghi cụ thể hoặc Lệnh xóa toàn bộ).
    \item D5: Dữ liệu danh sách thẻ (History Cards) để hiển thị lên màn hình.
    \item D6: Hiển thị kết quả lịch sử hoặc Thông báo xóa thành công.
\end{itemize}

Thuật toán:
\begin{enumerate}
    \item Nhận D1 (Yêu cầu thao tác) từ người dùng và D2 từ thiết bị nhập.
    \item Kết nối CSDL, xác định danh tính User qua Token.
    \item \textbf{Trường hợp "Xem danh sách":}
    \begin{itemize}
        \item Truy vấn lấy D3 (Danh sách lịch sử) sắp xếp theo thời gian mới nhất.
        \item Trả D5 để vẽ danh sách lên giao diện và D6 hiển thị cho người dùng.
    \end{itemize}
    \item \textbf{Trường hợp "Xóa bản ghi":}
    \begin{itemize}
        \item Kiểm tra tính hợp lệ của bản ghi.
        \item Gửi D4 (Lệnh xóa) xuống CSDL để loại bỏ dữ liệu.
        \item Trả D5 (Làm mới danh sách) và D6 (Thông báo xóa thành công).
    \end{itemize}
    \item Đóng kết nối.
\end{enumerate}


\subsubsection{Nhóm Use Case Dành Cho Quản Trị Viên (Admin)}

\textbf{- Quản lý danh sách thành phần}

\begin{figure}[H]
    \centering
    \includegraphics[width=0.75\linewidth]{UC06.png}
\end{figure}

\vspace{0.5cm}
Mô tả luồng dữ liệu:
\begin{itemize}
    \item D1: Thông tin thành phần (Thêm/Sửa/Xóa).
    \item D2: Không có.
    \item D3: Danh sách thành phần hiện có.
    \item D4: Dữ liệu cập nhật.
    \item D5: Không có.
    \item D6: Thông báo kết quả.
\end{itemize}

Thuật toán:
\begin{enumerate}
    \item Nhận D1 từ Admin.
    \item Kết nối CSDL.
    \item Đọc D3.
    \item Kiểm tra trùng hoặc lỗi dữ liệu.
    \item Nếu lỗi $\rightarrow$ trả D6 $\rightarrow$ Kết thúc.
    \item Lưu D4 vào CSDL.
    \item Trả D6 (Thành công).
    \item Đóng kết nối.
\end{enumerate}
\textbf{ - Quản lý luật phân tích (Rules)}


\begin{figure}[H]
    \centering
    \includegraphics[width=0.75\linewidth]{UC07.png}
\end{figure}
\vspace{0.5cm}

Mô tả luồng dữ liệu:
\begin{itemize}
    \item D1: Thông tin luật.
    \item D2: Không có.
    \item D3: Danh sách luật hiện có.
    \item D4: Luật cập nhật.
    \item D5: Không có.
    \item D6: Thông báo kết quả.
\end{itemize}

Thuật toán:
\begin{enumerate}
    \item Nhận D1 từ Admin.
    \item Kết nối CSDL.
    \item Đọc D3.
    \item Kiểm tra trùng luật.
    \item Nếu trùng $\rightarrow$ trả D6 $\rightarrow$ Kết thúc.
    \item Lưu D4 vào CSDL.
    \item Trả D6 (Thành công).
    \item Đóng kết nối.
\end{enumerate}

\textbf{- Quản lý danh mục sản phẩm}


\begin{figure}[H]
    \centering
    \includegraphics[width=0.8\linewidth]{UC09.png}
\end{figure}

Mô tả luồng dữ liệu:
\begin{itemize}
    \item D1: Thông tin sản phẩm.
    \item D2: Không có.
    \item D3: Danh sách sản phẩm hiện có.
    \item D4: Dữ liệu cập nhật.
    \item D5: Không có.
    \item D6: Thông báo kết quả.
\end{itemize}

Thuật toán:
\begin{enumerate}
    \item Nhận D1 từ Admin.
    \item Kết nối CSDL.
    \item Đọc D3.
    \item Kiểm tra hợp lệ dữ liệu.
    \item Nếu lỗi $\rightarrow$ trả D6 $\rightarrow$ Kết thúc.
    \item Lưu D4 vào CSDL.
    \item Trả D6 (Thành công).
    \item Đóng kết nối.
\end{enumerate}

\textbf{- Quản lý người dùng}


\begin{figure}[H]
    \centering
    \includegraphics[width=0.8\linewidth]{quanlynguoidung.png}
\end{figure}

Mô tả luồng dữ liệu:
\begin{itemize}
    \item D1: Thông tin người dùng (Thêm/Sửa/Xóa: Username, Password, Role, Loại da).
    \item D2: Không có.
    \item D3: Danh sách người dùng hiện có (dùng để kiểm tra trùng lặp và hiển thị).
    \item D4: Thông tin người dùng hợp lệ cần lưu trữ.
    \item D5: Thông báo kết quả (Thành công/Thất bại).
    \item D6: Không có.
\end{itemize}

Thuật toán:
\begin{itemize}
    \item B1: Nhận D1 (thông tin người dùng) từ Quản trị viên.
    \item B2: Kết nối cơ sở dữ liệu.
    \item B3: Đọc D3 (danh sách người dùng hiện có) từ bộ nhớ phụ.
    \item B4: Kiểm tra tính hợp lệ:
    \begin{itemize}
        \item Nếu Thêm: Kiểm tra Username không được trùng.
        \item Nếu Sửa: Kiểm tra dữ liệu cập nhật hợp lệ.
        \item Nếu Xóa: Kiểm tra người dùng tồn tại trong hệ thống.
    \end{itemize}
    \item B5: Nếu vi phạm quy định, trả D5 (Thông báo lỗi) và kết thúc.
    \item B6: Lưu D4 (thông tin người dùng) vào bảng Users.
    \item B7: Trả D5 (Thông báo thành công) cho Quản trị viên.
    \item B8: Đóng kết nối cơ sở dữ liệu.
    \item B9: Kết thúc.
\end{itemize}

\textbf{- Xem báo cáo thống kê}


\begin{figure}[H]
    \centering
    \includegraphics[width=0.8\linewidth]{baocaothongke.png}
\end{figure}
Mô tả luồng dữ liệu:
\begin{itemize}
    \item D1: Tiêu chí báo cáo (Ngày bắt đầu, Ngày kết thúc, loại báo cáo nếu có).
    \item D2: Không có.
    \item D3: Dữ liệu cần thống kê (Danh sách hóa đơn / lịch sử phân tích từ cơ sở dữ liệu).
    \item D4: Không có.
    \item D5: Kết quả báo cáo (Tổng hợp dữ liệu, biểu đồ, danh sách thống kê).
    \item D6: D5 (Hiển thị kết quả báo cáo cho người dùng).
\end{itemize}

Thuật toán:
\begin{itemize}
    \item B1: Nhận D1 (tiêu chí báo cáo) từ Quản trị viên.
    \item B2: Kết nối cơ sở dữ liệu.
    \item B3: Truy vấn dữ liệu (D3) từ bộ nhớ phụ dựa trên tiêu chí đã chọn.
    \item B4: Kiểm tra dữ liệu:
    \begin{itemize}
        \item Nếu không có dữ liệu phù hợp, trả D5 (Thông báo “Không có dữ liệu”) và kết thúc.
    \end{itemize}
    \item B5: Xử lý và tổng hợp dữ liệu:
    \begin{itemize}
        \item Tính toán tổng số lượng, tổng giá trị (nếu có).
        \item Gom nhóm dữ liệu theo tiêu chí (ngày, sản phẩm, loại da,...).
    \end{itemize}
    \item B6: Tạo kết quả báo cáo (D5) dưới dạng bảng hoặc biểu đồ.
    \item B7: Trả D6 (hiển thị kết quả báo cáo) cho người dùng.
    \item B8: Đóng kết nối cơ sở dữ liệu.
    \item B9: Kết thúc.
\end{itemize}

\textbf{- Quản lý dữ liệu qua tệp Excel (Import / Export)}

\begin{figure}[H]
    \centering
    \includegraphics[width=0.8\linewidth]{Screenshot 2026-05-06 202239.png} 
    \caption{Sơ đồ luồng dữ liệu chức năng Import / Export Excel}
\end{figure}

\vspace{0.5cm}

Mô tả luồng dữ liệu:
\begin{itemize}
    \item D1: Yêu cầu tải file (.xlsx) lên (Import), hoặc Yêu cầu xuất dữ liệu (Export).
    \item D2: Tệp tin vật lý Excel (.xlsx) truyền vào từ thiết bị nhập.
    \item D3: Toàn bộ dữ liệu hiện có trong CSDL (Ingredients, Rules, Products).
    \item D4: Dữ liệu hàng loạt từ file Excel cần ghi vào CSDL.
    \item D5: Bảng thống kê kết quả Import hoặc luồng tải (Stream) file Export.
    \item D6: Báo cáo Import hoàn tất hoặc File Excel được tải về máy.
\end{itemize}

Thuật toán:
\begin{enumerate}
    \item Nhận D1 (Yêu cầu thao tác) từ Quản trị viên (Admin).
    \item Kết nối CSDL.
    \item \textbf{Nếu là yêu cầu Export:}
    \begin{itemize}
        \item Truy vấn đọc D3 (Toàn bộ dữ liệu) từ các bảng tương ứng.
        \item Sử dụng thư viện \textit{xlsx} chuyển đổi D3 thành dạng bảng tính Excel.
        \item Trả D5 (Luồng dữ liệu) và D6 (File .xlsx) để trình duyệt tự động tải về.
    \end{itemize}
    \item \textbf{Nếu là yêu cầu Import:}
    \begin{itemize}
        \item Nhận D2 (File Excel) đưa vào bộ nhớ đệm (Memory Storage).
        \item Đọc File, bóc tách và đối chiếu dữ liệu trong File với D3 từ CSDL.
        \item Thực hiện cơ chế \textit{Upsert}: Cập nhật nếu trùng tên, tạo mới nếu chưa có.
        \item Lưu D4 (Dữ liệu hàng loạt) vào CSDL.
        \item Trả D5 (Thống kê số dòng thành công/lỗi) và D6 (Báo cáo chi tiết) cho Admin.
    \end{itemize}
    \item Đóng kết nối.
\end{enumerate}


\subsection{Sơ đồ Sử dụng Chức năng (Use Case Diagram)}

\begin{figure}[H]
    \centering
    \includegraphics[width=0.8\linewidth]{Screenshot 2026-05-06 221718.png}
\end{figure}

\subsubsection{Danh sách các Actor}

Trong hệ thống SKINMATE, các tác nhân (Actors) được xác định dựa trên vai trò và quyền hạn tương tác với các chức năng của phần mềm. Hệ thống bao gồm hai tác nhân chính như sau:

\begin{center}
\renewcommand{\arraystretch}{1.5}
\rowcolors{2}{gray!10}{white}
\begin{tabular}{|c|p{3cm}|p{9cm}|}
\hline
\rowcolor{blue!20}
\textbf{STT} & \textbf{Tên Actor} & \textbf{Mô tả vai trò và trách nhiệm} \\ \hline
1 & \textbf{Người dùng} & Là đối tượng sử dụng hệ thống để phân tích thành phần INCI. Có trách nhiệm quản lý hồ sơ cá nhân (loại da), theo dõi và quản lý lịch sử phân tích, đồng thời xem các sản phẩm gợi ý an toàn dựa trên kết quả kiểm tra. \\ \hline
2 & \textbf{Quản trị viên} & Là người có quyền hạn cao nhất, chịu trách nhiệm quản trị toàn bộ dữ liệu hệ thống (Ingredients, Rules, Products). Admin có quyền quản lý người dùng (khóa/mở khóa/xóa tài khoản), thực hiện xuất/nhập dữ liệu qua file Excel và theo dõi hiệu quả hoạt động thông qua các báo cáo thống kê trực quan. \\ \hline
\end{tabular}
\end{center}

\textit{Ghi chú: Hệ thống sử dụng cơ chế xác thực tập trung. Cả hai tác nhân đều cần thực hiện hành động Đăng nhập để được cấp mã xác thực (JWT), từ đó hệ thống sẽ phân quyền truy cập vào các chức năng tương ứng theo vai trò (Role-based Access Control).}


\subsubsection{Sơ đồ Use Case}

\subsubsection{Đặc tả các Use Case}

Dưới đây là bảng đặc tả chi tiết cho từng chức năng (Use Case) của hệ thống SKINMATE, mô tả cách thức tương tác giữa các tác nhân và hệ thống để đạt được mục tiêu cụ thể.

\vspace{0.5cm}

\textbf{*Nhóm Use Case dùng chung}


\vspace{0.5cm}
\textbf{Đặc tả Use Case: Đăng nhập}

\begin{center}
\renewcommand{\arraystretch}{1.4}
\begin{tabular}{|p{4cm}|p{10cm}|}
\hline
\rowcolor{blue!15}
\textbf{Tên Use Case} & \textbf{Đăng nhập} \\ \hline
\textbf{Tác nhân} & Người dùng, Quản trị viên \\ \hline
\textbf{Mô tả} & Xác thực danh tính tác nhân thông qua cơ chế so khớp mật khẩu băm và kiểm tra trạng thái tài khoản để cấp quyền truy cập (JWT Token). \\ \hline
\textbf{Tiền điều kiện} & Tác nhân đã có tài khoản hợp lệ trong hệ thống. \\ \hline
\textbf{Hậu điều kiện} & Hệ thống cấp mã xác thực JWT lưu tại LocalStorage; tác nhân truy cập vào các tính năng theo đúng vai trò (User/Admin). \\ \hline
\textbf{Luồng sự kiện chính} & 
\begin{enumerate}
    \item Tác nhân nhập Tên đăng nhập và Mật khẩu vào Biểu mẫu đăng nhập (BM 2).
    \item Hệ thống truy vấn CSDL để tìm tài khoản tương ứng với Username.
    \item Hệ thống kiểm tra trạng thái hoạt động của tài khoản (trường \textit{isActive}).
    \item Hệ thống thực hiện so khớp mật khẩu nhập vào với mật khẩu đã mã hóa trong CSDL bằng thuật toán \textbf{Bcrypt}.
    \item Sau khi xác thực thành công, hệ thống tạo một mã \textbf{JSON Web Token (JWT)} chứa mã người dùng, vai trò và loại da.
    \item Hệ thống phản hồi mã Token về Frontend để lưu trữ vào \textit{AuthContext}.
    \item Frontend gọi hàm \texttt{login()} để lưu \texttt{token} và \texttt{user} vào \textit{AuthContext}/LocalStorage, sau đó chuyển hướng về trang chủ (\texttt{/}). Từ Navbar, tài khoản \texttt{ADMIN} có thêm liên kết truy cập trang quản trị.
\end{enumerate} \\ \hline
\textbf{Luồng sự kiện phụ / Ngoại lệ} & 
\begin{itemize}
    \item \textit{Tài khoản bị khóa:} Nếu trường \textit{isActive} là \textit{false}, hệ thống trả về thông báo lỗi: "Tài khoản của bạn đã bị khóa bởi Quản trị viên." và dừng quá trình.
    \item \textit{Sai thông tin:} Nếu Username không tồn tại hoặc mật khẩu không khớp, hệ thống báo lỗi chung: "Tên đăng nhập hoặc mật khẩu không hợp lệ." để đảm bảo tính bảo mật.
    \item \textit{Để trống thông tin:} Hệ thống yêu cầu nhập đầy đủ các trường bắt buộc trước khi gửi yêu cầu lên Server.
\end{itemize} \\ \hline
\end{tabular}
\end{center}



\vspace{0.5cm}
\textbf{Đặc tả Use Case: Cập nhật hồ sơ cá nhân}

\begin{center}
\renewcommand{\arraystretch}{1.4}
\begin{tabular}{|p{4cm}|p{10cm}|}
\hline
\rowcolor{blue!15}
\textbf{Tên Use Case} & \textbf{Cập nhật hồ sơ cá nhân} \\ \hline
\textbf{Tác nhân} & Người dùng, Quản trị viên \\ \hline
\textbf{Mô tả} & Cho phép người dùng thay đổi Tên hiển thị (Display Name) trên hệ thống và cập nhật Loại da hiện tại để điều chỉnh kết quả phân tích AI. \\ \hline
\textbf{Tiền điều kiện} & Người dùng đã đăng nhập thành công vào hệ thống. \\ \hline
\textbf{Hậu điều kiện} & Thông tin hồ sơ mới được ghi đè vào CSDL và tự động cập nhật ngay lập tức trên Navbar (Reactive UI). \\ \hline
\textbf{Luồng sự kiện chính} & 
\begin{enumerate}
    \item Người dùng truy cập vào trang "Hồ sơ" (/profile).
    \item Hệ thống gọi API lấy thông tin hiện tại (Tên đăng nhập, Tên hiển thị, Loại da) và điền sẵn vào Form. (Ghi chú: Tên đăng nhập bị vô hiệu hóa, không cho sửa).
    \item Người dùng chỉnh sửa Tên hiển thị (Tùy chọn) và chọn Loại da mới từ Dropdown.
    \item Người dùng nhấn nút "Lưu thay đổi".
    \item Hệ thống gọi API \texttt{PUT /api/v1/users/profile}.
    \item Hệ thống cập nhật dữ liệu vào bảng \textit{Users}.
    \item Hệ thống gọi lại hàm \texttt{login()} trong \textit{AuthContext} ở phía Frontend để cập nhật giao diện ngay lập tức mà không cần reload trang.
    \item Hiển thị thông báo "Cập nhật thành công" (Màu xanh).
\end{enumerate} \\ \hline
\textbf{Luồng sự kiện phụ / Ngoại lệ} & 
\begin{itemize}
    \item \textit{Để trống loại da:} Hệ thống báo lỗi yêu cầu chọn loại da hợp lệ.
    \item \textit{Không có thay đổi:} Nút "Lưu thay đổi" sẽ bị vô hiệu hóa (Disabled) nếu người dùng chưa sửa thông tin nào so với ban đầu.
\end{itemize} \\ \hline
\end{tabular}
\end{center}

\vspace{0.5cm}
\textbf{Đặc tả Use Case: Thay đổi mật khẩu}

\begin{center}
\renewcommand{\arraystretch}{1.4}
\begin{tabular}{|p{4cm}|p{10cm}|}
\hline
\rowcolor{blue!15}
\textbf{Tên Use Case} & \textbf{Thay đổi mật khẩu} \\ \hline
\textbf{Tác nhân} & Người dùng, Quản trị viên \\ \hline
\textbf{Mô tả} & Cho phép người dùng cập nhật mật khẩu mới để bảo vệ tài khoản. \\ \hline
\textbf{Tiền điều kiện} & Người dùng đã đăng nhập thành công. \\ \hline
\textbf{Hậu điều kiện} & Mật khẩu băm trong CSDL được cập nhật. \\ \hline
\textbf{Luồng sự kiện chính} & 
\begin{enumerate}
    \item Tại trang Hồ sơ, người dùng mở rộng khu vực "Cài đặt bảo mật".
    \item Người dùng nhập Mật khẩu hiện tại, Mật khẩu mới và Xác nhận mật khẩu mới.
    \item Người dùng nhấn "Cập nhật mật khẩu".
    \item Hệ thống kiểm tra phía Client: Mật khẩu mới trùng khớp và $\ge$ 6 ký tự.
    \item Gửi API \texttt{PUT /api/v1/users/change-password}.
    \item Backend dùng Bcrypt so sánh "Mật khẩu hiện tại" với DB. Nếu đúng, tiến hành băm (hash) mật khẩu mới và lưu vào bảng \textit{Users}.
    \item Hệ thống thông báo đổi mật khẩu thành công và tự động thu gọn Form.
\end{enumerate} \\ \hline
\textbf{Luồng sự kiện phụ / Ngoại lệ} & 
\begin{itemize}
    \item \textit{Sai mật khẩu cũ:} Backend trả về lỗi 400. Hệ thống hiển thị "Invalid current password".
    \item \textit{Mật khẩu mới không khớp:} Bắt lỗi ngay tại Client, không gửi request lên Server.
\end{itemize} \\ \hline
\end{tabular}
\end{center}

\textbf{*Nhóm Use Case Dành cho người dùng (user)}


\vspace{0.5cm}
\textbf{Đặc tả Use Case: Đăng ký tài khoản}

\begin{center}
\renewcommand{\arraystretch}{1.4}
\begin{tabular}{|p{4cm}|p{10cm}|}
\hline
\rowcolor{blue!15}
\textbf{Tên Use Case} & \textbf{Đăng ký tài khoản} \\ \hline
\textbf{Tác nhân} & Người dùng (Khách) \\ \hline
\textbf{Mô tả} & Cho phép người dùng mới tạo tài khoản trên hệ thống để lưu trữ thông tin cá nhân và hồ sơ loại da. \\ \hline
\textbf{Tiền điều kiện} & Người dùng truy cập vào trang web và chưa đăng nhập vào hệ thống. \\ \hline
\textbf{Hậu điều kiện} & Một tài khoản mới được tạo trong cơ sở dữ liệu; hệ thống đăng nhập tự động và chuyển người dùng về trang chủ khi lấy Token thành công. \\ \hline
\textbf{Luồng sự kiện chính} & 
\begin{enumerate}
    \item Người dùng nhấn nút "Đăng ký" trên thanh điều hướng hoặc trang chủ.
    \item Hệ thống hiển thị Biểu mẫu đăng ký (BM 1) bao gồm các trường: Tên đăng nhập, Mật khẩu, Nhập lại mật khẩu và Lựa chọn loại da.
    \item Người dùng điền đầy đủ các thông tin cần thiết.
    \item Người dùng nhấn nút "Xác nhận đăng ký".
    \item Hệ thống thực hiện kiểm tra các quy định (QĐ 1) về tính duy nhất của Username và độ dài mật khẩu.
    \item Hệ thống lưu thông tin người dùng vào bảng \textit{Users} trong cơ sở dữ liệu.
    \item Frontend tự gửi yêu cầu đăng nhập bằng thông tin vừa tạo; nếu thành công lưu Token và chuyển về trang chủ (\texttt{/}), nếu đăng nhập tự động thất bại thì chuyển sang trang \texttt{/login}.
\end{enumerate} \\ \hline
\textbf{Luồng sự kiện phụ / Ngoại lệ} & 
\begin{itemize}
    \item \textit{Trường hợp bỏ trống thông tin:} Nếu có trường dữ liệu bắt buộc bị để trống, hệ thống báo lỗi "Vui lòng điền đầy đủ thông tin".
    \item \textit{Trùng tên đăng nhập:} Nếu Username đã tồn tại trong CSDL, hệ thống báo lỗi "Tên đăng nhập đã được sử dụng" và yêu cầu chọn tên khác.
    \item \textit{Mật khẩu không khớp:} Nếu "Mật khẩu" và "Nhập lại mật khẩu" không giống nhau, hệ thống báo lỗi và yêu cầu nhập lại.
    \item \textit{Mật khẩu yếu:} Nếu mật khẩu dưới 6 ký tự, hệ thống yêu cầu người dùng tăng độ dài mật khẩu.
\end{itemize} \\ \hline
\end{tabular}
\end{center}






% UC 03

\vspace{0.5cm}
% \textbf{Đặc tả Use Case: Cập nhật hồ sơ loại da}

% \begin{center}
% \renewcommand{\arraystretch}{1.4}
% \begin{tabular}{|p{4cm}|p{10cm}|}
% \hline
% \rowcolor{blue!15}
% \textbf{Tên Use Case} & \textbf{Cập nhật hồ sơ loại da} \\ \hline
% \textbf{Tác nhân} & Người dùng \\ \hline
% \textbf{Mô tả} & Cho phép người dùng thay đổi hoặc cập nhật tình trạng loại da hiện tại để hệ thống điều chỉnh các luật phân tích thành phần mỹ phẩm tương ứng. \\ \hline
% \textbf{Tiền điều kiện} & Người dùng đã đăng nhập thành công vào hệ thống. \\ \hline
% \textbf{Hậu điều kiện} & Thông tin loại da mới được ghi đè vào cơ sở dữ liệu và áp dụng cho các lần phân tích tiếp theo. \\ \hline
% \textbf{Luồng sự kiện chính} & 
% \begin{enumerate}
%     \item Người dùng truy cập vào trang "Hồ sơ cá nhân" hoặc "Cài đặt tài khoản".
%     \item Hệ thống truy vấn và hiển thị thông tin loại da hiện tại của người dùng.
%     \item Người dùng chọn loại da mới từ danh sách thả xuống (Dropdown) trong Biểu mẫu cập nhật (BM 3). Các lựa chọn bao gồm: Da dầu, Da khô, Da nhạy cảm, Da hỗn hợp.
%     \item Người dùng nhấn nút "Cập nhật hồ sơ".
%     \item Hệ thống kiểm tra tính hợp lệ của dữ liệu theo Quy định 3 (QĐ 3).
%     \item Hệ thống cập nhật trường \textit{skin\_type} trong bảng \textit{Users} của cơ sở dữ liệu.
%     \item Hệ thống hiển thị thông báo "Cập nhật loại da thành công".
% \end{enumerate} \\ \hline
% \textbf{Luồng sự kiện phụ / Ngoại lệ} & 
% \begin{itemize}
%     \item \textit{Không có thay đổi:} Nếu người dùng nhấn cập nhật nhưng không thay đổi giá trị so với hiện tại, hệ thống thông báo "Thông tin không có sự thay đổi".
%     \item \textit{Lỗi lưu trữ:} Nếu quá trình ghi dữ liệu vào CSDL gặp sự cố, hệ thống thông báo "Không thể cập nhật, vui lòng thử lại sau" và giữ nguyên thông tin cũ.
%     \item \textit{Phiên đăng nhập hết hạn:} Nếu người dùng để trang chờ quá lâu, hệ thống yêu cầu đăng nhập lại trước khi thực hiện cập nhật.
% \end{itemize} \\ \hline
% \end{tabular}
% \end{center}






% UC 04

\vspace{0.5cm}
\textbf{Đặc tả Use Case: Phân tích thành phần mỹ phẩm (Core Flow)}

\begin{center}
\renewcommand{\arraystretch}{1.4}
\begin{tabular}{|p{4cm}|p{10cm}|}
\hline
\rowcolor{blue!15}
\textbf{Tên Use Case} & \textbf{Phân tích thành phần mỹ phẩm (Hybrid Analysis)} \\ \hline
\textbf{Tác nhân} & Người dùng \\ \hline
\textbf{Mô tả} & Hệ thống tiếp nhận chuỗi INCI, thực hiện phân tích (Đối chiếu CSDL nội bộ và gọi AI Gemini cho hoạt chất lạ) để đưa ra cảnh báo an toàn và gợi ý sản phẩm phù hợp. \\ \hline
\textbf{Tiền điều kiện} & Người dùng đã đăng nhập và thiết lập hồ sơ loại da (\textit{SkinType}). \\ \hline
\textbf{Hậu điều kiện} & Kết quả được lưu vào Lịch sử phân tích (\textit{AnalysisHistory}); hệ thống tự động làm giàu CSDL nếu có hoạt chất mới từ AI. \\ \hline
\textbf{Luồng sự kiện chính} & 
\begin{enumerate}
    \item Người dùng dán chuỗi danh sách thành phần vào ô nhập liệu và nhấn "Phân tích".
    % \item Hệ thống thực hiện kiểm tra giới hạn lượt dùng (\textbf{Rate Limit 25 lần/24h}).
    \item Hệ thống chuẩn hóa dữ liệu: Tách mảng theo dấu phẩy, \textit{trim} khoảng trắng và chuyển về chữ thường (\textit{lowercase}).
    \item Hệ thống truy vấn CSDL nội bộ để tìm các thành phần và luật an toàn khớp với loại da của người dùng.
    \item Với các hoạt chất không có trong CSDL, hệ thống tự động gọi \textbf{Gemini Flash API} (model \texttt{gemini-flash-latest}) để đánh giá tính an toàn.
    \item \textbf{Tự động cập nhật (Auto-caching):} Các kết quả từ AI được hệ thống tự động lưu (\textit{upsert}) ngược lại vào bảng \textit{Ingredient} và \textit{IngredientRule} để phục vụ tức thời cho các lần sau.
    \item Hệ thống lưu chuỗi nhập liệu gốc vào bảng \textbf{AnalysisHistory} gắn với mã người dùng.
    \item Hệ thống phản hồi kết quả mã màu (Đỏ/Xanh/Xám) và kích hoạt thuật toán \textbf{Safety-first Recommendation} để lọc các sản phẩm an toàn tuyệt đối.
\end{enumerate} \\ \hline
\textbf{Luồng sự kiện phụ / Ngoại lệ} & 
\begin{itemize}
    \item \textit{Vượt quá giới hạn:} Nếu người dùng phân tích quá 25 lần trong 24 giờ, API trả lỗi 429 với thông báo: "Bạn đã đạt giới hạn phân tích trong ngày. Vui lòng thử lại sau, hoặc nâng cấp tài khoản Pro để có nhiều lượt hơn!".
    \item \textit{Lỗi kết nối AI:} Nếu Gemini API gặp sự cố hoặc hết hạn mức, hệ thống mặc định gán nhãn "Trung tính" (Neutral) để đảm bảo trải nghiệm không bị ngắt quãng.
    \item \textit{Dữ liệu rỗng:} Hệ thống báo lỗi "Vui lòng nhập danh sách thành phần" nếu ô văn bản trống.
\end{itemize} \\ \hline
\end{tabular}
\end{center}





%UC 05


\vspace{0.5cm}
\textbf{Đặc tả Use Case: Xem chi tiết mô tả thành phần}

\begin{center}
\renewcommand{\arraystretch}{1.4}
\begin{tabular}{|p{4cm}|p{10cm}|}
\hline
\rowcolor{blue!15}
\textbf{Tên Use Case} & \textbf{Xem chi tiết mô tả thành phần} \\ \hline
\textbf{Tác nhân} & Người dùng \\ \hline
\textbf{Mô tả} & Cung cấp thông tin chi tiết về công dụng, tính chất khoa học của một hóa chất cụ thể sau khi hệ thống đã thực hiện phân tích. \\ \hline
\textbf{Tiền điều kiện} & Người dùng đã thực hiện UC04 (Phân tích thành phần) và đang ở màn hình hiển thị kết quả. \\ \hline
\textbf{Hậu điều kiện} & Một cửa sổ phụ (Modal/Popup) hiển thị nội dung chi tiết về thành phần được chọn. \\ \hline
\textbf{Luồng sự kiện chính} & 
\begin{enumerate}
    \item Tại danh sách các thành phần đã được phân tích màu sắc, người dùng nhấn chuột (click) vào tên của một thành phần cụ thể.
    \item Frontend sử dụng trường \textit{description} đã có trong phần tử kết quả phân tích đang được chọn để hiển thị chi tiết, không phát sinh request tra cứu mới khi mở Modal.
    \item Hệ thống hiển thị một hộp thoại (Modal) nổi lên trên màn hình hiện tại.
    \item Nội dung hiển thị bao gồm: Tên thành phần (danh pháp quốc tế), công dụng chính, và các lưu ý đặc biệt (nếu có).
    \item Người dùng nhấn nút "Đóng" hoặc click ra vùng ngoài hộp thoại để quay lại màn hình kết quả.
\end{enumerate} \\ \hline
\textbf{Luồng sự kiện phụ / Ngoại lệ} & 
\begin{itemize}
    \item \textit{Dữ liệu trống:} Nếu thành phần chưa có mô tả chi tiết trong CSDL, hệ thống hiển thị thông báo: "Thông tin đang được cập nhật".
    \item \textit{Không có kết quả phân tích:} Modal chỉ được mở khi người dùng đã chọn một thành phần trong mảng kết quả trả về.
\end{itemize} \\ \hline
\end{tabular}
\end{center}



\textbf{Đặc tả Use Case: Xem gợi ý sản phẩm}

\begin{center}
\renewcommand{\arraystretch}{1.4}
\begin{tabular}{|p{4cm}|p{10cm}|}
\hline
\rowcolor{blue!15}
\textbf{Tên Use Case} & \textbf{Xem gợi ý sản phẩm} \\ \hline
\textbf{Tác nhân} & Người dùng \\ \hline
\textbf{Mô tả} & Hệ thống tự động đề xuất sản phẩm trong danh mục quản lý, lọc theo loại da người dùng và hiển thị sau khi hoàn tất phân tích. \\ \hline
\textbf{Tiền điều kiện} & Người dùng đã thực hiện thành công UC04 (Phân tích thành phần). \\ \hline
\textbf{Hậu điều kiện} & Danh sách các sản phẩm phù hợp được hiển thị trực quan ở cuối trang kết quả. \\ \hline
\textbf{Luồng sự kiện chính} & 
\begin{enumerate}
    \item Sau khi hệ thống hoàn tất hiển thị kết quả phân tích màu sắc ở UC04, thuật toán gợi ý được kích hoạt ngầm.
    \item Hệ thống truy vấn danh sách sản phẩm từ bảng \textit{Products}.
    \item Đối với mỗi sản phẩm, hệ thống đối chiếu chuỗi \textit{inci\_list} của sản phẩm đó với loại da của người dùng hiện tại.
    \item Hệ thống lọc ra các sản phẩm \textbf{không chứa bất kỳ thành phần nào bị đánh giá là "Bad" (Màu đỏ)} đối với loại da đó.
    \item Hệ thống chấm điểm sản phẩm an toàn theo số thành phần "Good", lấy nhóm tối đa 6 sản phẩm điểm cao nhất, xáo trộn ngẫu nhiên và hiển thị tối đa 3 sản phẩm.
    \item Người dùng xem các thông tin gồm: Tên sản phẩm, Thương hiệu và Hình ảnh minh họa.
\end{enumerate} \\ \hline
\textbf{Luồng sự kiện phụ / Ngoại lệ} & 
\begin{itemize}
    \item \textit{Không tìm thấy sản phẩm phù hợp:} Nếu không có sản phẩm thỏa điều kiện an toàn cho loại da của người dùng, giao diện hiển thị trạng thái trống thông báo chưa tìm thấy sản phẩm phù hợp.
    \item \textit{Dữ liệu sản phẩm trống:} Nếu bảng \textit{Products} chưa có dữ liệu, khu vực gợi ý vẫn hiển thị trạng thái trống với nội dung chưa tìm thấy sản phẩm phù hợp.
\end{itemize} \\ \hline
\end{tabular}
\end{center}


% NOTE: BỎ USE CASE LƯU LỊCH SỬ PHÂN TÍCH
%UC 07
% \textbf{Đặc tả Use Case: Lưu lịch sử phân tích}

% \begin{center}
% \renewcommand{\arraystretch}{1.4}
% \begin{tabular}{|p{4cm}|p{10cm}|}
% \hline
% \rowcolor{blue!15}
% \textbf{Tên Use Case} & \textbf{Lưu lịch sử phân tích} \\ \hline
% \textbf{Tác nhân} & Người dùng \\ \hline
% \textbf{Mô tả} & Hệ thống hỗ trợ người dùng ghi lại kết quả các lần kiểm tra thành phần mỹ phẩm để tra cứu và theo dõi tình trạng da lâu dài. \\ \hline
% \textbf{Tiền điều kiện} & Người dùng đã đăng nhập thành công và vừa hoàn thành quy trình phân tích thành phần (UC04). \\ \hline
% \textbf{Hậu điều kiện} & Một bản ghi mới chứa thông tin phân tích được lưu trữ vào bảng \textit{History} trong cơ sở dữ liệu. \\ \hline
% \textbf{Luồng sự kiện chính} & 
% \begin{enumerate}
%     \item Tại màn hình hiển thị kết quả phân tích, người dùng nhấn vào nút "Lưu vào lịch sử".
%     \item Hệ thống tự động thu thập các thông tin bao gồm: Mã người dùng (User ID), danh sách thành phần đã nhập và kết quả phân loại (Đỏ/Xanh/Xám).
%     \item Hệ thống thực hiện ghi nhận dữ liệu vào cơ sở dữ liệu cùng với dấu thời gian thực tế.
%     \item Hệ thống hiển thị thông báo "Lưu kết quả thành công" và cập nhật danh sách trong mục lịch sử cá nhân.
% \end{enumerate} \\ \hline
% \textbf{Luồng sự kiện phụ / Ngoại lệ} & 
% \begin{itemize}
%     \item \textit{Phiên đăng nhập hết hạn:} Nếu người dùng để trang chờ quá lâu trước khi nhấn lưu, hệ thống yêu cầu đăng nhập lại để xác thực danh tính.
%     \item \textit{Lỗi kết nối CSDL:} Nếu quá trình lưu gặp sự cố kỹ thuật, hệ thống hiển thị thông báo "Không thể lưu kết quả lúc này, vui lòng thử lại sau" để người dùng biết.
% \end{itemize} \\ \hline
% \end{tabular}
% \end{center}



\vspace{0.5cm}
\textbf{Đặc tả Use Case: Quản lý lịch sử phân tích}

\begin{center}
\renewcommand{\arraystretch}{1.4}
\begin{tabular}{|p{4cm}|p{10cm}|}
\hline
\rowcolor{blue!15}
\textbf{Tên Use Case} & \textbf{Quản lý lịch sử phân tích} \\ \hline
\textbf{Tác nhân} & Người dùng \\ \hline
\textbf{Mô tả} & Người dùng có thể xem lại danh sách các bảng thành phần đã từng phân tích, thực hiện lệnh phân tích lại (Re-analyze), hoặc xóa bỏ lịch sử. \\ \hline
\textbf{Tiền điều kiện} & Người dùng đã đăng nhập thành công. \\ \hline
\textbf{Hậu điều kiện} & Giao diện hiển thị danh sách lịch sử. Nếu có thao tác xóa, dữ liệu trong bảng \textit{AnalysisHistory} sẽ bị xóa tương ứng. \\ \hline
\textbf{Luồng sự kiện chính} & 
\begin{enumerate}
    \item Người dùng truy cập trang "Lịch sử" (/history).
    \item Hệ thống gọi API \texttt{GET /api/v1/history}, lấy danh sách lịch sử từ CSDL sắp xếp theo thời gian mới nhất (DESC).
    \item Hệ thống hiển thị danh sách dưới dạng các thẻ (Cards).
    \item \textbf{Nếu user nhấn "Phân tích lại":} Hệ thống chuyển hướng sang trang Analysis với chuỗi INCI cũ trong query \texttt{inci}; chuỗi được điền sẵn để người dùng thực hiện phân tích.
    \item \textbf{Nếu user nhấn "Xóa 1 thẻ":} Bật xác nhận $\rightarrow$ Gọi API \texttt{DELETE /api/v1/history/:id} $\rightarrow$ Cập nhật giao diện.
    \item \textbf{Nếu user nhấn "Xóa tất cả":} Bật cảnh báo $\rightarrow$ Gọi API \texttt{DELETE /api/v1/history} $\rightarrow$ Làm sạch toàn bộ dữ liệu trên giao diện.
\end{enumerate} \\ \hline
\textbf{Luồng sự kiện phụ / Ngoại lệ} & 
\begin{itemize}
    \item \textit{Chưa có lịch sử:} Nếu DB trả về mảng rỗng, hệ thống hiển thị màn hình Trống (Empty State) kèm nút "Phân tích ngay".
    \item \textit{Hủy thao tác xóa:} Khi hộp thoại Confirm hiện lên, nếu User chọn Cancel, hệ thống không làm gì cả.
\end{itemize} \\ \hline
\end{tabular}
\end{center}






\vspace{0.5cm}
\textbf{Đặc tả Use Case: Trích xuất thành phần bằng OCR}

\begin{center}
\renewcommand{\arraystretch}{1.4}
\begin{tabular}{|p{4cm}|p{10cm}|}
\hline
\rowcolor{blue!15}
\textbf{Tên Use Case} & \textbf{Trích xuất thành phần bằng OCR} \\ \hline
\textbf{Tác nhân} & Người dùng \\ \hline
\textbf{Mô tả} & Cho phép người dùng tải lên ảnh chụp bảng thành phần, hệ thống nhận dạng ký tự (OCR) và tự động điền chuỗi INCI vào ô phân tích. \\ \hline
\textbf{Tiền điều kiện} & Người dùng đã đăng nhập và đang ở màn hình phân tích. \\ \hline
\textbf{Hậu điều kiện} & Chuỗi INCI trích xuất được hiển thị trong textarea, sẵn sàng để phân tích. \\ \hline
\textbf{Luồng sự kiện chính} & 
\begin{enumerate}
    \item Tại màn hình Phân tích, người dùng click vào vùng “Trích xuất từ hình ảnh” hoặc nút tải lên.
    \item Hệ thống mở hộp thoại chọn file (chấp nhận JPG/PNG).
    \item Người dùng chọn ảnh chụp bảng thành phần, hệ thống hiển thị ảnh preview kèm trạng thái “Đang phân tích”.
    \item Frontend gửi file ảnh đến API \texttt{POST /api/ocr/ingredients} (dùng FormData).
    \item Backend chuyển ảnh thành base64, gọi OCR.space API để nhận dạng văn bản.
    \item Kết quả thô được xử lý bởi \texttt{ingredientsExtractor.ts}: tìm từ khóa “ingredients”, tách theo dấu phẩy/chấm phẩy, chuẩn hóa về chữ thường.
    \item Frontend nhận kết quả và tự động gán vào ô nhập INCI.
    \item Người dùng có thể chỉnh sửa chuỗi trước khi nhấn “Phân tích”.
\end{enumerate} \\ \hline
\textbf{Luồng sự kiện phụ / Ngoại lệ} & 
\begin{itemize}
    \item \textit{Không tìm thấy thành phần:} OCR trả về rỗng hoặc không tìm được từ khóa → thông báo lỗi “Không tìm thấy thành phần trong ảnh”.
    \item \textit{Định dạng ảnh không hỗ trợ:} Hệ thống báo lỗi ngay khi chọn file sai.
    \item \textit{Ảnh quá lớn:} Giới hạn dung lượng ở phía client và server.
\end{itemize} \\ \hline
\end{tabular}
\end{center}











\vspace{0.5cm}
\textbf{Đặc tả Use Case: Báo cáo phân loại sai thành phần}

\begin{center}
\renewcommand{\arraystretch}{1.4}
\begin{tabular}{|p{4cm}|p{10cm}|}
\hline
\rowcolor{blue!15}
\textbf{Tên Use Case} & \textbf{Báo cáo phân loại sai thành phần} \\ \hline
\textbf{Tác nhân} & Người dùng \\ \hline
\textbf{Mô tả} & Người dùng gửi yêu cầu xem xét lại đánh giá an toàn của một thành phần, đề xuất phân loại mới kèm lý do và bằng chứng. \\ \hline
\textbf{Tiền điều kiện} & Người dùng đã đăng nhập, đã thực hiện phân tích (UC04) và đang xem chi tiết một thành phần. \\ \hline
\textbf{Hậu điều kiện} & Một báo cáo mới được lưu vào bảng \texttt{IngredientReport} với trạng thái \texttt{PENDING}. \\ \hline
\textbf{Luồng sự kiện chính} & 
\begin{enumerate}
    \item Tại modal chi tiết thành phần, người dùng nhấn nút “Báo cáo”.
    \item Hệ thống hiển thị form (BM8) với: phân loại đề xuất (GOOD/BAD/NEUTRAL), ô nhập lý do, ô link dẫn chứng (tùy chọn).
    \item Người dùng chọn phân loại mong muốn, nhập lý do ($\ge$ 20 ký tự).
    \item Nhấn “Gửi báo cáo”.
    \item Frontend gọi API \texttt{POST /api/v1/reports} với: ingredientId (tìm qua tên INCI), skinType của người dùng, reportedEffect, reason, evidenceUrl.
    \item Backend lưu báo cáo, trả về thành công.
    \item Hệ thống thông báo “Báo cáo thành công! Cảm ơn bạn đã đóng góp.”
\end{enumerate} \\ \hline
\textbf{Luồng sự kiện phụ / Ngoại lệ} & 
\begin{itemize}
    \item \textit{Lý do quá ngắn:} Client kiểm tra độ dài < 20 → cảnh báo, không gửi lên server.
    \item \textit{Thành phần không tồn tại:} Khi search tên INCI không có trong DB → thông báo lỗi, đề nghị admin thêm thành phần trước.
    \item \textit{Lỗi server:} Hiển thị thông báo lỗi màu đỏ, cho phép thử lại.
\end{itemize} \\ \hline
\end{tabular}
\end{center}





\vspace{0.5cm}
\textbf{Đặc tả Use Case: Bình chọn báo cáo cộng đồng}

\begin{center}
\renewcommand{\arraystretch}{1.4}
\begin{tabular}{|p{4cm}|p{10cm}|}
\hline
\rowcolor{blue!15}
\textbf{Tên Use Case} & \textbf{Bình chọn báo cáo cộng đồng} \\ \hline
\textbf{Tác nhân} & Người dùng \\ \hline
\textbf{Mô tả} & Người dùng có thể bình chọn (UP/DOWN) cho các báo cáo phân loại đang chờ duyệt để thể hiện sự đồng thuận, giúp Admin ưu tiên xử lý. \\ \hline
\textbf{Tiền điều kiện} & Người dùng đã đăng nhập, truy cập trang “Cộng đồng” (/community/reports). \\ \hline
\textbf{Hậu điều kiện} & Số phiếu của báo cáo được cập nhật; người dùng có thể thay đổi hoặc hủy phiếu. \\ \hline
\textbf{Luồng sự kiện chính} & 
\begin{enumerate}
    \item Trang tải danh sách báo cáo đang chờ (PENDING) kèm tổng điểm (up - down) và trạng thái bình chọn của người dùng hiện tại.
    \item Người dùng nhấn nút UP hoặc DOWN trên một báo cáo.
    \item Frontend gọi API \texttt{POST /api/v1/reports/vote} với reportId và voteType (UP/DOWN).
    \item Nếu người dùng đã bình chọn trùng loại → phiếu được hủy (toggle off). Nếu khác loại → cập nhật thành loại mới.
    \item Backend cập nhật bảng \texttt{ReportVote}, trả về số liệu up/down mới và trạng thái bình chọn của user.
    \item Frontend cập nhật giao diện ngay lập tức (optimistic update).
\end{enumerate} \\ \hline
\textbf{Luồng sự kiện phụ / Ngoại lệ} & 
\begin{itemize}
    \item \textit{Lỗi mạng:} Nếu API thất bại, giao diện khôi phục trạng thái trước đó để không mất dữ liệu tạm thời.
    \item \textit{Chưa đăng nhập:} Hệ thống chặn và yêu cầu đăng nhập.
\end{itemize} \\ \hline
\end{tabular}
\end{center}








\vspace{0.5cm}
\textbf{Đặc tả Use Case: Xem thông báo}

\begin{center}
\renewcommand{\arraystretch}{1.4}
\begin{tabular}{|p{4cm}|p{10cm}|}
\hline
\rowcolor{blue!15}
\textbf{Tên Use Case} & \textbf{Xem thông báo} \\ \hline
\textbf{Tác nhân} & Người dùng \\ \hline
\textbf{Mô tả} & Người dùng xem các thông báo hệ thống (báo cáo được duyệt/từ chối, tin nhắn từ Admin), đánh dấu đã đọc và điều hướng đến liên kết liên quan. \\ \hline
\textbf{Tiền điều kiện} & Người dùng đã đăng nhập. \\ \hline
\textbf{Hậu điều kiện} & Các thông báo được đánh dấu \texttt{isRead = true}. \\ \hline
\textbf{Luồng sự kiện chính} & 
\begin{enumerate}
    \item Trên thanh điều hướng, người dùng bấm vào biểu tượng chuông (NotificationBell).
    \item Hệ thống gọi API \texttt{GET /api/v1/notifications} khi người dùng đăng nhập, mỗi lần mở dropdown và định kỳ 30 giây khi dropdown đang đóng.
    \item Danh sách thông báo (tối đa 50 gần nhất) hiển thị kèm: tiêu đề, nội dung, thời gian, trạng thái đọc/chưa đọc.
    \item Thông báo chưa đọc được tô nền nhạt và có chấm đỏ.
    \item Người dùng có thể:
        \begin{itemize}
            \item Bấm vào thông báo → đánh dấu đã đọc qua API \texttt{PATCH /api/v1/notifications/:id/read} và nếu có liên kết sẽ chuyển hướng.
            \item Bấm “Đánh dấu tất cả đã đọc” → gọi \texttt{PATCH /api/v1/notifications/read-all}.
        \end{itemize}
    \item Số lượng thông báo chưa đọc hiển thị badge trên biểu tượng chuông.
\end{enumerate} \\ \hline
\textbf{Luồng sự kiện phụ / Ngoại lệ} & 
\begin{itemize}
    \item \textit{Không có thông báo:} Hiển thị giao diện trống rỗng kèm icon.
    \item \textit{Lỗi tải:} Dropdown vẫn hiển thị dữ liệu cũ, badge không đổi; lần sau mở lại sẽ tự động thử lại.
\end{itemize} \\ \hline
\end{tabular}
\end{center}


\textbf{*Nhóm Use Case dành cho Quản trị viên (Admin)}


\vspace{0.5cm}
\textbf{Đặc tả Use Case: Quản lý danh sách thành phần}

\begin{center}
\renewcommand{\arraystretch}{1.4}
\begin{tabular}{|p{4cm}|p{10cm}|}
\hline
\rowcolor{blue!15}
\textbf{Tên Use Case} & \textbf{Quản lý danh sách thành phần} \\ \hline
\textbf{Tác nhân} & Quản trị viên (Admin) \\ \hline
\textbf{Mô tả} & Cho phép Admin thực hiện các thao tác quản trị dữ liệu (Thêm, Sửa, Xóa) đối với danh mục các thành phần hóa học trong cơ sở dữ liệu của hệ thống. \\ \hline
\textbf{Tiền điều kiện} & Tác nhân đã đăng nhập vào hệ thống với quyền Quản trị viên. \\ \hline
\textbf{Hậu điều kiện} & Cơ sở dữ liệu bảng \textit{Ingredients} được cập nhật thông tin mới nhất. \\ \hline
\textbf{Luồng sự kiện chính} & 
\begin{enumerate}
    \item Admin truy cập vào menu "Quản lý thành phần" trên trang quản trị.
    \item Hệ thống hiển thị danh sách toàn bộ các thành phần hiện có kèm theo công cụ tìm kiếm và bộ lọc.
    \item \textbf{Trường hợp Thêm mới:} Admin nhấn "Thêm", nhập tên thành phần và mô tả công dụng vào biểu mẫu, sau đó nhấn "Lưu".
    \item \textbf{Trường hợp Chỉnh sửa:} Admin chọn một thành phần, thay đổi nội dung mô tả và nhấn "Cập nhật".
    \item \textbf{Trường hợp Xóa:} Admin nhấn vào biểu tượng xóa tại dòng tương ứng của thành phần cần loại bỏ.
    \item Hệ thống kiểm tra tính hợp lệ của dữ liệu (không để trống tên thành phần).
    \item Hệ thống thực hiện cập nhật vào CSDL và thông báo "Thao tác thành công".
\end{enumerate} \\ \hline
\textbf{Luồng sự kiện phụ / Ngoại lệ} & 
\begin{itemize}
    \item \textit{Trùng tên thành phần:} Nếu Admin thêm một thành phần đã tồn tại trong CSDL, hệ thống báo lỗi "Thành phần này đã có trong hệ thống".
    \item \textit{Xác nhận xóa:} Khi Admin nhấn xóa, hệ thống hiển thị cảnh báo "Bạn có chắc chắn muốn xóa thành phần này không?" để tránh thao tác nhầm.
    \item \textit{Ràng buộc dữ liệu:} Nếu thành phần đang được liên kết với các "Luật phân tích" (UC07), hệ thống sẽ yêu cầu xóa các luật liên quan trước khi xóa thành phần đó để đảm bảo tính toàn vẹn dữ liệu.
\end{itemize} \\ \hline
\end{tabular}
\end{center}


% UC 07


\vspace{0.5cm}
\textbf{Đặc tả Use Case: Quản lý luật phân tích (Rules)}

\begin{center}
\renewcommand{\arraystretch}{1.4}
\begin{tabular}{|p{4cm}|p{10cm}|}
\hline
\rowcolor{blue!15}
\textbf{Tên Use Case} & \textbf{Quản lý luật phân tích (Rules)} \\ \hline
\textbf{Tác nhân} & Quản trị viên (Admin) \\ \hline
\textbf{Mô tả} & Cho phép Admin thiết lập, điều chỉnh hoặc xóa bỏ các quy tắc ánh xạ giữa một thành phần hóa học và một loại da cụ thể để đưa ra đánh giá (Good/Bad/Neutral). \\ \hline
\textbf{Tiền điều kiện} & Admin đã đăng nhập thành công; các thành phần hóa học liên quan đã tồn tại trong bảng \textit{Ingredients}. \\ \hline
\textbf{Hậu điều kiện} & Tập luật trong bảng \textit{Ingredient\_Rules} được cập nhật, làm thay đổi trực tiếp kết quả phân tích của người dùng ở UC04. \\ \hline
\textbf{Luồng sự kiện chính} & 
\begin{enumerate}
    \item Admin truy cập vào menu "Quản lý Luật phân tích".
    \item Hệ thống hiển thị danh sách các luật hiện có (Thành phần - Loại da - Đánh giá).
    \item \textbf{Thêm luật mới:} Admin chọn tên thành phần từ danh sách có sẵn, chọn loại da tương ứng (Bình thường, Dầu, Khô, Nhạy cảm, Hỗn hợp) và chọn mức độ đánh giá (Xanh, Đỏ, Xám).
    \item Admin nhấn "Lưu luật".
    \item Hệ thống kiểm tra quy định (QĐ 4): Mỗi cặp (Thành phần, Loại da) chỉ được có một mức độ đánh giá duy nhất.
    \item Hệ thống cập nhật dữ liệu vào bảng \textit{Ingredient\_Rules} và thông báo thành công.
\end{enumerate} \\ \hline
\textbf{Luồng sự kiện phụ / Ngoại lệ} & 
\begin{itemize}
    \item \textit{Luật đã tồn tại:} Nếu Admin thêm một luật cho cặp Thành phần - Loại da đã có trong hệ thống, hệ thống báo lỗi "Luật này đã tồn tại" và gợi ý sử dụng chức năng "Sửa".
    \item \textit{Sửa đổi luật:} Admin có thể thay đổi mức độ đánh giá của một luật hiện có (ví dụ: chuyển từ Neutral sang Bad) và nhấn "Cập nhật".
    \item \textit{Xóa luật:} Khi xóa một luật, hệ thống sẽ gỡ bỏ quy tắc đánh giá đó. Sau khi xóa, nếu người dùng phân tích thành phần này, hệ thống sẽ mặc định trả về màu Xám (Trung tính).
\end{itemize} \\ \hline
\end{tabular}
\end{center}




%UC 08


\vspace{0.5cm}



%UC 09


\vspace{0.5cm}
\textbf{Đặc tả Use Case: Quản lý danh mục sản phẩm}

\begin{center}
\renewcommand{\arraystretch}{1.4}
\begin{tabular}{|p{4cm}|p{10cm}|}
\hline
\rowcolor{blue!15}
\textbf{Tên Use Case} & \textbf{Quản lý danh mục sản phẩm} \\ \hline
\textbf{Tác nhân} & Quản trị viên (Admin) \\ \hline
\textbf{Mô tả} & Cho phép Admin quản lý thông tin các sản phẩm mỹ phẩm (Thêm, Sửa, Xóa) bao gồm tên, thương hiệu, hình ảnh và danh sách thành phần để hệ thống thực hiện gợi ý cho người dùng. \\ \hline
\textbf{Tiền điều kiện} & Admin đã đăng nhập thành công vào hệ thống. \\ \hline
\textbf{Hậu điều kiện} & Thông tin sản phẩm trong bảng \textit{Products} được cập nhật chính xác. \\ \hline
\textbf{Luồng sự kiện chính} & 
\begin{enumerate}
    \item Admin truy cập vào menu "Quản lý sản phẩm" trên trang quản trị.
    \item Hệ thống hiển thị danh sách các sản phẩm hiện có cùng các tùy chọn chức năng.
    \item \textbf{Thêm sản phẩm:} Admin nhập các thông tin: Tên sản phẩm, Thương hiệu, Link hình ảnh và chuỗi INCI (danh sách thành phần cách nhau bằng dấu phẩy) vào biểu mẫu.
    \item \textbf{Chỉnh sửa sản phẩm:} Admin chọn sản phẩm cần sửa, cập nhật lại các thông tin thay đổi và nhấn "Lưu".
    \item \textbf{Xóa sản phẩm:} Admin chọn sản phẩm không còn phù hợp và nhấn "Xóa".
    \item Hệ thống kiểm tra tính hợp lệ của dữ liệu đầu vào.
    \item Hệ thống thực thi cập nhật vào CSDL và thông báo "Thao tác thành công".
\end{enumerate} \\ \hline
\textbf{Luồng sự kiện phụ / Ngoại lệ} & 
\begin{itemize}
    \item \textit{Dữ liệu không hợp lệ:} Nếu Admin để trống Tên sản phẩm hoặc Chuỗi thành phần (INCI), hệ thống sẽ báo lỗi và yêu cầu nhập bổ sung.
    \item \textit{Định dạng INCI sai:} Nếu chuỗi thành phần nhập vào không đúng định dạng (thiếu dấu phẩy), hệ thống hiển thị cảnh báo để Admin kiểm tra lại nhằm đảm bảo thuật toán gợi ý ở UC08 hoạt động chính xác.
    \item \textit{Xác nhận xóa:} Hệ thống hiển thị hộp thoại xác nhận trước khi xóa vĩnh viễn dữ liệu sản phẩm.
\end{itemize} \\ \hline
\end{tabular}
\end{center}



\vspace{0.5cm}

\textbf{Đặc tả Use Case: Quản lý người dùng}

\begin{center}
\renewcommand{\arraystretch}{1.4}
\begin{tabular}{|p{4cm}|p{10cm}|}
\hline
\rowcolor{blue!15}
\textbf{Tên Use Case} & \textbf{Quản lý người dùng (User Management)} \\ \hline
\textbf{Tác nhân} & Quản trị viên (Admin) \\ \hline
\textbf{Mô tả} & Quản trị viên kiểm soát danh sách người dùng, thực hiện khóa/mở khóa tài khoản hoặc xóa vĩnh viễn dữ liệu người dùng vi phạm để bảo đảm an ninh hệ thống. \\ \hline
\textbf{Tiền điều kiện} & Tác nhân đã đăng nhập với vai trò "ADMIN". \\ \hline
\textbf{Hậu điều kiện} & Trạng thái tài khoản được cập nhật; nếu xóa, toàn bộ lịch sử phân tích của người dùng đó cũng sẽ bị loại bỏ khỏi CSDL (\textit{Cascade Delete}). \\ \hline
\textbf{Luồng sự kiện chính} & 
\begin{enumerate}
    \item Admin truy cập vào menu "User Management".
    \item Hệ thống truy vấn CSDL và hiển thị danh sách người dùng (không bao gồm mật khẩu băm) kèm trạng thái hoạt động (\textit{isActive}).
    \item Admin lựa chọn một trong các hành động:
        \begin{itemize}
            \item \textbf{Khóa/Mở khóa:} Đổi trạng thái từ \textit{Active} sang \textit{Locked} và ngược lại.
            \item \textbf{Xóa tài khoản:} Loại bỏ hoàn toàn bản ghi người dùng.
        \end{itemize}
    \item Hệ thống hiển thị \textit{Confirmation Modal} để ngăn chặn các thao tác sai lầm.
    \item Admin nhấn "Xác nhận".
    \item Hệ thống thực thi lệnh cập nhật qua Prisma:
        \begin{itemize}
            \item Nếu khóa: Cập nhật trường \texttt{isActive = false}.
            \item Nếu xóa: Thực hiện lệnh \texttt{delete} kèm cơ chế xóa dây chuyền lịch sử.
        \end{itemize}
    \item Hệ thống thông báo thành công và cập nhật lại danh sách trên giao diện.
\end{enumerate} \\ \hline
\textbf{Luồng sự kiện phụ / Ngoại lệ} & 
\begin{itemize}
    \item \textit{Bảo vệ trên giao diện:} Các nút Khóa/Mở khóa/Xóa không được hiển thị đối với các tài khoản có vai trò \texttt{ADMIN}; thao tác với tài khoản người dùng được xác nhận trước khi gửi yêu cầu.
    \item \textit{Hệ quả của việc Khóa:} Người dùng có tài khoản bị khóa (\textit{isActive: false}) sẽ lập tức bị từ chối ở bước Đăng nhập với lỗi: "Tài khoản của bạn đã bị khóa bởi Quản trị viên."
    \item \textit{Hủy thao tác:} Nếu Admin nhấn "Hủy", hệ thống đóng modal và không thay đổi dữ liệu.
\end{itemize} \\ \hline
\end{tabular}
\end{center}


\vspace{0.5cm}
\vspace{0.5cm}
\textbf{Đặc tả Use Case: Xem báo cáo thống kê (Admin Dashboard)}

\begin{center}
\renewcommand{\arraystretch}{1.4}
\begin{tabular}{|p{4cm}|p{10cm}|}
\hline
\rowcolor{blue!15}
\textbf{Tên Use Case} & \textbf{Xem báo cáo thống kê} \\ \hline
\textbf{Tác nhân} & Quản trị viên (Admin) \\ \hline
\textbf{Mô tả} & Hệ thống tự động tổng hợp các chỉ số vận hành từ CSDL và trực quan hóa dưới dạng biểu đồ để Admin theo dõi sự tăng trưởng người dùng và phân bố nhân khẩu học (loại da). \\ \hline
\textbf{Tiền điều kiện} & Admin đã đăng nhập thành công và truy cập vào đường dẫn \textit{/admin/reports}. \\ \hline
\textbf{Hậu điều kiện} & Toàn bộ chỉ số thực tế được hiển thị; không làm thay đổi dữ liệu hệ thống. \\ \hline
\textbf{Luồng sự kiện chính} & 
\begin{enumerate}
    \item Admin truy cập vào mục "Reports" trên thanh điều hướng.
    \item Giao diện gửi yêu cầu \texttt{GET /api/v1/admin/reports} kèm Token xác thực.
    \item \textbf{Xử lý tại Backend (Service):}
        \begin{itemize}
            \item Sử dụng hàm \texttt{count()} để thống kê tổng số Người dùng (\textit{User}) và tổng số lượt Phân tích (\textit{AnalysisHistory}).
            \item Sử dụng hàm \texttt{groupBy()} để đếm số lượng người dùng theo từng nhóm \textit{SkinType} (Dầu, Khô, Hỗn hợp, Nhạy cảm).
        \end{itemize}
    \item Backend phản hồi dữ liệu tổng hợp dưới dạng JSON.
    \item \textbf{Xử lý tại Frontend (View):}
        \begin{itemize}
            \item Hiển thị các chỉ số tổng quát (Summary Cards) kèm hiệu ứng \textit{hover}.
            \item Sử dụng thư viện \textbf{Recharts} để vẽ biểu đồ tròn (\textit{PieChart}) thể hiện tỷ lệ phân bố loại da.
            \item Hiển thị chú thích (\textit{Legend}) và gợi ý thông tin (\textit{Tooltip}) khi Admin di chuyển chuột vào biểu đồ.
        \end{itemize}
\end{enumerate} \\ \hline
\textbf{Luồng sự kiện phụ / Ngoại lệ} & 
\begin{itemize}
    \item \textit{Dữ liệu mới (Trống):} Nếu hệ thống chưa có dữ liệu, các thẻ tổng số sẽ hiển thị giá trị "0" và biểu đồ sẽ không hiển thị các lát cắt, kèm thông báo "Chưa có dữ liệu thống kê".
    \item \textit{Lỗi đồng bộ:} Nếu API không phản hồi, hệ thống hiển thị thông báo lỗi màu đỏ và nút "Thử lại".
\end{itemize} \\ \hline
\end{tabular}
\end{center}




\vspace{0.5cm}
\textbf{Đặc tả Use Case: Quản lý dữ liệu qua tệp Excel (Import/Export)}

\begin{center}
\renewcommand{\arraystretch}{1.4}
\begin{tabular}{|p{4cm}|p{10cm}|}
\hline
\rowcolor{blue!15}
\textbf{Tên Use Case} & \textbf{Quản lý dữ liệu qua tệp Excel} \\ \hline
\textbf{Tác nhân} & Quản trị viên (Admin) \\ \hline
\textbf{Mô tả} & Cho phép Admin xuất dữ liệu hiện có ra tệp Excel để lưu trữ hoặc nhập dữ liệu hàng loạt từ tệp Excel vào hệ thống để cập nhật tri thức nhanh chóng. \\ \hline
\textbf{Tiền điều kiện} & Tác nhân đã đăng nhập với vai trò "ADMIN". \\ \hline
\textbf{Hậu điều kiện} & Tệp dữ liệu được tải về máy cục bộ (Export) hoặc Cơ sở dữ liệu được cập nhật hàng loạt (Import). \\ \hline
\textbf{Luồng sự kiện chính (Import)} & 
\begin{enumerate}
    \item Admin truy cập trang "Import / Export".
    \item Admin chọn danh mục cần nhập (Thành phần, Quy tắc hoặc Sản phẩm).
    \item Admin chọn tệp tin từ máy tính (định dạng \textit{.xlsx} hoặc \textit{.xls}) và nhấn "Upload".
    \item Hệ thống đọc tệp dữ liệu vào bộ nhớ đệm (\textit{Memory Storage}).
    \item Hệ thống duyệt từng dòng dữ liệu:
        \begin{itemize}
            \item \textbf{Thành phần:} Thực hiện \textit{Upsert} (Cập nhật nếu trùng tên, tạo mới nếu chưa có).
            \item \textbf{Sản phẩm:} Bóc tách chuỗi INCI trong tệp Excel để tự động liên kết với các mã thành phần tương ứng.
        \end{itemize}
    \item Hệ thống tổng hợp kết quả: Số dòng thành công, Số dòng lỗi, Số dòng cập nhật.
    \item Hệ thống hiển thị báo cáo chi tiết cho Admin.
\end{enumerate} \\ \hline
\textbf{Luồng sự kiện chính (Export)} & 
\begin{enumerate}
    \item Admin nhấn nút "Export" tại danh mục tương ứng.
    \item Hệ thống truy vấn toàn bộ dữ liệu từ bảng liên quan.
    \item Hệ thống sử dụng thư viện \textit{XLSX} để chuyển đổi dữ liệu JSON thành bảng tính, tự động căn chỉnh độ rộng cột.
    \item Hệ thống trả về luồng dữ liệu (\textit{Buffer}) và trình duyệt tự động tải xuống tệp tin.
\end{enumerate} \\ \hline
\textbf{Luồng sự kiện phụ / Ngoại lệ} & 
\begin{itemize}
    \item \textit{Sai định dạng tệp:} Nếu Admin tải lên file không phải Excel, hệ thống báo lỗi: "Only Excel files (.xlsx, .xls) are allowed".
    \item \textit{Thiếu cột bắt buộc:} Nếu file Excel thiếu các cột định danh (như tên thành phần), hệ thống sẽ bỏ qua dòng đó và ghi nhận vào danh sách lỗi.
    \item \textit{Lỗi dữ liệu logic:} Nếu nhập Quy tắc cho một hoạt chất không tồn tại, hệ thống báo lỗi chi tiết theo từng dòng để Admin điều chỉnh.
\end{itemize} \\ \hline
\end{tabular}
\end{center}





\vspace{0.5cm}
\textbf{Đặc tả Use Case: Duyệt báo cáo cộng đồng (Admin)}

\begin{center}
\renewcommand{\arraystretch}{1.4}
\begin{tabular}{|p{4cm}|p{10cm}|}
\hline
\rowcolor{blue!15}
\textbf{Tên Use Case} & \textbf{Duyệt báo cáo cộng đồng} \\ \hline
\textbf{Tác nhân} & Quản trị viên (Admin) \\ \hline
\textbf{Mô tả} & Admin xem xét các báo cáo phân loại đang chờ, có thể chấp nhận (APPROVED) để cập nhật luật phân tích, hoặc từ chối (REJECTED). Hệ thống tự động gửi thông báo đến người gửi báo cáo. \\ \hline
\textbf{Tiền điều kiện} & Admin đã đăng nhập, có quyền ADMIN. \\ \hline
\textbf{Hậu điều kiện} & Trạng thái báo cáo thay đổi; nếu APPROVED, luật \texttt{IngredientRule} tương ứng được tạo/cập nhật; một thông báo được gửi đến người báo cáo. \\ \hline
\textbf{Luồng sự kiện chính} & 
\begin{enumerate}
    \item Admin truy cập “Admin” → “BC Cộng đồng” (/admin/community-reports).
    \item Hệ thống hiển thị bảng các báo cáo PENDING, sắp xếp theo số phiếu (nổi bật nhất) hoặc mới nhất.
    \item Admin xem chi tiết từng báo cáo (thành phần, đề xuất, lý do, dẫn chứng, số up/down).
    \item Admin nhấn nút ``duyệt'' hoặc ``từ chối'' → mở modal xác nhận.
    \item Tại modal, Admin có thể nhập ghi chú (tuỳ chọn), sau đó xác nhận.
    \item Frontend gọi API \texttt{POST /api/v1/reports/resolve} với: reportId, status (APPROVED/REJECTED), adminNote.
    \item Backend cập nhật báo cáo, nếu APPROVED thì gọi \texttt{createOrUpdateRule()} để thiết lập luật \texttt{IngredientRule} dựa trên reportedEffect và skinType.
    \item Backend tạo một bản ghi \texttt{Notification} gửi đến userId của người báo cáo với tiêu đề và nội dung tương ứng.
    \item Frontend xoá dòng báo cáo khỏi danh sách, hiển thị thông báo thành công.
\end{enumerate} \\ \hline
\textbf{Luồng sự kiện phụ / Ngoại lệ} & 
\begin{itemize}
    \item \textit{Báo cáo đã bị xử lý bởi Admin khác:} Nếu report không còn PENDING, API trả lỗi → frontend thông báo và tải lại danh sách.
    \item \textit{Admin huỷ modal:} Không có thay đổi.
\end{itemize} \\ \hline
\end{tabular}
\end{center}







\vspace{0.5cm}
\textbf{Đặc tả Use Case: Gửi thông báo đến người dùng (Admin)}

\begin{center}
\renewcommand{\arraystretch}{1.4}
\begin{tabular}{|p{4cm}|p{10cm}|}
\hline
\rowcolor{blue!15}
\textbf{Tên Use Case} & \textbf{Gửi thông báo đến người dùng} \\ \hline
\textbf{Tác nhân} & Hệ thống (kích hoạt bởi Admin) \\ \hline
\textbf{Mô tả} & Khi Admin duyệt (APPROVED) hoặc từ chối (REJECTED) một báo cáo phân loại, hệ thống sẽ tự động tạo thông báo gửi đến người dùng đã gửi báo cáo đó. \\ \hline
\textbf{Tiền điều kiện} & Admin đã thực hiện thao tác duyệt/từ chối báo cáo (Use Case “Duyệt báo cáo cộng đồng”) và quá trình cập nhật trạng thái báo cáo thành công. \\ \hline
\textbf{Hậu điều kiện} & Một bản ghi \texttt{Notification} được lưu vào cơ sở dữ liệu, người dùng tương ứng nhìn thấy thông báo trong chuông thông báo. \\ \hline
\textbf{Luồng sự kiện chính} & 
\begin{enumerate}
    \item Admin nhấn “Xác nhận” duyệt hoặc từ chối báo cáo (trong modal duyệt báo cáo).
    \item Hệ thống gọi API \texttt{POST /api/v1/reports/resolve} để cập nhật trạng thái báo cáo và (nếu được duyệt) cập nhật luật \texttt{IngredientRule}.
    \item Ngay sau khi cập nhật báo cáo thành công, \textbf{hệ thống tự động} tạo một thông báo mới trong bảng \texttt{Notification} với các thông tin:
        \begin{itemize}
            \item \texttt{userId}: ID của người đã gửi báo cáo.
            \item \texttt{type}: \texttt{REPORT\_RESOLVED}
            \item \texttt{title}: “Báo cáo của bạn đã được duyệt” hoặc “Báo cáo của bạn đã bị từ chối”
            \item \texttt{message}: Nội dung mặc định hoặc ghi chú của Admin (nếu có).
            \item \texttt{link}: \texttt{/community/reports} (trang báo cáo cộng đồng)
        \end{itemize}
    \item Backend lưu thông báo và trả về phản hồi thành công cho Admin (như một phần của response duyệt báo cáo).
    \item Khi người dùng đăng nhập hoặc mở chuông thông báo, thông báo mới sẽ hiển thị kèm trạng thái chưa đọc.
\end{enumerate} \\ \hline
\textbf{Luồng sự kiện phụ / Ngoại lệ} & 
\begin{itemize}
    \item \textit{Lỗi tạo thông báo}: Nếu vì lý do kỹ thuật không thể tạo thông báo, hệ thống vẫn hoàn tất việc duyệt báo cáo và ghi log lỗi. Người dùng có thể không nhận được thông báo nhưng trạng thái báo cáo vẫn được cập nhật.
    \item \textit{Admin không nhập ghi chú}: Hệ thống sử dụng nội dung mặc định như mô tả ở trên.
\end{itemize} \\ \hline
\end{tabular}
\end{center}












\subsection{Mô hình Hoạt động (Activity Diagram)}

\subsubsection{Sơ đồ hoạt động Đăng nhập và Cập nhật hồ sơ}
Sơ đồ này mô tả luồng kiểm tra rẽ nhánh cơ bản. Nếu người dùng nhập sai thông tin đăng nhập, hệ thống sẽ yêu cầu nhập lại. Khi đăng nhập thành công, người dùng có quyền truy cập vào chức năng thay đổi loại da để cập nhật thông tin trong cơ sở dữ liệu.

\begin{figure}[H]
    \centering
    \includegraphics[width=0.35\linewidth]{loginandupdate.png}
    \caption{Sơ đồ hoạt động Đăng nhập và Cập nhật hồ sơ}
    \label{fig:placeholder}
\end{figure}



\subsubsection{Sơ đồ hoạt động Phân tích thành phần mỹ phẩm (Core Flow)}
Đây là luồng hoạt động phức tạp nhất, bao gồm các vòng lặp (Loop) và nút quyết định (Decision node). 
Hệ thống sẽ lặp qua từng thành phần trong danh sách người dùng nhập. Với mỗi thành phần, nếu tồn tại trong CSDL, hệ thống tiếp tục rẽ nhánh kiểm tra tập luật theo loại da; nếu không tồn tại, tự động gán nhãn "Trung tính". Cuối cùng, kết thúc vòng lặp và hiển thị giao diện kết quả tổng hợp.
\begin{figure}[H]
    \centering
    \includegraphics[width=1\linewidth]{coreflow.png}
    \caption{Sơ đồ hoạt động Phân tích thành phần (Core Flow)}
    \label{fig:placeholder}
\end{figure}


\subsubsection{Sơ đồ hoạt động Quản lý dữ liệu (Dành cho Admin)}
Sơ đồ thể hiện luồng thao tác của quản trị viên đối với cơ sở dữ liệu (CRUD). Tùy thuộc vào hành động Admin chọn (Thêm, Sửa hoặc Xóa), hệ thống sẽ thực thi luồng tương ứng, lưu vào CSDL và tải lại trang quản trị.

\begin{figure}[H]
    \centering
    \includegraphics[width=0.95\linewidth]{datacontrol1.png}
    \caption{Sơ đồ hoạt động Quản trị dữ liệu hệ thống}
    \label{fig:placeholder}
\end{figure}


\begin{figure}[H]
    \centering
    \includegraphics[width=0.65\linewidth]{datacontrol2.png}
    \caption{Sơ đồ hoạt động Quản trị dữ liệu hệ thống}
    \label{fig:placeholder}
\end{figure}


\subsection{Mô hình Tuần tự (Sequence Diagram)}


\subsubsection{Sơ đồ tuần tự Đăng ký tài khoản}
Quy trình đăng ký đảm bảo tính an toàn dữ liệu ngay từ bước khởi tạo. Mật khẩu của người dùng không bao giờ được lưu dưới dạng văn bản thuần túy mà luôn được xử lý qua thuật toán băm bảo mật.

\begin{figure}[H]
    \centering
    \includegraphics[width=0.95\linewidth]{sequence_register.png}
    \caption{Sơ đồ tuần tự quy trình Đăng ký tài khoản}
\end{figure}



\subsubsection{Sơ đồ tuần tự Đăng nhập}
Sơ đồ tuần tự thể hiện sự tương tác theo thứ tự thời gian giữa các đối tượng: Người dùng (Actor), Giao diện Web (View), Bộ điều khiển (Controller) và Cơ sở dữ liệu (Database). Quá trình bắt đầu từ khi User nhấn nút Submit cho đến khi Server trả về thông báo trạng thái.

\begin{figure}[H]
    \centering
    \includegraphics[width=1\linewidth]{sequence_login.png}
    \caption{Sơ đồ tuần tự quy trình Đăng nhập}
    \label{fig:placeholder}
\end{figure}




\subsubsection{Sơ đồ tuần tự Phân tích thành phần}
Mô tả chi tiết cách giao diện Frontend gửi chuỗi INCI qua API phân tích xuống Backend. Controller lấy loại da mới nhất của người dùng, dịch vụ đối chiếu Ingredients/Rules và trả về kết quả phân loại; sau khi nhận kết quả, Frontend gọi riêng API gợi ý sản phẩm để hiển thị các thẻ sản phẩm an toàn.
\begin{figure}[H]
    \centering
    \includegraphics[width=0.95\linewidth]{sequence_core_1.png}
    \caption{Sơ đồ tuần tự luồng Phân tích thành phần và Gợi ý}
    \label{fig:placeholder}
\end{figure}



\subsubsection{Sơ đồ tuần tự Đề xuất sản phẩm an toàn}
Sơ đồ mô tả quy trình hệ thống tự động lọc và đề xuất sản phẩm dựa trên loại da của người dùng. Điểm nhấn của luồng này nằm ở tầng Business Logic (Service), nơi hệ thống áp dụng bộ lọc an toàn tuyệt đối (Safety-first filter) để loại bỏ mọi sản phẩm chứa chất gây kích ứng, sau đó tính điểm và xáo trộn ngẫu nhiên (shuffle) để trả về tối đa 3 sản phẩm từ nhóm điểm cao.

\begin{figure}[H]
    \centering
    \includegraphics[width=0.95\linewidth]{sequence_product_recommendation.png} % Thay tên file ảnh xuất từ Mermaid vào đây
    \caption{Sơ đồ tuần tự quy trình Lọc và Đề xuất sản phẩm an toàn}
\end{figure}






\subsubsection{Sơ đồ tuần tự Quản trị dữ liệu (Admin CRUD)}
Sơ đồ tuần tự này mô tả chi tiết luồng tương tác theo thời gian khi Quản trị viên (Admin) thực hiện các thao tác quản lý dữ liệu cốt lõi như Thêm mới, Chỉnh sửa, hoặc Xóa (CRUD) thông tin thành phần hóa học và tập luật. Quá trình bắt đầu từ lúc Admin gửi yêu cầu thao tác từ giao diện Web (View), hệ thống truyền dữ liệu qua Bộ điều khiển (Controller) để kiểm tra tính hợp lệ và xác thực quyền truy cập, sau đó thực thi các câu lệnh truy vấn tương ứng xuống Cơ sở dữ liệu (Database). Cuối cùng, kết quả thành công hay thất bại sẽ được phản hồi ngược lại và cập nhật trên giao diện của Admin.
\begin{figure}[H]
    \centering
    \includegraphics[width=0.95\linewidth]{sequence_admin_1.png}
    \caption{Sơ đồ tuần tự luồng Quản trị dữ liệu hệ thống (Thêm/Sửa/Xóa)}
    \label{fig:placeholder}
\end{figure}



\subsubsection{Sơ đồ tuần tự Quản lý người dùng và Phân quyền}
Sơ đồ minh họa cơ chế bảo mật đa tầng của hệ thống khi thực hiện các tác vụ nhạy cảm như Khóa/Mở khóa hoặc Xóa tài khoản người dùng. Request bắt buộc phải vượt qua hai "lớp khiên" bảo vệ liên tiếp: \textit{authMiddleware} (kiểm tra tính hợp lệ của phiên đăng nhập) và \textit{adminMiddleware} (kiểm tra quyền hạn Admin). Nếu vượt qua, tầng Service mới thực thi việc đảo ngược trạng thái \textit{isActive} trong CSDL.

\begin{figure}[H]
    \centering
    \includegraphics[width=0.95\linewidth]{sequence_admin_user_management.png} % Thay tên file ảnh xuất từ Mermaid vào đây
    \caption{Sơ đồ tuần tự quy trình Quản lý tài khoản người dùng và Phân quyền}
\end{figure}




\subsubsection{Sơ đồ tuần tự Nhập dữ liệu từ Excel}
Tính năng này cho phép Quản trị viên cập nhật hàng loạt danh mục thành phần hoặc sản phẩm mà không cần thao tác thủ công từng bản ghi, giúp tối ưu hóa hiệu suất quản trị.

\begin{figure}[H]
    \centering
    \includegraphics[width=0.95\linewidth]{sequence_excel_import.png}
    \caption{Sơ đồ tuần tự quy trình Import dữ liệu bằng Excel}
\end{figure}

\textbf{Cơ chế xử lý dữ liệu lớn:}
\begin{itemize}
    \item \textbf{Xử lý tại bộ nhớ đệm:} File Excel được tải lên dưới dạng \textit{Buffer} và bóc tách ngay tại bộ nhớ (\textit{Memory Storage}) để tăng tốc độ xử lý.
    \item \textbf{Chiến lược Upsert:} Thay vì chỉ thêm mới, hệ thống thông minh tự động nhận diện nếu tên thành phần đã tồn tại thì sẽ cập nhật mô tả, nếu chưa có sẽ tạo mới, giúp dữ liệu luôn đồng bộ.
\end{itemize}







\subsubsection{Sơ đồ tuần tự chức năng OCR (Trích xuất INCI từ ảnh)}
Sơ đồ dưới đây mô tả luồng xử lý khi người dùng tải ảnh chụp bảng thành phần. Hệ thống gửi ảnh đến OCR.space API, nhận kết quả nhận dạng, sau đó dùng bộ lọc \texttt{ingredientsExtractor} để chuẩn hóa chuỗi INCI và tự động điền vào ô nhập liệu.

\begin{figure}[H]
    \centering
    \includegraphics[width=0.95\linewidth]{sequence_ocr.png}
    \caption{Sơ đồ tuần tự chức năng OCR}
    \label{fig:seq_ocr}
\end{figure}

\subsubsection{Sơ đồ tuần tự báo cáo phân loại sai thành phần}
Quy trình bắt đầu khi người dùng mở modal chi tiết thành phần, chọn phân loại đề xuất và nhập lý do. Hệ thống tìm \texttt{ingredientId} qua API tìm kiếm, sau đó tạo báo cáo với trạng thái \texttt{PENDING}.

\begin{figure}[H]
    \centering
    \includegraphics[width=0.95\linewidth]{sequence_report.png}
    \caption{Sơ đồ tuần tự báo cáo phân loại sai thành phần}
    \label{fig:seq_report}
\end{figure}


\subsubsection{Sơ đồ tuần tự bình chọn báo cáo (Vote)}
Minh họa cơ chế \textit{optimistic update}: giao diện được cập nhật ngay khi người dùng bấm nút, sau đó gửi yêu cầu lên server. Nếu thất bại, hệ thống tự động rollback để đảm bảo nhất quán dữ liệu.

\begin{figure}[H]
    \centering
    \includegraphics[width=0.9\linewidth]{sequence_vote.png}
    \caption{Sơ đồ tuần tự bình chọn báo cáo}
    \label{fig:seq_vote}
\end{figure}



\subsubsection{Sơ đồ tuần tự duyệt báo cáo (Admin)}
Quản trị viên xem danh sách báo cáo đang chờ, chọn chấp nhận hoặc từ chối. Nếu duyệt, hệ thống tự động cập nhật \texttt{IngredientRule} (thông qua \texttt{adminService}) và gửi thông báo đến người dùng đã báo cáo.

\begin{figure}[H]
    \centering
    \includegraphics[width=0.95\linewidth]{sequence_resolve_report.png}
    \caption{Sơ đồ tuần tự duyệt báo cáo (Admin)}
    \label{fig:seq_resolve}
\end{figure}



\subsubsection{Sơ đồ tuần tự xem thông báo và đánh dấu đã đọc}
Hệ thống tải thông báo khi người dùng đăng nhập và khi mở dropdown; khi dropdown đang đóng, hệ thống polling mỗi 30 giây để lấy thông báo mới. Các thao tác đánh dấu đã đọc (từng cái hoặc tất cả) được gửi về server và cập nhật giao diện tức thì.

\begin{figure}[H]
    \centering
    \includegraphics[width=0.95\linewidth]{sequence_notification.png}
    \caption{Sơ đồ tuần tự xem thông báo và đánh dấu đã đọc}
    \label{fig:seq_notification}
\end{figure}







\subsection{Mô hình Trạng thái (State Machine Diagram)}

Sơ đồ trạng thái dùng để mô tả sự thay đổi vòng đời của một đối tượng cụ thể trong hệ thống khi có các biến cố (event) tác động. 

\subsubsection{Sơ đồ trạng thái vòng đời tài khoản người dùng}

Sơ đồ dưới đây minh họa sự chuyển đổi trạng thái của thực thể User dựa trên biến cờ \texttt{isActive} trong Cơ sở dữ liệu. Nhờ cơ chế này, hệ thống bảo đảm an ninh bằng cách chặn đứng các tài khoản vi phạm (Locked) ngay tại cửa ngõ đăng nhập, đồng thời dọn dẹp sạch sẽ dữ liệu rác (Cascade Delete) khi tài khoản bị xóa bỏ.

\begin{figure}[H]
    \centering
    \includegraphics[width=0.95\linewidth]{state_user_lifecycle.jpg} % Đổi tên file ảnh bạn tải về từ PlantUML vào đây
    \caption{Sơ đồ trạng thái: Vòng đời hoạt động của Tài khoản Người dùng}
\end{figure}



\subsubsection{Sơ đồ trạng thái của IngredientReport}
Sơ đồ mô tả vòng đời của một báo cáo phân loại sai thành phần, từ khi người dùng gửi (PENDING) đến khi Quản trị viên duyệt (APPROVED) hoặc từ chối (REJECTED). Khi được duyệt, hệ thống tự động cập nhật IngredientRule tương ứng.

\begin{figure}[H]
    \centering
    \includegraphics[width=0.95\linewidth]{state_ingredientreport.png}
    \caption{Sơ đồ trạng thái của IngredientReport}
    \label{fig:state_ingredientreport}
\end{figure}


\subsubsection{Sơ đồ trạng thái của Notification}
Thông báo hệ thống được tạo ở trạng thái UNREAD. Người dùng có thể đánh dấu đã đọc từng thông báo (markAsRead) hoặc tất cả (markAllAsRead), chuyển sang trạng thái READ.

\begin{figure}[H]
    \centering
    \includegraphics[width=0.35\linewidth]{state_notification.png}
    \caption{Sơ đồ trạng thái của Notification}
    \label{fig:state_notification}
\end{figure}


\subsubsection{Sơ đồ trạng thái của AnalysisHistory (Tuỳ chọn)}
Minh họa tính năng “Phân tích lại” từ lịch sử. Một bản ghi lịch sử ban đầu ở trạng thái NEW. Khi người dùng bấm nút “Phân tích lại”, hệ thống tạo một bản phân tích mới nhưng vẫn giữ bản ghi cũ, đánh dấu trạng thái là RE\_ANALYZED (có thể dùng để thống kê hành vi người dùng).

\begin{figure}[H]
    \centering
    \includegraphics[width=0.9\linewidth]{state_analysishistory.png}
    \caption{Sơ đồ trạng thái của AnalysisHistory (Tuỳ chọn)}
    \label{fig:state_analysishistory}
\end{figure}




\section{Thiết kế hệ thống}
\subsection{Kiến trúc hệ thống}

\subsubsection{Mô hình kiến trúc}

Hệ thống SKINMATE được xây dựng theo mô hình \textbf{Decoupled Monorepo} dựa trên kiến trúc Client--Server hiện đại. Điểm nổi bật là sự tách biệt hoàn toàn giữa giao diện người dùng (Frontend) và xử lý nghiệp vụ (Backend), giao tiếp với nhau thông qua chuẩn RESTful API.

\begin{itemize}
    \item \textbf{Client (Frontend - Next.js 16):} \\
    Được phát triển bằng framework Next.js 16 kết hợp với React 19 và TypeScript. Thay vì sử dụng CSS truyền thống, hệ thống ứng dụng \textbf{Tailwind CSS v4} để xây dựng giao diện Luxury theo phong cách "Glassmorphism". Client thực hiện các tác vụ như quản lý trạng thái đăng nhập (AuthContext), gửi chuỗi INCI qua API, tải ảnh cho OCR, hiển thị báo cáo thống kê trực quan bằng thư viện Recharts, giao diện báo cáo cộng đồng và chuông thông báo.

    \item \textbf{Server (Backend - Node.js/Express):} \\
    Được xây dựng bằng Node.js kết hợp Express.js trên nền tảng TypeScript để đảm bảo tính an toàn về kiểu dữ liệu (Type-safety). Backend đóng vai trò là "bộ não" điều phối toàn bộ luồng nghiệp vụ: từ xác thực JWT, xử lý logic tách chuỗi INCI, gọi \textbf{Gemini AI API} cho các thành phần lạ, xử lý nhận dạng ảnh \textbf{OCR}, quản lý báo cáo phân loại và bình chọn từ cộng đồng, tự động tạo thông báo khi duyệt báo cáo, đến việc quản lý xuất/nhập file Excel.

    \item \textbf{Cơ sở dữ liệu (PostgreSQL \& Prisma ORM):} \\
    Hệ thống sử dụng \textbf{PostgreSQL 15} làm hệ quản trị CSDL chính. Điểm khác biệt lớn nhất là việc ứng dụng \textbf{Prisma ORM} để thay thế các câu lệnh SQL thuần. Prisma đóng vai trò trung gian giúp ánh xạ các Model trong mã nguồn trực tiếp xuống các bảng dữ liệu, bao gồm các bảng mở rộng như IngredientReport, ReportVote, Notification, đảm bảo tốc độ truy vấn cao và tính toàn vẹn dữ liệu thông qua cơ chế Schema Migration.
\end{itemize}

\paragraph{Lý do lựa chọn mô hình:}Việc sử dụng TypeScript đồng bộ cho cả Frontend và Backend giúp kiểm soát kiểu dữ liệu khi trao đổi dữ liệu. Mô hình này cho phép triển khai ứng dụng trên các nền tảng Container (Docker), dễ dàng tích hợp thêm các dịch vụ AI bên ngoài cũng như các module mới (OCR, báo cáo, thông báo) mà không làm thay đổi cấu trúc cốt lõi của phần mềm.

\paragraph{Đánh giá:}Hệ thống tuân thủ nguyên tắc \textbf{SoC (Separation of Concerns)}. Backend tách các lớp Routes, Middlewares, Controllers và Services; riêng OCR được tổ chức theo module, Report có service nghiệp vụ, còn Notification hiện được xử lý trực tiếp trong Controller thông qua Prisma.

\subsection{Mô tả các thành phần trong hệ thống}

\renewcommand{\arraystretch}{1.3}
\begin{table}[H]
\centering
\rowcolors{2}{gray!10}{white}
\begin{tabularx}{\textwidth}{|c|>{\centering\arraybackslash}p{3.5cm}|X|}
\hline
\rowcolor{gray!30}
\textbf{STT} & \textbf{Thành phần} & \textbf{Mô tả} \\
\hline

1 & Frontend (Next.js 16) & Giao diện Luxury sử dụng Tailwind v4, xử lý logic phía client bằng React Hooks và Context API. \\
\hline

2 & RESTful API & Cổng giao tiếp chuẩn hóa giữa Frontend và Backend, truyền tải dữ liệu dưới dạng JSON. \\
\hline

3 & Auth/Admin Middleware & \texttt{authMiddleware} kiểm tra JWT và \texttt{adminMiddleware} kiểm tra vai trò \texttt{ADMIN} cho các endpoint quản trị. \\
\hline

4 & Analysis Service (AI Integrated) & Dịch vụ lõi xử lý INCI. Tích hợp Google Gemini Flash (\texttt{gemini-flash-latest}) để phân tích và tự động cập nhật dữ liệu thành phần mới. \\
\hline

5 & Excel Service & Xử lý logic Import/Export dữ liệu quy mô lớn thông qua thư viện \textit{xlsx}. \\
\hline

6 & OCR Service & Xử lý nhận dạng ký tự từ ảnh chụp bảng thành phần, tích hợp API OCR.space và bộ trích xuất INCI. \\
\hline

7 & Report Service & Quản lý báo cáo phân loại sai thành phần từ cộng đồng, xử lý bình chọn (vote) và duyệt báo cáo (Admin). \\
\hline

8 & Notification Controller & Cung cấp API lấy/đánh dấu đã đọc/gửi thông báo; thông báo duyệt báo cáo được tạo từ \texttt{reportService} và hiển thị qua chuông thông báo. \\
\hline

9 & Internal Layered Architecture & Cấu trúc mã nguồn Backend chia làm 4 lớp: \textbf{Routes} (Định tuyến) $\rightarrow$ \textbf{Middlewares} (Bảo mật) $\rightarrow$ \textbf{Controllers} (Điều phối) $\rightarrow$ \textbf{Services} (Logic nghiệp vụ). \\
\hline

10 & Prisma Client & Tầng truy xuất dữ liệu (Data Access Layer), thực hiện các truy vấn PostgreSQL một cách Type-safe, bao gồm các model mới: IngredientReport, ReportVote, Notification. \\
\hline

11 & PostgreSQL Database & Hệ quản trị CSDL quan hệ lưu trữ: Users, Ingredients, Ingredient\_Rules, Products, AnalysisHistory, IngredientReport, ReportVote, Notification. \\
\hline

\end{tabularx}
\caption{Mô tả các thành phần kiến trúc hệ thống}
\end{table}



\section{Thiết kế dữ liệu}
\subsection{Sơ đồ lớp}




\begin{figure}[H]
    \centering
    \includegraphics[width=0.9\linewidth]{class_diag_09_05.png}
    \caption{Sơ đồ Lớp (Class Diagram) tổng quát của hệ thống}
    \label{fig:placeholder}
\end{figure}




\subsection{Danh sách các lớp đối tượng}

Hệ thống được thiết kế dựa trên các lớp đối tượng cốt lõi sau đây nhằm đảm bảo tính toàn vẹn của luồng phân tích Hybrid và quản trị dữ liệu tập trung:

\begin{center}
\renewcommand{\arraystretch}{1.5}
\rowcolors{2}{gray!10}{white}
\begin{tabularx}{\textwidth}{|c|l|X|}
\hline
\rowcolor{blue!20}
\textbf{STT} & \textbf{Tên lớp} & \textbf{Mô tả} \\ \hline

1 & \textbf{User} & Lưu trữ thông tin định danh người dùng, bao gồm: tên đăng nhập, mật khẩu đã băm (Bcrypt), vai trò (User/Admin), loại da (\textit{SkinType}) và trạng thái tài khoản (\textit{isActive}). \\ \hline

2 & \textbf{Ingredient} & Đại diện cho bộ từ điển các hoạt chất hóa mỹ phẩm. Mỗi đối tượng chứa tên hoạt chất chuẩn hóa (viết thường) và mô tả chi tiết về tính chất khoa học của chất đó. \\ \hline

3 & \textbf{IngredientRule} & Lớp định nghĩa logic cốt lõi cho hệ thống \textit{Rule-based}. Nó thiết lập mối quan hệ giữa một \texttt{Ingredient} và một \texttt{SkinType} để xác định mức độ an toàn (GOOD, BAD, NEUTRAL). \\ \hline

4 & \textbf{Product} & Lưu trữ thông tin về các sản phẩm mỹ phẩm dùng cho mục đích gợi ý, bao gồm tên sản phẩm, thương hiệu (\textit{brand}) và đường dẫn hình ảnh minh họa. \\ \hline

5 & \textbf{ProductIngredient} & Lớp trung gian giải quyết quan hệ nhiều-nhiều giữa \texttt{Product} và \texttt{Ingredient}. Lớp này lưu trữ thêm thông tin về vị trí (\textit{position}) của hoạt chất trong bảng thành phần sản phẩm. \\ \hline

6 & \textbf{AnalysisHistory} & Ghi lại nhật ký mỗi lần người dùng thực hiện phân tích. Lưu trữ chuỗi INCI gốc (\textit{rawInput}), mã người dùng và dấu thời gian (\textit{createdAt}) để hỗ trợ tính năng xem lại lịch sử. \\ \hline



7 & \textbf{IngredientReport} & Lưu trữ các báo cáo của người dùng về phân loại sai của một thành phần. Bao gồm: ingredientId, userId, skinType, reportedEffect, lý do, trạng thái (PENDING/APPROVED/REJECTED), thông tin người duyệt. \\ \hline
8 & \textbf{ReportVote} & Ghi nhận lượt bình chọn (UP/DOWN) của người dùng cho một báo cáo. Mỗi user chỉ được một vote cho một report. \\ \hline
9 & \textbf{Notification} & Lưu thông báo hệ thống gửi đến từng người dùng: loại thông báo, tiêu đề, nội dung, liên kết, trạng thái đã đọc. \\ \hline
\end{tabularx}
\end{center}

% \textbf{Ghi chú kỹ thuật:} 
% \begin{itemize}
%     \item Toàn bộ các lớp đối tượng trên được ánh xạ (\textit{Mapping}) trực tiếp từ Prisma Models vào cơ sở dữ liệu PostgreSQL.
%     \item Các lớp đều hỗ trợ cơ chế xóa dây chuyền (\textit{onDelete: Cascade}) để đảm bảo không có dữ liệu mồ côi khi một đối tượng cha bị xóa khỏi hệ thống.
% \end{itemize}



\subsection{Danh sách các quan hệ}

Các lớp đối tượng trong hệ thống SKINMATE được liên kết chặt chẽ thông qua các ràng buộc khóa ngoại, đảm bảo tính toàn vẹn dữ liệu và hỗ trợ các truy vấn phức tạp từ Prisma ORM.

\begin{center}
\renewcommand{\arraystretch}{1.5}
\rowcolors{2}{gray!10}{white}
\begin{tabularx}{\textwidth}{|c|l|X|}
\hline
\rowcolor{blue!20}
\textbf{STT} & \textbf{Tên quan hệ} & \textbf{Mô tả} \\ \hline

1 & \textbf{User -- AnalysisHistory} & Quan hệ Một -- Nhiều (1:N). Một người dùng có thể thực hiện nhiều lượt phân tích và lưu lại lịch sử. Khi xóa một \texttt{User}, toàn bộ \texttt{AnalysisHistory} liên quan sẽ bị xóa theo (\textit{Cascade Delete}). \\ \hline

2 & \textbf{Ingredient -- IngredientRule} & Quan hệ Một -- Nhiều (1:N). Một hoạt chất có thể có nhiều quy tắc an toàn khác nhau tương ứng với từng loại da. Mỗi quy tắc (\texttt{IngredientRule}) bắt buộc phải thuộc về một \texttt{Ingredient} duy nhất. \\ \hline

3 & \textbf{Product -- ProductIngredient} & Quan hệ Một -- Nhiều (1:N). Một sản phẩm chứa nhiều bản ghi liên kết thành phần. Đây là nhánh trái của quan hệ Nhiều -- Nhiều giữa \texttt{Product} và \texttt{Ingredient}. \\ \hline

4 & \textbf{Ingredient -- ProductIngredient} & Quan hệ Một -- Nhiều (1:N). Một hoạt chất có thể xuất hiện trong bảng thành phần của nhiều sản phẩm khác nhau. Đây là nhánh phải của quan hệ Nhiều -- Nhiều giữa \texttt{Product} và \texttt{Ingredient}. \\ \hline

5 & \textbf{IngredientRule -- SkinType} & Quan hệ logic dựa trên \textit{Enum}. Mỗi \texttt{IngredientRule} được ánh xạ với một giá trị cụ thể trong danh sách \texttt{SkinType} để xác định ngữ cảnh phân tích. \\ \hline

6 & \textbf{User -- IngredientReport} & Quan hệ Một -- Nhiều (1:N). Một người dùng có thể gửi nhiều báo cáo phân loại sai thành phần. Khi xóa \texttt{User}, các báo cáo liên quan cũng bị xóa (\textit{Cascade Delete}). \\ \hline

7 & \textbf{Ingredient -- IngredientReport} & Quan hệ Một -- Nhiều (1:N). Một hoạt chất có thể bị nhiều người dùng báo cáo phân loại sai. Mỗi báo cáo luôn gắn với một \texttt{Ingredient} duy nhất. \\ \hline

8 & \textbf{IngredientReport -- ReportVote} & Quan hệ Một -- Nhiều (1:N). Một báo cáo có thể nhận được nhiều phiếu bình chọn (UP/DOWN) từ cộng đồng. Khi xóa báo cáo, các phiếu tương ứng cũng bị xóa (\textit{Cascade Delete}). \\ \hline

9 & \textbf{User -- ReportVote} & Quan hệ Một -- Nhiều (1:N). Một người dùng có thể bình chọn cho nhiều báo cáo khác nhau. Mỗi phiếu bình chọn thuộc về một người dùng xác định. \\ \hline

10 & \textbf{User -- Notification} & Quan hệ Một -- Nhiều (1:N). Mỗi người dùng có thể nhận được nhiều thông báo hệ thống (duyệt báo cáo, tin nhắn từ Admin). Khi xóa \texttt{User}, thông báo liên quan cũng bị xóa (\textit{Cascade Delete}). \\ \hline

\end{tabularx}
\end{center}



\subsection{Mô tả từng lớp đối tượng}

Phần này mô tả chi tiết các thuộc tính, kiểu dữ liệu và các ràng buộc kỹ thuật của từng lớp đối tượng được triển khai trong cơ sở dữ liệu PostgreSQL của hệ thống SKINMATE.

% --- LỚP USER ---
\subsubsection{Lớp đối tượng User}
Lưu trữ thông tin tài khoản, danh tính hiển thị và cấu hình loại da của người dùng.

\begin{center}
\renewcommand{\arraystretch}{1.4}
\begin{tabularx}{\textwidth}{|c|l|l|l|X|}
\hline
\rowcolor{blue!15}
\textbf{STT} & \textbf{Thuộc tính} & \textbf{Kiểu dữ liệu} & \textbf{Ràng buộc} & \textbf{Mô tả} \\ \hline
1 & id & String (UUID) & PK, Default & Định danh duy nhất (UUID) cho mỗi người dùng. \\ \hline
2 & username & String & Unique, Not Null & Tên đăng nhập dùng để xác thực hệ thống. \\ \hline
3 & displayName & String & Nullable & Họ và tên hiển thị trên giao diện (Tùy chọn). \\ \hline
4 & passwordHash & String & Not Null & Mật khẩu đã được băm bằng thuật toán Bcrypt. \\ \hline
5 & skinType & Enum & Nullable, Default: NORMAL & Các loại da: OILY, DRY, SENSITIVE, COMBINATION, NORMAL. \\ \hline
6 & role & Enum & Default: USER & Vai trò hệ thống: USER hoặc ADMIN. \\ \hline
7 & isActive & Boolean & Default: true & Trạng thái tài khoản (Nếu false là bị khóa). \\ \hline
8 & createdAt & DateTime & Default: now() & Thời điểm khởi tạo tài khoản. \\ \hline
9 & updatedAt & DateTime & Auto Update & Thời điểm cập nhật thông tin gần nhất. \\ \hline
\end{tabularx}
\end{center}

% --- LỚP INGREDIENT ---
\subsubsection{Lớp đối tượng Ingredient}
Danh mục các hoạt chất hóa mỹ phẩm trong hệ thống.

\begin{center}
\renewcommand{\arraystretch}{1.4}
\begin{tabularx}{\textwidth}{|c|l|l|l|X|}
\hline
\rowcolor{blue!15}
\textbf{STT} & \textbf{Thuộc tính} & \textbf{Kiểu dữ liệu} & \textbf{Ràng buộc} & \textbf{Mô tả} \\ \hline
1 & id & Int & PK, Auto Inc & Định danh duy nhất cho hoạt chất. \\ \hline
2 & name & String & Unique, Not Null & Tên hoạt chất (luôn lưu ở dạng chữ thường). \\ \hline
3 & description & Text & Nullable & Mô tả khoa học hoặc công dụng của chất. \\ \hline
\end{tabularx}
\end{center}

% --- LỚP INGREDIENTRULE ---
\subsubsection{Lớp đối tượng IngredientRule}
Định nghĩa quy tắc an toàn của hoạt chất đối với từng loại da.

\begin{center}
\renewcommand{\arraystretch}{1.4}
\begin{tabularx}{\textwidth}{|c|l|l|l|X|}
\hline
\rowcolor{blue!15}
\textbf{STT} & \textbf{Thuộc tính} & \textbf{Kiểu dữ liệu} & \textbf{Ràng buộc} & \textbf{Mô tả} \\ \hline
1 & id & Int & PK, Auto Inc & Định danh duy nhất cho quy tắc. \\ \hline
2 & ingredientId & Int & FK (Ingredient) & Liên kết với hoạt chất tương ứng. \\ \hline
3 & skinType & Enum & Not Null & Loại da áp dụng quy tắc này. \\ \hline
4 & effect & Enum & Not Null & Đánh giá an toàn: GOOD, BAD, NEUTRAL. \\ \hline
- & \textit{(Composite)} & \textit{Unique} & \textit{ingredientId + skinType} & Đảm bảo không trùng lặp cặp (Chất, Loại da). \\ \hline
\end{tabularx}
\end{center}

% --- LỚP PRODUCT ---
\subsubsection{Lớp đối tượng Product}
Thông tin về các sản phẩm mỹ phẩm dùng để gợi ý.

\begin{center}
\renewcommand{\arraystretch}{1.4}
\begin{tabularx}{\textwidth}{|c|l|l|l|X|}
\hline
\rowcolor{blue!15}
\textbf{STT} & \textbf{Thuộc tính} & \textbf{Kiểu dữ liệu} & \textbf{Ràng buộc} & \textbf{Mô tả} \\ \hline
1 & id & String (UUID) & PK, Default & Định danh duy nhất cho sản phẩm. \\ \hline
2 & name & String & Not Null & Tên thương mại của sản phẩm. \\ \hline
3 & brand & String & Not Null & Thương hiệu sản xuất. \\ \hline
4 & imageUrl & String & Nullable & Đường dẫn đến hình ảnh sản phẩm. \\ \hline
5 & createdAt & DateTime & Default: now() & Thời điểm thêm sản phẩm vào hệ thống. \\ \hline
\end{tabularx}
\end{center}

% --- LỚP PRODUCTINGREDIENT ---
\subsubsection{Lớp đối tượng ProductIngredient}
Bảng trung gian liên kết sản phẩm với các thành phần của nó.

\begin{center}
\renewcommand{\arraystretch}{1.4}
\begin{tabularx}{\textwidth}{|c|l|l|l|X|}
\hline
\rowcolor{blue!15}
\textbf{STT} & \textbf{Thuộc tính} & \textbf{Kiểu dữ liệu} & \textbf{Ràng buộc} & \textbf{Mô tả} \\ \hline
1 & productId & String & PK, FK (Product) & Liên kết tới mã sản phẩm. \\ \hline
2 & ingredientId & Int & PK, FK (Ingredient) & Liên kết tới mã hoạt chất. \\ \hline
3 & position & Int & Not Null & Thứ tự của thành phần trong bảng INCI. \\ \hline
\end{tabularx}
\end{center}

% --- LỚP ANALYSISHISTORY ---
\subsubsection{Lớp đối tượng AnalysisHistory}
Nhật ký các lần thực hiện phân tích của người dùng.

\begin{center}
\renewcommand{\arraystretch}{1.4}
\begin{tabularx}{\textwidth}{|c|l|l|l|X|}
\hline
\rowcolor{blue!15}
\textbf{STT} & \textbf{Thuộc tính} & \textbf{Kiểu dữ liệu} & \textbf{Ràng buộc} & \textbf{Mô tả} \\ \hline
1 & id & String (UUID) & PK, Default & Định danh duy nhất cho bản ghi lịch sử. \\ \hline
2 & userId & String & FK (User) & Liên kết tới người dùng thực hiện. \\ \hline
3 & rawInput & Text & Not Null & Chuỗi INCI gốc mà người dùng đã nhập. \\ \hline
4 & createdAt & DateTime & Default: now() & Thời điểm thực hiện phân tích. \\ \hline
\end{tabularx}
\end{center}




% --- LỚP INGREDIENTREPORT ---
\subsubsection{Lớp đối tượng IngredientReport}
Lưu trữ các báo cáo của người dùng về việc phân loại sai mức độ an toàn của một thành phần đối với một loại da cụ thể.

\begin{center}
\renewcommand{\arraystretch}{1.4}
\begin{tabularx}{\textwidth}{|c|l|l|l|X|}
\hline
\rowcolor{blue!15}
\textbf{STT} & \textbf{Thuộc tính} & \textbf{Kiểu dữ liệu} & \textbf{Ràng buộc} & \textbf{Mô tả} \\ \hline
1 & id & Int & PK, Auto Inc & Định danh duy nhất cho mỗi báo cáo. \\ \hline
2 & ingredientId & Int & FK (Ingredient), Not Null & Liên kết tới thành phần bị báo cáo. \\ \hline
3 & userId & String & FK (User), Not Null & Người dùng gửi báo cáo. \\ \hline
4 & skinType & Enum (SkinType) & Not Null & Loại da của người báo cáo tại thời điểm báo cáo (ảnh hưởng đến luật sẽ cập nhật). \\ \hline
5 & reportedEffect & Enum (SafetyEffect) & Not Null & Mức độ an toàn mà người dùng đề xuất: GOOD, BAD, NEUTRAL. \\ \hline
6 & reason & Text & Not Null & Lý do chi tiết cho đề xuất; giao diện yêu cầu tối thiểu 20 ký tự trước khi gửi. \\ \hline
7 & evidenceUrl & String & Nullable & Đường dẫn đến tài liệu tham khảo (nghiên cứu, bài báo,...). \\ \hline
8 & status & Enum (ReportStatus) & Default: PENDING & Trạng thái xử lý: PENDING, APPROVED, REJECTED. \\ \hline
9 & createdAt & DateTime & Default: now() & Thời điểm gửi báo cáo. \\ \hline
10 & resolvedAt & DateTime & Nullable & Thời điểm Admin duyệt/từ chối. \\ \hline
11 & resolvedBy & String & Nullable & ID của Admin thực hiện duyệt; Prisma schema hiện lưu chuỗi ID, không khai báo quan hệ khóa ngoại riêng. \\ \hline
12 & adminNote & Text & Nullable & Ghi chú của Admin khi duyệt (có thể dùng làm nội dung thông báo). \\ \hline
\end{tabularx}
\end{center}

% --- LỚP REPORTVOTE ---
\subsubsection{Lớp đối tượng ReportVote}
Ghi nhận phiếu bình chọn (UP/DOWN) của người dùng cho các báo cáo đang chờ duyệt, giúp cộng đồng ưu tiên các báo cáo đáng tin cậy.

\begin{center}
\renewcommand{\arraystretch}{1.4}
\begin{tabularx}{\textwidth}{|c|l|l|l|X|}
\hline
\rowcolor{blue!15}
\textbf{STT} & \textbf{Thuộc tính} & \textbf{Kiểu dữ liệu} & \textbf{Ràng buộc} & \textbf{Mô tả} \\ \hline
1 & id & Int & PK, Auto Inc & Định danh duy nhất cho phiếu bình chọn. \\ \hline
2 & reportId & Int & FK (IngredientReport), Not Null & Liên kết tới báo cáo được bình chọn. \\ \hline
3 & userId & String & FK (User), Not Null & Người dùng bình chọn. \\ \hline
4 & voteType & Enum (VoteType) & Not Null & Loại bình chọn: UP hoặc DOWN. \\ \hline
5 & createdAt & DateTime & Default: now() & Thời điểm bình chọn. \\ \hline
- & \textit{(Composite)} & \textit{Unique} & \textit{reportId + userId} & Mỗi người dùng chỉ được bình chọn một lần cho một báo cáo. \\ \hline
\end{tabularx}
\end{center}

% --- LỚP NOTIFICATION ---
\subsubsection{Lớp đối tượng Notification}
Lưu trữ các thông báo hệ thống gửi đến người dùng, bao gồm kết quả duyệt báo cáo, tin nhắn từ quản trị viên, v.v.

\begin{center}
\renewcommand{\arraystretch}{1.4}
\begin{tabularx}{\textwidth}{|c|l|l|l|X|}
\hline
\rowcolor{blue!15}
\textbf{STT} & \textbf{Thuộc tính} & \textbf{Kiểu dữ liệu} & \textbf{Ràng buộc} & \textbf{Mô tả} \\ \hline
1 & id & String (UUID) & PK, Default & Định danh duy nhất cho thông báo. \\ \hline
2 & userId & String & FK (User), Not Null & Người dùng nhận thông báo. \\ \hline
3 & type & String & Not Null & Loại thông báo, ví dụ: \texttt{REPORT\_RESOLVED}, \texttt{ADMIN\_MESSAGE}. \\ \hline
4 & title & String & Not Null & Tiêu đề ngắn gọn của thông báo. \\ \hline
5 & message & Text & Not Null & Nội dung chi tiết của thông báo. \\ \hline
6 & link & String & Nullable & Đường dẫn điều hướng khi người dùng nhấn vào thông báo (ví dụ: \texttt{/community/reports}). \\ \hline
7 & isRead & Boolean & Default: false & Trạng thái đã đọc của thông báo. \\ \hline
8 & createdAt & DateTime & Default: now() & Thời điểm tạo thông báo. \\ \hline
\end{tabularx}
\end{center}


\section{Thiết kế giao diện}

\subsection{Danh sách các màn hình}
Hệ thống được thiết kế theo phong cách Luxury Clean Beauty, tối ưu hóa trải nghiệm trên cả thiết bị di động và máy tính cá nhân. Dưới đây là danh sách các màn hình chính thực thi các nghiệp vụ của hệ thống:

\begin{center}
\renewcommand{\arraystretch}{1.5}
\rowcolors{2}{gray!10}{white}
\begin{tabularx}{\textwidth}{|c|l|l|X|}
\hline
\rowcolor{blue!20}
\textbf{STT} & \textbf{Màn hình} & \textbf{Loại màn hình} & \textbf{Chức năng} \\ \hline

1 & \textbf{Trang chủ} & Trang chính & Giới thiệu dịch vụ, kiểm tra trạng thái kết nối hệ thống và điều hướng người dùng. \\ \hline

2 & \textbf{Đăng ký} & Xác thực & Cho phép người dùng tạo tài khoản và thiết lập loại da ban đầu. \\ \hline

3 & \textbf{Đăng nhập} & Xác thực & Kiểm tra thông tin đăng nhập, nhận JWT và truy cập các tính năng cá nhân hóa. \\ \hline

4 & \textbf{Phân tích (Analysis)} & Nghiệp vụ chính & Nhập chuỗi INCI hoặc trích xuất từ ảnh bằng OCR, hiển thị kết quả phân loại màu sắc và đề xuất sản phẩm an toàn. \\ \hline

5 & \textbf{Hồ sơ (Profile)} & Cá nhân hóa & Xem thông tin định danh, cập nhật tên hiển thị, loại da hiện tại và đổi mật khẩu. \\ \hline

6 & \textbf{Lịch sử (History)} & Lưu trữ & Xem lại danh sách các lần phân tích cũ, thực hiện tái phân tích hoặc xóa lịch sử. \\ \hline

7 & \textbf{Admin Dashboard} & Quản trị & Trang tổng quan dành cho quản trị viên dẫn đến các tiểu mục quản lý. \\ \hline

8 & \textbf{Quản lý Thành phần} & Quản trị & Thực hiện CRUD (Thêm, Sửa, Xóa) danh mục các hoạt chất hóa học. \\ \hline

9 & \textbf{Quản lý Quy tắc} & Quản trị & Thiết lập mức độ an toàn (Good/Bad) cho từng cặp Chất -- Loại da. \\ \hline

10 & \textbf{Quản lý Sản phẩm} & Quản trị & Quản lý danh mục sản phẩm phục vụ thuật toán gợi ý. \\ \hline

11 & \textbf{Quản lý Người dùng} & Quản trị & Theo dõi danh sách thành viên, thực hiện Khóa hoặc Xóa tài khoản. \\ \hline

12 & \textbf{Báo cáo thống kê} & Quản trị & Xem biểu đồ phân bố loại da và các chỉ số tăng trưởng hệ thống. \\ \hline

13 & \textbf{Import / Export} & Tiện ích & Xử lý dữ liệu hàng loạt thông qua tệp tin Excel (.xlsx). \\ \hline

14 & \textbf{Chi tiết thành phần} & Popup (Modal) & Hiển thị mô tả chi tiết và lời khuyên khoa học cho một hoạt chất cụ thể. \\ \hline

15 & \textbf{Xác nhận thao tác} & Popup (Modal) & Cảnh báo và yêu cầu xác nhận trước khi thực hiện các thao tác xóa dữ liệu. \\ \hline



\end{tabularx}
\end{center}




\begin{center}
\renewcommand{\arraystretch}{1.5}
\rowcolors{2}{gray!10}{white}
\begin{tabularx}{\textwidth}{|c|l|l|X|}
\hline
\rowcolor{blue!20}
\textbf{STT} & \textbf{Màn hình} & \textbf{Loại màn hình} & \textbf{Chức năng} \\ \hline
16 & \textbf{Cộng đồng - Báo cáo} & Nghiệp vụ & Xem danh sách các báo cáo phân loại sai đang chờ duyệt, bình chọn (UP/DOWN) để ủng hộ hoặc phản đối các đề xuất từ cộng đồng. \\ \hline

17 & \textbf{Admin - Duyệt báo cáo} & Quản trị & Xem danh sách báo cáo cần duyệt (đã sắp xếp theo số phiếu), phê duyệt hoặc từ chối kèm ghi chú. Tự động cập nhật luật và gửi thông báo. \\ \hline

18 & \textbf{Thông báo} & Popup (Dropdown) & Hiển thị danh sách thông báo hệ thống (kết quả duyệt báo cáo, tin nhắn từ Admin). Đánh dấu đã đọc hoặc điều hướng đến liên kết liên quan. \\ \hline

19 & \textbf{Báo cáo phân loại} & Popup (Modal) & Cho phép người dùng đề xuất hiệu ứng đúng của một thành phần, nhập lý do và dẫn chứng để gửi cộng đồng xét duyệt. \\ \hline

\end{tabularx}
\end{center}


\textbf{Nguyên tắc thiết kế hệ thống:}
\begin{itemize}
    \item \textbf{Tính nhất quán (Consistency):} Sử dụng bộ mã màu Rose-Sage xuyên suốt để định hướng cảm xúc và hành vi người dùng.
    \item \textbf{Phòng ngừa lỗi (Error Prevention):} Sử dụng các Modal xác nhận cho các hành động xóa dữ liệu quan trọng ở phía Admin.
    \item \textbf{Khả năng đáp ứng (Responsiveness):} Toàn bộ các màn hình được xây dựng bằng Tailwind CSS v4, đảm bảo hiển thị tối ưu trên nhiều kích thước màn hình.
\end{itemize}


\subsection{Mô tả các màn hình}
\subsubsection{Trang chủ}

\textbf{a) Giao diện}

\begin{figure}[H]
    \centering
    \includegraphics[width=0.9\linewidth]{ui_home.pdf} % Thay bằng ảnh chụp màn hình thực tế của bạn
    \caption{Giao diện Trang chủ hệ thống SKINMATE}
    \label{fig:ui_home}
\end{figure}

% \includepdf[pages=1]{ui_home.pdf}


\textbf{b) Mô tả các đối tượng trên màn hình}

\begin{center}
\renewcommand{\arraystretch}{1.5}
\rowcolors{2}{gray!10}{white}
\begin{tabularx}{\textwidth}{|c|l|l|l|X|}
\hline
\rowcolor{blue!20}
\textbf{STT} & \textbf{Tên đối tượng} & \textbf{Kiểu} & \textbf{Ràng buộc} & \textbf{Chức năng} \\ \hline

1 & \textbf{Logo Skinmate} & Link & - & Đưa người dùng quay về đầu trang chủ. \\ \hline

2 & \textbf{Navbar Menu} & Navigation & - & Chứa liên kết Phân tích, Cộng đồng; khi đăng nhập có menu Hồ sơ, Lịch sử và liên kết Quản trị dành cho Admin. \\ \hline

3 & \textbf{CTA Button} & Button & Auth state & Nếu chưa đăng nhập: hiển thị "Bắt Đầu Ngay" (dẫn đến Đăng ký). Nếu đã đăng nhập: hiển thị "Phân Tích Ngay" (dẫn đến Phân tích) và "Hồ Sơ Của Tôi". \\ \hline

4 & \textbf{Sign In Link} & Link & - & Chuyển hướng người dùng đến màn hình Đăng nhập. \\ \hline

5 & \textbf{Feature Cards} & Container & - & Hiển thị 03 thẻ: Phân tích AI, Lọc độ phù hợp và Gợi ý sản phẩm. \\ \hline

6 & \textbf{Health Status Bar} & Label & Cấu hình health check & Hiển thị trạng thái đồng bộ; chỉ gọi endpoint kiểm tra máy chủ khi biến \texttt{NEXT\_PUBLIC\_ENABLE\_HEALTH\_CHECK} được bật. \\ \hline

7 & \textbf{Sign Out Button} & Button & Đã Login & Thực hiện đăng xuất, xóa Token khỏi LocalStorage. \\ \hline

\end{tabularx}
\end{center}

\textbf{c) Danh sách biến cố và xử lý}

\begin{center}
\renewcommand{\arraystretch}{1.5}
\rowcolors{2}{gray!10}{white}
\begin{tabularx}{\textwidth}{|c|l|X|}
\hline
\rowcolor{blue!20}
\textbf{STT} & \textbf{Biến cố} & \textbf{Xử lý} \\ \hline

1 & \textbf{Truy cập Trang chủ} & Giao diện khởi tạo trạng thái sẵn sàng. Nếu \texttt{NEXT\_PUBLIC\_ENABLE\_HEALTH\_CHECK=true}, Hook \texttt{useEffect} gọi endpoint \texttt{/health}; lỗi kết nối được hiển thị tại thanh trạng thái. \\ \hline

2 & \textbf{Click CTA chính} & Theo trạng thái \texttt{user} trong \texttt{AuthContext}, nút hiển thị sẵn liên kết \texttt{/register} cho khách hoặc \texttt{/analysis} cho người đã đăng nhập. \\ \hline

3 & \textbf{Click "Sign In"} & Chuyển hướng người dùng sang trang \texttt{/login}. \\ \hline

4 & \textbf{Click "Logout"} & Gọi hàm \texttt{logout()} từ \texttt{useAuth}. Xóa toàn bộ dữ liệu người dùng trong LocalStorage, chuyển trạng thái hệ thống về khách (Guest) và cập nhật lại Navbar. \\ \hline

\end{tabularx}
\end{center}




\subsubsection{Màn hình Đăng ký}

\textbf{a) Giao diện}

\begin{figure}[H]
    \centering
    \includegraphics[width=0.9\linewidth]{ui_register.pdf} % Thay bằng ảnh chụp màn hình thực tế
    \caption{Giao diện Màn hình Đăng ký hệ thống SKINMATE}
    \label{fig:ui_register}
\end{figure}

\textbf{b) Mô tả các đối tượng trên màn hình}

\begin{center}
\renewcommand{\arraystretch}{1.5}
\rowcolors{2}{gray!10}{white}
\begin{tabularx}{\textwidth}{|c|l|l|l|X|}
\hline
\rowcolor{blue!20}
\textbf{STT} & \textbf{Tên đối tượng} & \textbf{Kiểu} & \textbf{Ràng buộc} & \textbf{Chức năng} \\ \hline

1 & \textbf{Input Display Name} & Textbox & Tùy chọn & Nhập họ và tên hoặc tên hiển thị của người dùng trên hệ thống. \\ \hline

2 & \textbf{Input Username} & Textbox & Bắt buộc, duy nhất & Nhập tên tài khoản dùng để đăng nhập hệ thống. \\ \hline

3 & \textbf{Select Skin Type} & Dropdown & Bắt buộc & Chọn loại da thực tế của người dùng (Bình thường, Dầu, Khô, Nhạy cảm, Hỗn hợp). \\ \hline

4 & \textbf{Input Password} & Password & Tối thiểu 6 ký tự & Nhập mật khẩu bảo mật cho tài khoản. \\ \hline

5 & \textbf{Confirm Password} & Password & Khớp Password & Nhập lại mật khẩu để kiểm tra tính chính xác tránh gõ nhầm. \\ \hline

6 & \textbf{Register Button} & Button & Không bỏ trống form & Gửi yêu cầu đăng ký lên máy chủ. Hiển thị trạng thái "Đang tạo tài khoản..." khi xử lý. \\ \hline

7 & \textbf{Error Alert} & Label & Hiển thị khi có lỗi & Thông báo các cảnh báo lỗi như: Trùng Username, mật khẩu không khớp, v.v. \\ \hline

8 & \textbf{Login Link} & Link & - & Chuyển hướng sang màn hình Đăng nhập nếu người dùng đã có tài khoản. \\ \hline

\end{tabularx}
\end{center}

\textbf{c) Danh sách biến cố và xử lý}

\begin{center}
\renewcommand{\arraystretch}{1.5}
\rowcolors{2}{gray!10}{white}
\begin{tabularx}{\textwidth}{|c|l|X|}
\hline
\rowcolor{blue!20}
\textbf{STT} & \textbf{Biến cố} & \textbf{Xử lý} \\ \hline

1 & \textbf{Thay đổi thông tin} & Hệ thống cập nhật các \textit{state} tương ứng (\texttt{username, password, skinType}) ngay khi người dùng nhập liệu. \\ \hline

2 & \textbf{Click "Create Account"} & 
\begin{itemize}
    \item Kiểm tra mật khẩu khớp nhau.
    \item Kiểm tra độ dài mật khẩu $\ge$ 6.
    \item Nếu hợp lệ: Gọi API \texttt{POST /api/v1/auth/register}.
    \item Nếu thất bại: Hiển thị thông báo lỗi màu Rose.
\end{itemize} \\ \hline

3 & \textbf{Đăng ký thành công} & Hệ thống tự động kích hoạt luồng Đăng nhập ngầm (\texttt{POST /api/v1/auth/login}) bằng thông tin vừa tạo để lấy Token, sau đó chuyển hướng về trang chủ (\texttt{/}); nếu đăng nhập ngầm thất bại thì chuyển đến \texttt{/login}. \\ \hline

4 & \textbf{Click "Sign In"} & Chuyển hướng người dùng sang trang \texttt{/login} mà không lưu lại các dữ liệu đã nhập. \\ \hline

\end{tabularx}
\end{center}




\subsubsection{Màn hình Đăng nhập}

\textbf{a) Giao diện}

\begin{figure}[H]
    \centering
    \includegraphics[width=0.9\linewidth]{ui_login.pdf} % Thay bằng ảnh chụp màn hình thực tế
    \caption{Giao diện Màn hình Đăng nhập hệ thống SKINMATE}
    \label{fig:ui_login}
\end{figure}
\textbf{b) Mô tả các đối tượng trên màn hình}

\begin{center}
\renewcommand{\arraystretch}{1.5}
\rowcolors{2}{gray!10}{white}
\begin{tabularx}{\textwidth}{|c|l|l|l|X|}
\hline
\rowcolor{blue!20}
\textbf{STT} & \textbf{Tên đối tượng} & \textbf{Kiểu} & \textbf{Ràng buộc} & \textbf{Chức năng} \\ \hline

1 & \textbf{Input Username} & Textbox & Bắt buộc & Người dùng nhập tên tài khoản đã đăng ký. \\ \hline

2 & \textbf{Input Password} & Password & Bắt buộc & Người dùng nhập mật khẩu bảo mật. \\ \hline

3 & \textbf{Sign In Button} & Button & - & Gửi thông tin đăng nhập và chuyển sang trạng thái "Đang đăng nhập..." khi đang xử lý. \\ \hline

4 & \textbf{Error Message} & Label & Khi sai thông tin & Hiển thị thông báo lỗi (ví dụ: Sai mật khẩu, tài khoản bị khóa) kèm biểu tượng cảnh báo. \\ \hline

5 & \textbf{Register Link} & Link & - & Chuyển hướng người dùng sang trang \texttt{/register} (Join Skinmate). \\ \hline

6 & \textbf{Decorative Image} & Image & - & Hình ảnh phong cách Luxury skincare bên trái màn hình giúp tăng tính thẩm mỹ. \\ \hline

\end{tabularx}
\end{center}

\textbf{c) Danh sách biến cố và xử lý}

\begin{center}
\renewcommand{\arraystretch}{1.5}
\rowcolors{2}{gray!10}{white}
\begin{tabularx}{\textwidth}{|c|l|X|}
\hline
\rowcolor{blue!20}
\textbf{STT} & \textbf{Biến cố} & \textbf{Xử lý} \\ \hline

1 & \textbf{Submit Form} & Hệ thống ngăn chặn hành động load lại trang mặc định, thiết lập trạng thái \texttt{loading = true} và xóa các thông báo lỗi cũ. \\ \hline

2 & \textbf{Gọi API Đăng nhập} & Gửi yêu cầu \texttt{POST} tới đầu cuối \texttt{/api/v1/auth/login}. \\ \hline

3 & \textbf{Đăng nhập thành công} & Hệ thống nhận về gói dữ liệu gồm \texttt{token} và \texttt{user}. Gọi hàm \texttt{login()} trong \texttt{AuthContext} để lưu dữ liệu vào LocalStorage và chuyển hướng về Trang chủ (\texttt{/}). \\ \hline

4 & \textbf{Thất bại (Sai thông tin)} & Hệ thống bắt lỗi \texttt{401} hoặc \texttt{400}, hiển thị nội dung tiếng Việt từ Server hoặc thông báo "Đăng nhập thất bại" lên \textbf{Error Message}. \\ \hline

5 & \textbf{Thất bại (Tài khoản khóa)} & Nếu nhận mã lỗi liên quan đến tài khoản bị khóa, hệ thống hiển thị thông báo yêu cầu liên hệ quản trị viên. \\ \hline

\end{tabularx}
\end{center}







\subsubsection{Màn hình Phân tích thành phần (INCI Analysis)}

\textbf{a) Giao diện}

\begin{figure}[H]
    \centering
    \includegraphics[width=0.9\linewidth]{ui_analysis.pdf} % Thay bằng ảnh chụp màn hình thực tế
    \caption{Giao diện Màn hình Phân tích thành phần hệ thống SKINMATE}
    \label{fig:ui_analysis}
\end{figure}

\textbf{b) Mô tả các đối tượng trên màn hình}

\begin{center}
\renewcommand{\arraystretch}{1.5}
\rowcolors{2}{gray!10}{white}
\begin{tabularx}{\textwidth}{|c|l|l|l|X|}
\hline
\rowcolor{blue!20}
\textbf{STT} & \textbf{Tên đối tượng} & \textbf{Kiểu} & \textbf{Ràng buộc} & \textbf{Chức năng} \\ \hline

1 & \textbf{OCR Image Uploader} & File Input & Ảnh tối đa 5 MB & Tải ảnh bảng thành phần lên API \texttt{/api/ocr/ingredients}; chuỗi trích xuất thành công được điền vào ô INCI. \\ \hline

2 & \textbf{Textarea Input (INCI)} & Text Area & Không rỗng & Nơi người dùng dán hoặc nhận chuỗi danh sách thành phần mỹ phẩm (INCI string) cần kiểm tra. \\ \hline

3 & \textbf{Skin Type Badge} & Label & Dữ liệu User & Hiển thị loại da hiện tại của người dùng để nhắc nhở kết quả phân tích sẽ dựa trên loại da này. \\ \hline

4 & \textbf{Analyze Button} & Button & Có dữ liệu input & Kích hoạt quá trình gửi chuỗi INCI lên máy chủ để phân tích. \\ \hline

5 & \textbf{Result Badges/Cards} & List Item & Chỉ hiện sau phân tích & Hiển thị kết quả từng thành phần kèm mã màu: Xanh (Tốt/GOOD), Đỏ (Xấu/BAD), Xám (Trung tính/NEUTRAL). \\ \hline

6 & \textbf{Ingredient Description} & Label & Dữ liệu từ DB/AI & Hiển thị công dụng hoặc cảnh báo chi tiết cho từng thành phần tương ứng với loại da. \\ \hline

7 & \textbf{Recommendation Section} & Container & Phụ thuộc kết quả & Hiển thị tối đa 3 sản phẩm gợi ý không chứa thành phần BAD, được chọn từ nhóm sản phẩm an toàn có điểm GOOD cao. \\ \hline

8 & \textbf{Product Card} & Card & - & Hiển thị hình ảnh, tên và thương hiệu của sản phẩm được gợi ý. \\ \hline

\end{tabularx}
\end{center}

\textbf{c) Danh sách biến cố và xử lý}

\begin{center}
\renewcommand{\arraystretch}{1.5}
\rowcolors{2}{gray!10}{white}
\begin{tabularx}{\textwidth}{|c|l|X|}
\hline
\rowcolor{blue!20}
\textbf{STT} & \textbf{Biến cố} & \textbf{Xử lý} \\ \hline

1 & \textbf{Truy cập màn hình} & Trang kiểm tra trạng thái \texttt{user} từ \texttt{AuthContext}. Nếu chưa đăng nhập, hiển thị vùng thông báo hạn chế truy cập kèm liên kết Đăng nhập/Đăng ký; nếu đã đăng nhập, hiển thị biểu mẫu phân tích. \\ \hline

2 & \textbf{Tải ảnh OCR} & Người dùng chọn ảnh tại \texttt{ImageOCRUploader}; Frontend gửi \texttt{POST /api/ocr/ingredients} bằng \texttt{FormData}. Nếu nhận được \texttt{ingredients}, hệ thống điền chuỗi vào \texttt{inciString}; nếu lỗi, hiển thị cảnh báo trích xuất. \\ \hline

3 & \textbf{Click "Analyze"} & 
\begin{itemize}
    \item Kiểm tra chuỗi nhập vào. Nếu hợp lệ, đổi trạng thái nút thành "Analyzing..." (Loading).
    \item Gọi API \texttt{POST /api/v1/analysis/check} kèm dữ liệu INCI.
    \item Hệ thống Backend sẽ tra cứu Database, nếu gặp thành phần lạ sẽ tự động gọi \textbf{Gemini AI} để đánh giá và cache lại.
    \item Lịch sử phân tích tự động được Backend lưu vào bảng \texttt{AnalysisHistory}.
\end{itemize} \\ \hline

4 & \textbf{Xử lý kết quả phân tích} & Nhận mảng kết quả từ máy chủ. Hệ thống tiến hành \textit{render} danh sách thành phần. Tùy thuộc vào trường \texttt{effect}, áp dụng CSS class tương ứng: \texttt{bg-emerald-50} (GOOD), \texttt{bg-rose-50} (BAD), \texttt{bg-gray-50} (NEUTRAL). \\ \hline

5 & \textbf{Tự động tải Gợi ý sản phẩm} & Ngay sau khi có kết quả phân tích thành công, Frontend tự động gọi API \texttt{GET /api/v1/products/recommendations}. Backend loại sản phẩm có chất BAD, chấm điểm GOOD, lấy nhóm tối đa 6 sản phẩm điểm cao và xáo trộn chọn tối đa 3 sản phẩm hiển thị. \\ \hline

6 & \textbf{Hạn mức Rate Limit} & Nếu người dùng thường vượt quá 25 lần phân tích trong 24 giờ, API trả về lỗi 429; tài khoản \texttt{ADMIN} được bỏ qua giới hạn này. Frontend hiển thị thông báo lỗi do API trả về. \\ \hline

\end{tabularx}
\end{center}









\subsubsection{Màn hình Hồ sơ cá nhân (Profile)}

\textbf{a) Giao diện}

\begin{figure}[H]
    \centering
    \includegraphics[width=0.9\linewidth]{ui_profile.pdf} % Thay bằng ảnh chụp màn hình thực tế
    \caption{Giao diện Màn hình Hồ sơ cá nhân hệ thống SKINMATE}
    \label{fig:ui_profile}
\end{figure}

\textbf{b) Mô tả các đối tượng trên màn hình}

\begin{center}
\renewcommand{\arraystretch}{1.5}
\rowcolors{2}{gray!10}{white}
\begin{tabularx}{\textwidth}{|c|l|l|l|X|}
\hline
\rowcolor{blue!20}
\textbf{STT} & \textbf{Tên đối tượng} & \textbf{Kiểu} & \textbf{Ràng buộc} & \textbf{Chức năng} \\ \hline

1 & \textbf{Input Username} & Textbox & Read-only & Hiển thị tên đăng nhập của người dùng. Trường này bị vô hiệu hóa (không cho phép chỉnh sửa). \\ \hline

2 & \textbf{Input Display Name} & Textbox & Tùy chọn & Cho phép cập nhật tên hiển thị của người dùng trên thanh điều hướng (Navbar). \\ \hline

3 & \textbf{Select Skin Type} & Dropdown & Bắt buộc & Cho phép thay đổi loại da. Đây là trường dữ liệu cốt lõi ảnh hưởng trực tiếp đến kết quả phân tích AI và gợi ý sản phẩm. \\ \hline

4 & \textbf{Save Changes} & Button & Khóa nếu chưa đổi & Nút lưu cập nhật hồ sơ (Chỉ sáng lên khi người dùng có thay đổi dữ liệu so với ban đầu). \\ \hline

5 & \textbf{Toggle Password} & Button & - & Nút "Đổi Mật Khẩu" / "Hủy" để mở rộng hoặc thu gọn khu vực Cài đặt bảo mật. \\ \hline

6 & \textbf{Password Inputs} & Password & Bắt buộc khi đổi & Gồm 3 trường: Mật khẩu hiện tại, Mật khẩu mới, Xác nhận mật khẩu mới. \\ \hline

7 & \textbf{Update Password} & Button & Pass $\ge$ 6 ký tự & Nút xác nhận gửi yêu cầu thay đổi mật khẩu lên máy chủ. \\ \hline

8 & \textbf{Notification Banner} & Label & Tùy ngữ cảnh & Thanh thông báo trạng thái (Xanh: Thành công, Đỏ: Thất bại) xuất hiện phía trên form. \\ \hline

\end{tabularx}
\end{center}

\textbf{c) Danh sách biến cố và xử lý}

\begin{center}
\renewcommand{\arraystretch}{1.5}
\rowcolors{2}{gray!10}{white}
\begin{tabularx}{\textwidth}{|c|l|X|}
\hline
\rowcolor{blue!20}
\textbf{STT} & \textbf{Biến cố} & \textbf{Xử lý} \\ \hline

1 & \textbf{Truy cập màn hình} & Component \texttt{ProtectedRoute} kiểm tra người dùng trong \texttt{AuthContext}; nếu chưa đăng nhập thì chuyển về \texttt{/login}. Khi đã đăng nhập, gọi API \texttt{GET /api/v1/users/profile} để lấy dữ liệu cá nhân (Skin Type, Display Name) và điền sẵn vào các trường (Pre-fill). Bật trạng thái \textit{Loading} trong lúc chờ phản hồi. \\ \hline

2 & \textbf{Click "Lưu Thay Đổi"} & 
\begin{itemize}
    \item Gọi API \texttt{PUT /api/v1/users/profile} kèm dữ liệu mới.
    \item Nếu thành công: Hiển thị thông báo màu Xanh. \textbf{Đồng thời gọi hàm \texttt{login()} trong \texttt{AuthContext}} để cập nhật lại dữ liệu \texttt{user} trong bộ nhớ cục bộ, giúp tên và loại da trên Navbar thay đổi ngay lập tức mà không cần tải lại trang (Reactive UI).
\end{itemize} \\ \hline

3 & \textbf{Click "Cập Nhật Mật Khẩu"} & 
\begin{itemize}
    \item Kiểm tra Client-side: Xác nhận mật khẩu mới khớp nhau và độ dài $\ge$ 6 ký tự. Nếu sai, hiển thị thông báo lỗi ngay lập tức.
    \item Nếu hợp lệ: Gọi API \texttt{PUT /api/v1/users/change-password}.
    \item Hệ thống Backend sẽ so sánh mã băm (hash) của mật khẩu cũ. Nếu đúng, tiến hành cập nhật mật khẩu mới.
    \item Nếu thành công: Hiển thị thông báo, xóa rỗng các ô nhập liệu và tự động thu gọn khu vực đổi mật khẩu.
\end{itemize} \\ \hline

4 & \textbf{Click "Quay lại"} & Hệ thống kích hoạt hàm \texttt{window.history.back()} để đưa người dùng trở về trang trước đó một cách mượt mà. \\ \hline

\end{tabularx}
\end{center}





\subsubsection{Màn hình Lịch sử phân tích (History)}

\textbf{a) Giao diện}

\begin{figure}[H]
    \centering
    \includegraphics[width=0.9\linewidth]{ui_history.pdf} % Thay bằng ảnh chụp màn hình thực tế
    \caption{Giao diện Màn hình Lịch sử hệ thống SKINMATE}
    \label{fig:ui_history}
\end{figure}

\textbf{b) Mô tả các đối tượng trên màn hình}

\begin{center}
\renewcommand{\arraystretch}{1.5}
\rowcolors{2}{gray!10}{white}
\begin{tabularx}{\textwidth}{|c|l|l|l|X|}
\hline
\rowcolor{blue!20}
\textbf{STT} & \textbf{Tên đối tượng} & \textbf{Kiểu} & \textbf{Ràng buộc} & \textbf{Chức năng} \\ \hline

1 & \textbf{Back to Home Link} & Link & - & Nút "Về Trang Chủ" giúp điều hướng nhanh về màn hình chính. \\ \hline

2 & \textbf{Clear All Button} & Button & Có dữ liệu lịch sử & Nút "Xóa Tất Cả" cho phép xóa toàn bộ bản ghi phân tích của người dùng hiện tại. \\ \hline

3 & \textbf{History Card} & Container & - & Thẻ hiển thị thông tin tóm tắt của một lần phân tích (Bao gồm thời gian và danh sách thành phần INCI đã nhập). \\ \hline

4 & \textbf{Re-analyze Button} & Button & - & Nút "Phân Tích Lại" trên từng thẻ lịch sử để chạy lại quá trình phân tích với dữ liệu cũ. \\ \hline

5 & \textbf{Delete Item Button} & Button & - & Nút "Xóa" (biểu tượng thùng rác) để xóa một bản ghi cụ thể. \\ \hline

6 & \textbf{Empty/Guest State} & Container & Không có dữ liệu / Chưa login & Hiển thị thông báo yêu cầu đăng nhập (kèm nút Đăng nhập) hoặc thông báo chưa có lịch sử (kèm nút Phân tích ngay). \\ \hline

7 & \textbf{Error Banner} & Label & Khi lỗi API & Hiển thị thông báo lỗi nếu quá trình tải hoặc xóa dữ liệu gặp sự cố. \\ \hline

\end{tabularx}
\end{center}

\textbf{c) Danh sách biến cố và xử lý}

\begin{center}
\renewcommand{\arraystretch}{1.5}
\rowcolors{2}{gray!10}{white}
\begin{tabularx}{\textwidth}{|c|l|X|}
\hline
\rowcolor{blue!20}
\textbf{STT} & \textbf{Biến cố} & \textbf{Xử lý} \\ \hline

1 & \textbf{Truy cập màn hình} & 
\begin{itemize}
    \item Nếu chưa đăng nhập: Trình bày màn hình \textit{Guest State} chặn truy cập.
    \item Nếu đã đăng nhập: Hiển thị \textit{Loading Spinner} và gọi API \texttt{GET /api/v1/history} kèm Token. Sau khi có dữ liệu, \textit{render} danh sách các History Card theo thứ tự mới nhất xếp trên.
\end{itemize} \\ \hline

2 & \textbf{Click "Phân Tích Lại"} & Hệ thống lấy chuỗi thành phần gốc (\texttt{rawInput}) của thẻ tương ứng, mã hóa nó (\texttt{encodeURIComponent}) và đẩy vào URL Query. Chuyển hướng người dùng sang trang \texttt{/analysis?inci=...} để điền sẵn chuỗi cũ; người dùng nhấn phân tích để thực hiện lại. \\ \hline

3 & \textbf{Click "Xóa" (Một thẻ)} & 
\begin{itemize}
    \item Bật hộp thoại xác nhận \texttt{confirm()} của trình duyệt. Nếu đồng ý, đổi icon thùng rác thành \textit{Loading Spinner} cho riêng thẻ đó.
    \item Gọi API \texttt{DELETE /api/v1/history/:id}. 
    \item Nếu thành công, tự động loại bỏ thẻ đó khỏi danh sách giao diện bằng cách lọc mảng trạng thái (filter state) mà không cần tải lại trang.
\end{itemize} \\ \hline

4 & \textbf{Click "Xóa Tất Cả"} & Bật hộp thoại cảnh báo nghiêm trọng. Nếu đồng ý, gọi API \texttt{DELETE /api/v1/history}. Sau khi thành công, dọn sạch mảng dữ liệu và hiển thị giao diện \textit{Empty State} (Chưa có lịch sử). \\ \hline

\end{tabularx}
\end{center}











\subsubsection{Màn hình Admin Dashboard (Bảng điều khiển Quản trị)}

\textbf{a) Giao diện}

\begin{figure}[H]
    \centering
    \includegraphics[width=0.9\linewidth]{ui_admin_dashboard.pdf} % Thay bằng ảnh chụp màn hình thực tế
    \caption{Giao diện Màn hình Admin Dashboard hệ thống SKINMATE}
    \label{fig:ui_admin_dashboard}
\end{figure}

\textbf{b) Mô tả các đối tượng trên màn hình}

\begin{center}
\renewcommand{\arraystretch}{1.5}
\rowcolors{2}{gray!10}{white}
\begin{tabularx}{\textwidth}{|c|l|l|l|X|}
\hline
\rowcolor{blue!20}
\textbf{STT} & \textbf{Tên đối tượng} & \textbf{Kiểu} & \textbf{Ràng buộc} & \textbf{Chức năng} \\ \hline

1 & \textbf{Sidebar Menu} & Navigation & Cố định bên trái & Danh mục điều hướng chính của quản trị viên (Bảng điều khiển, Thành phần, Quy tắc an toàn, Sản phẩm, Người dùng, Báo cáo, BC Cộng đồng, Nhập/Xuất). \\ \hline

2 & \textbf{Active Menu Item} & Label & Phụ thuộc URL & Làm nổi bật (Highlight màu xanh lá) module mà quản trị viên đang truy cập. \\ \hline

3 & \textbf{Main Content Area} & Container & - & Vùng hiển thị nội dung động của từng phân hệ quản trị tương ứng khi người dùng bấm vào Sidebar. \\ \hline

4 & \textbf{Dashboard Title} & Label & - & Hiển thị tiêu đề "Bảng Điều Khiển". \\ \hline

5 & \textbf{Quick Overview Cards} & Card & - & Bộ 3 thẻ thông tin (Thành phần, Quy tắc an toàn, Sản phẩm) giúp quản trị viên nắm bắt nhanh các cấu phần cốt lõi của hệ thống. \\ \hline

\end{tabularx}
\end{center}

\textbf{c) Danh sách biến cố và xử lý}

\begin{center}
\renewcommand{\arraystretch}{1.5}
\rowcolors{2}{gray!10}{white}
\begin{tabularx}{\textwidth}{|c|l|X|}
\hline
\rowcolor{blue!20}
\textbf{STT} & \textbf{Biến cố} & \textbf{Xử lý} \\ \hline

1 & \textbf{Truy cập phân hệ Admin} & 
\begin{itemize}
    \item Hệ thống kích hoạt Component \texttt{AdminProtectedRoute}.
    \item \textbf{Kiểm tra Auth:} Đợi trạng thái \texttt{isLoading} hoàn tất. Nếu \texttt{user = null}, tự động đẩy về trang \texttt{/login}.
    \item \textbf{Kiểm tra Role:} Component kiểm tra \texttt{user.role} trong \texttt{AuthContext}. Nếu \texttt{role !== 'ADMIN'}, giao diện không render nội dung quản trị và chuyển hướng về trang chủ (\texttt{/}); các API quản trị tiếp tục được bảo vệ bởi \texttt{adminMiddleware} ở Backend.
\end{itemize} \\ \hline

2 & \textbf{Chuyển hướng (Click Sidebar)} & 
\begin{itemize}
    \item Hệ thống sử dụng Next.js \texttt{Router} để thay đổi URL (ví dụ: chuyển sang \texttt{/admin/ingredients}).
    \item Hook \texttt{usePathname()} bắt được sự thay đổi URL, tự động cập nhật CSS để bôi đậm mục đang chọn (Active state) trên Sidebar.
    \item Nội dung của phân hệ mới được \textit{render} vào Main Content Area mà không cần tải lại (reload) toàn bộ trang (Nhờ cơ chế Layout của Next.js).
\end{itemize} \\ \hline

3 & \textbf{Tương tác thẻ (Hover Cards)} & Áp dụng hiệu ứng CSS Transitions: Khi di chuột qua các \textit{Quick Overview Cards}, thẻ tự động nổi lên (\texttt{-translate-y-0.5}) và đổ bóng đậm hơn để tăng trải nghiệm UX. \\ \hline

\end{tabularx}
\end{center}







\subsubsection{Màn hình Quản lý Thành phần (Ingredients)}

\textbf{a) Giao diện}

\begin{figure}[H]
    \centering
    \includegraphics[width=0.9\linewidth]{ui_admin_ingredients.png} % Thay bằng ảnh chụp màn hình thực tế
    \caption{Giao diện Màn hình Quản lý Thành phần hệ thống SKINMATE}
    \label{fig:ui_admin_ingredients}
\end{figure}

\textbf{b) Mô tả các đối tượng trên màn hình}

\begin{center}
\renewcommand{\arraystretch}{1.5}
\rowcolors{2}{gray!10}{white}
\begin{tabularx}{\textwidth}{|c|l|l|l|X|}
\hline
\rowcolor{blue!20}
\textbf{STT} & \textbf{Tên đối tượng} & \textbf{Kiểu} & \textbf{Ràng buộc} & \textbf{Chức năng} \\ \hline

1 & \textbf{Add New Button} & Button & - & Nút mở hộp thoại (Modal) để thêm một thành phần mỹ phẩm mới vào cơ sở dữ liệu. \\ \hline

2 & \textbf{Search Input} & Textbox & - & Thanh tìm kiếm gửi từ khóa lên endpoint danh sách để tra cứu thành phần theo tên. \\ \hline

3 & \textbf{Data Table} & Table & - & Bảng hiển thị danh sách thành phần. Các cột bao gồm: ID, Tên (Name), Mô tả (Description) và Cột thao tác (Actions). \\ \hline

4 & \textbf{Edit / Delete} & Button & Theo từng dòng & Các nút thao tác trên từng bản ghi. \textit{Edit} mở Modal điền sẵn dữ liệu cũ, \textit{Delete} yêu cầu xác nhận xóa. \\ \hline

5 & \textbf{Modal Window} & Popup & - & Hộp thoại nổi (Overlay) chứa Form dùng chung cho cả chức năng Thêm mới và Cập nhật. \\ \hline

6 & \textbf{Input Name (INCI)} & Textbox & Bắt buộc & Trường nhập tên thành phần mỹ phẩm (Hiển thị trong Modal). \\ \hline

7 & \textbf{Input Description} & Text Area & Tùy chọn & Trường nhập mô tả, công dụng chi tiết của thành phần (Hiển thị trong Modal). \\ \hline

8 & \textbf{Error/Alert Banner} & Label & Khi có lỗi & Thanh cảnh báo hiển thị mã lỗi trả về từ máy chủ (VD: Thành phần đã tồn tại). \\ \hline

9 & \textbf{Pagination Controls} & Button & Khi có nhiều hơn 1 trang & Chuyển trang danh sách; mỗi yêu cầu tải tối đa 15 thành phần. \\ \hline

\end{tabularx}
\end{center}

\textbf{c) Danh sách biến cố và xử lý}

\begin{center}
\renewcommand{\arraystretch}{1.5}
\rowcolors{2}{gray!10}{white}
\begin{tabularx}{\textwidth}{|c|l|X|}
\hline
\rowcolor{blue!20}
\textbf{STT} & \textbf{Biến cố} & \textbf{Xử lý} \\ \hline

1 & \textbf{Tải danh sách} & Khi truy cập, hệ thống gọi API \texttt{GET /api/v1/admin/ingredients} kèm \texttt{Bearer Token} và các query \texttt{page}, \texttt{limit=15}, \texttt{search}. Dữ liệu phân trang được lưu vào state và hiển thị lên bảng. \\ \hline

2 & \textbf{Nhập từ khóa tìm kiếm / chuyển trang} & Cập nhật \texttt{searchTerm} hoặc \texttt{currentPage}; sau độ trễ 250 ms, Frontend gọi lại API danh sách. Khi nhập từ khóa mới, trang được đặt lại về trang 1; Backend tìm tên không phân biệt hoa/thường. \\ \hline

3 & \textbf{Mở Modal Thêm/Sửa} & 
\begin{itemize}
    \item \textbf{Nếu Thêm mới:} Xóa rỗng \texttt{formData}, đặt biến trạng thái \texttt{editingIngredient = null}.
    \item \textbf{Nếu Sửa:} Đổ dữ liệu của dòng được chọn vào \texttt{formData}, lưu ID vào \texttt{editingIngredient}.
    \item Đổi trạng thái \texttt{isModalOpen = true} để bật Overlay mờ và hiển thị Form.
\end{itemize} \\ \hline

4 & \textbf{Submit Form (Lưu)} & 
Hệ thống chặn reload trang. Dựa vào trạng thái \texttt{editingIngredient}:
\begin{itemize}
    \item Gửi request \texttt{POST} (nếu thêm mới) hoặc \texttt{PUT} (nếu cập nhật) lên API tương ứng.
    \item Nếu máy chủ báo lỗi (Ví dụ: 409 Conflict - trùng tên), hiển thị thông báo lỗi màu đỏ.
    \item Nếu thành công, tự động đóng Modal, làm sạch form và gọi lại hàm \texttt{fetchIngredients()} để làm mới bảng dữ liệu.
\end{itemize} \\ \hline

5 & \textbf{Click "Delete"} & Hiển thị hàm cảnh báo \texttt{confirm()} của trình duyệt. Nếu Admin đồng ý, gọi API \texttt{DELETE /api/v1/admin/ingredients/:id}. Sau khi xóa thành công, tự động tải lại danh sách. \\ \hline

\end{tabularx}
\end{center}




\subsubsection{Màn hình Quản lý Quy tắc An toàn (Safety Rules)}

\textbf{a) Giao diện}

\begin{figure}[H]
    \centering
    \includegraphics[width=0.9\linewidth]{ui_admin_rules.png} % Thay bằng ảnh chụp màn hình thực tế
    \caption{Giao diện Màn hình Quản lý Quy tắc An toàn hệ thống SKINMATE}
    \label{fig:ui_admin_rules}
\end{figure}

\textbf{b) Mô tả các đối tượng trên màn hình}

\begin{center}
\renewcommand{\arraystretch}{1.5}
\rowcolors{2}{gray!10}{white}
\begin{tabularx}{\textwidth}{|c|l|l|l|X|}
\hline
\rowcolor{blue!20}
\textbf{STT} & \textbf{Tên đối tượng} & \textbf{Kiểu} & \textbf{Ràng buộc} & \textbf{Chức năng} \\ \hline

1 & \textbf{Search Input} & Textbox & - & Thanh tìm kiếm gửi từ khóa lên server để tra cứu các quy tắc theo tên thành phần. \\ \hline

2 & \textbf{Inline Form} & Form & Nằm cố định trên cùng & Khu vực nhập liệu trực tiếp để Thêm mới hoặc Cập nhật quy tắc (Không dùng Modal nổi). \\ \hline

3 & \textbf{Select Ingredient} & Dropdown & Bắt buộc & Chọn thành phần mỹ phẩm từ danh sách đã được lưu trong cơ sở dữ liệu. \\ \hline

4 & \textbf{Select Skin Type} & Dropdown & Bắt buộc & Chọn loại da cần áp dụng quy tắc (NORMAL, OILY, DRY, SENSITIVE, COMBINATION). \\ \hline

5 & \textbf{Select Effect} & Dropdown & Bắt buộc & Chọn mức độ đánh giá: GOOD (Tốt), BAD (Kích ứng), NEUTRAL (Trung tính). \\ \hline

6 & \textbf{Save Rule Button} & Button & - & Gửi thông tin (Thành phần + Loại da + Đánh giá) lên máy chủ để thiết lập quy tắc. \\ \hline

7 & \textbf{Data Table} & Table & - & Bảng danh sách các quy tắc hiện có. Cột \textit{Effect} được gắn nhãn màu trực quan (Xanh/Đỏ/Xám). \\ \hline

8 & \textbf{Edit / Delete} & Button & Theo từng dòng & Nhấn \textit{Edit} sẽ cuộn trang lên trên và đổ dữ liệu vào Form. \textit{Delete} sẽ xóa quy tắc. \\ \hline

9 & \textbf{Pagination Controls} & Button & Khi có nhiều hơn 1 trang & Chuyển trang bảng quy tắc; mỗi yêu cầu hiển thị tối đa 15 bản ghi. \\ \hline

\end{tabularx}
\end{center}

\textbf{c) Danh sách biến cố và xử lý}

\begin{center}
\renewcommand{\arraystretch}{1.5}
\rowcolors{2}{gray!10}{white}
\begin{tabularx}{\textwidth}{|c|l|X|}
\hline
\rowcolor{blue!20}
\textbf{STT} & \textbf{Biến cố} & \textbf{Xử lý} \\ \hline

1 & \textbf{Khởi tạo màn hình} & 
Hệ thống sử dụng kỹ thuật \texttt{Promise.all} để gọi song song 02 API:
\begin{itemize}
    \item \texttt{GET /api/v1/admin/ingredients} (Lấy danh sách thành phần đổ vào Dropdown).
    \item \texttt{GET /api/v1/admin/rules?page=\ldots\&limit=15\&search=\ldots} (Lấy danh sách quy tắc phân trang để vẽ Bảng).
\end{itemize}
Khi cả 2 luồng tải xong, giao diện mới được hiển thị, tự động chọn ID của thành phần đầu tiên vào biến \texttt{formData}. \\ \hline

2 & \textbf{Nhập từ khóa tìm kiếm / chuyển trang} & Cập nhật state và gọi lại endpoint \texttt{/api/v1/admin/rules} sau độ trễ 250 ms; Backend thực hiện tìm kiếm không phân biệt hoa/thường theo tên thành phần và trả dữ liệu phân trang. \\ \hline

3 & \textbf{Click "Edit" (Sửa)} & Hệ thống tự động đổ dữ liệu của dòng tương ứng (\texttt{ingredientId}, \texttt{skinType}, \texttt{effect}) vào trạng thái \texttt{formData} của Form phía trên. Đồng thời, gọi lệnh \texttt{window.scrollTo} để cuộn màn hình mượt mà lên trên cùng, giúp Admin dễ dàng thao tác. \\ \hline

4 & \textbf{Submit Form (Lưu)} & 
Gửi yêu cầu \texttt{POST /api/v1/admin/rules}. Tại tầng Backend, máy chủ sử dụng cơ chế \textbf{Upsert} (Tìm \texttt{ingredientId} và \texttt{skinType}):
\begin{itemize}
    \item Nếu cặp khóa này đã tồn tại: Tiến hành cập nhật \texttt{effect}.
    \item Nếu chưa tồn tại: Tạo bản ghi mới.
\end{itemize}
Sau khi có kết quả, Frontend hiển thị thông báo "Rule created/updated successfully!" và gọi lại hàm \texttt{fetchData()} để làm mới bảng. \\ \hline

5 & \textbf{Click "Delete" (Xóa)} & Bật thông báo xác nhận \texttt{confirm()}. Nếu đồng ý, gọi API \texttt{DELETE /api/v1/admin/rules/:id}. Sau khi thành công, bảng dữ liệu tự động được làm mới. \\ \hline

\end{tabularx}
\end{center}


\subsubsection{Màn hình Quản lý Sản phẩm (Products)}

\textbf{a) Giao diện}

Màn hình được triển khai tại đường dẫn \texttt{/admin/products}, gồm biểu mẫu thêm/chỉnh sửa sản phẩm, vùng nhập chuỗi INCI và bảng dữ liệu có tìm kiếm, phân trang.

\textbf{b) Mô tả các đối tượng trên màn hình}

\begin{center}
\renewcommand{\arraystretch}{1.5}
\rowcolors{2}{gray!10}{white}
\begin{tabularx}{\textwidth}{|c|l|l|l|X|}
\hline
\rowcolor{blue!20}
\textbf{STT} & \textbf{Tên đối tượng} & \textbf{Kiểu} & \textbf{Ràng buộc} & \textbf{Chức năng} \\ \hline

1 & \textbf{Search Input} & Textbox & - & Tìm kiếm sản phẩm theo tên hoặc thương hiệu và đưa danh sách về trang đầu tiên. \\ \hline

2 & \textbf{Product Form} & Form & Tên, thương hiệu bắt buộc & Nhập tên sản phẩm, thương hiệu và đường dẫn hình ảnh tùy chọn để tạo mới hoặc cập nhật sản phẩm. \\ \hline

3 & \textbf{INCI Textarea} & Text Area & Phân cách bằng dấu phẩy & Nhập danh sách thành phần của sản phẩm; hệ thống tách chuỗi thành mảng tên thành phần khi lưu. \\ \hline

4 & \textbf{Save / Cancel Edit Button} & Button & - & Lưu dữ liệu sản phẩm hoặc hủy chế độ chỉnh sửa và xóa dữ liệu đang nhập. \\ \hline

5 & \textbf{Products Table} & Table & - & Hiển thị tên, ảnh, thương hiệu, danh sách thành phần cùng thao tác Sửa/Xóa của từng sản phẩm. \\ \hline

6 & \textbf{Pagination Controls} & Button & Khi có nhiều hơn 1 trang & Điều hướng giữa các trang dữ liệu, mỗi trang tối đa 15 sản phẩm. \\ \hline

\end{tabularx}
\end{center}

\textbf{c) Danh sách biến cố và xử lý}

\begin{center}
\renewcommand{\arraystretch}{1.5}
\rowcolors{2}{gray!10}{white}
\begin{tabularx}{\textwidth}{|c|l|X|}
\hline
\rowcolor{blue!20}
\textbf{STT} & \textbf{Biến cố} & \textbf{Xử lý} \\ \hline

1 & \textbf{Tải hoặc tìm kiếm sản phẩm} & Gọi API \texttt{GET /api/v1/admin/products} kèm tham số \texttt{page}, \texttt{limit=15}, \texttt{search} và Bearer Token; dữ liệu trả về được đưa vào bảng và tính số trang. \\ \hline

2 & \textbf{Submit Form} & Nếu chưa chọn sản phẩm, gửi \texttt{POST /api/v1/admin/products}; nếu đang sửa, gửi \texttt{PUT /api/v1/admin/products/:id}. Chuỗi INCI được tách theo dấu phẩy trước khi gửi trường \texttt{ingredientNames}. \\ \hline

3 & \textbf{Click "Sửa"} & Điền dữ liệu sản phẩm và danh sách thành phần vào biểu mẫu, lưu \texttt{editingId}, sau đó cuộn lên đầu trang để Admin chỉnh sửa. \\ \hline

4 & \textbf{Click "Xóa"} & Hiển thị xác nhận của trình duyệt; khi đồng ý, gọi \texttt{DELETE /api/v1/admin/products/:id} và tải lại danh sách sau khi thành công. \\ \hline

\end{tabularx}
\end{center}


\subsubsection{Màn hình Quản lý Người dùng (Users)}

\textbf{a) Giao diện}

Màn hình tại \texttt{/admin/users} trình bày danh sách tài khoản và popup xác nhận cho thao tác khóa, mở khóa hoặc xóa người dùng.

\textbf{b) Mô tả các đối tượng trên màn hình}

\begin{center}
\renewcommand{\arraystretch}{1.5}
\rowcolors{2}{gray!10}{white}
\begin{tabularx}{\textwidth}{|c|l|l|l|X|}
\hline
\rowcolor{blue!20}
\textbf{STT} & \textbf{Tên đối tượng} & \textbf{Kiểu} & \textbf{Ràng buộc} & \textbf{Chức năng} \\ \hline

1 & \textbf{Users Table} & Table & - & Liệt kê tên người dùng, loại da, vai trò, trạng thái hoạt động và ngày tham gia. \\ \hline

2 & \textbf{Status Badge} & Label & Theo dữ liệu \texttt{isActive} & Thể hiện tài khoản đang Hoạt động hoặc Đã khóa bằng màu nhận diện. \\ \hline

3 & \textbf{Lock / Unlock Button} & Button & Không áp dụng cho ADMIN & Yêu cầu thay đổi trạng thái truy cập của tài khoản người dùng. \\ \hline

4 & \textbf{Delete Button} & Button & Không áp dụng cho ADMIN & Yêu cầu xóa vĩnh viễn tài khoản được chọn. \\ \hline

5 & \textbf{Confirmation Modal} & Popup & Có thao tác được chọn & Hiển thị người dùng và hành động cần xác nhận trước khi gửi yêu cầu thay đổi. \\ \hline

\end{tabularx}
\end{center}

\textbf{c) Danh sách biến cố và xử lý}

\begin{center}
\renewcommand{\arraystretch}{1.5}
\rowcolors{2}{gray!10}{white}
\begin{tabularx}{\textwidth}{|c|l|X|}
\hline
\rowcolor{blue!20}
\textbf{STT} & \textbf{Biến cố} & \textbf{Xử lý} \\ \hline

1 & \textbf{Khởi tạo màn hình} & Khi có Token, gọi \texttt{GET /api/v1/admin/users}; trong thời gian chờ hiển thị thông báo đang tải, sau đó dựng các dòng tài khoản. \\ \hline

2 & \textbf{Click "Khóa" / "Mở khóa"} & Mở Modal xác nhận. Khi Admin xác nhận, gọi \texttt{PATCH /api/v1/admin/users/:id/status}, tải lại danh sách và đóng Modal. \\ \hline

3 & \textbf{Click "Xóa"} & Mở Modal cảnh báo thao tác không thể hoàn tác. Khi xác nhận, gửi \texttt{DELETE /api/v1/admin/users/:id} và cập nhật bảng dữ liệu. \\ \hline

4 & \textbf{Tài khoản Admin} & Giao diện thay các nút thao tác bằng nhãn ``Hệ thống bảo vệ'', không cho phép khóa hoặc xóa tài khoản có vai trò \texttt{ADMIN}. \\ \hline

\end{tabularx}
\end{center}


\subsubsection{Màn hình Báo cáo Thống kê Hệ thống (Reports)}

\textbf{a) Giao diện}

Màn hình tại \texttt{/admin/reports} hiển thị các thẻ tổng hợp, biểu đồ phân bổ loại da và chức năng xuất dữ liệu báo cáo dưới dạng CSV.

\textbf{b) Mô tả các đối tượng trên màn hình}

\begin{center}
\renewcommand{\arraystretch}{1.5}
\rowcolors{2}{gray!10}{white}
\begin{tabularx}{\textwidth}{|c|l|l|l|X|}
\hline
\rowcolor{blue!20}
\textbf{STT} & \textbf{Tên đối tượng} & \textbf{Kiểu} & \textbf{Ràng buộc} & \textbf{Chức năng} \\ \hline

1 & \textbf{Total Users Card} & Card & Dữ liệu thống kê & Hiển thị tổng số người dùng của hệ thống. \\ \hline

2 & \textbf{Total Analyses Card} & Card & Dữ liệu thống kê & Hiển thị tổng số lượt phân tích thành phần đã thực hiện. \\ \hline

3 & \textbf{Skin Type Pie Chart} & Chart & Dữ liệu phân bổ & Biểu diễn tỷ lệ các loại da bằng biểu đồ tròn, chú giải và tooltip. \\ \hline

4 & \textbf{Insight Panel} & Container & - & Trình bày ghi chú về tăng trưởng và chất lượng dữ liệu của nền tảng. \\ \hline

5 & \textbf{Export CSV Button} & Button & Đã có dữ liệu & Tải tệp \texttt{skinmate\_reports.csv} chứa phân bổ loại da. \\ \hline

\end{tabularx}
\end{center}

\textbf{c) Danh sách biến cố và xử lý}

\begin{center}
\renewcommand{\arraystretch}{1.5}
\rowcolors{2}{gray!10}{white}
\begin{tabularx}{\textwidth}{|c|l|X|}
\hline
\rowcolor{blue!20}
\textbf{STT} & \textbf{Biến cố} & \textbf{Xử lý} \\ \hline

1 & \textbf{Truy cập màn hình} & Khi Token sẵn sàng, gọi \texttt{GET /api/v1/admin/reports}. Kết quả gồm \texttt{totalUsers}, \texttt{totalAnalyses} và \texttt{skinTypeDistribution}. \\ \hline

2 & \textbf{Hiển thị biểu đồ} & Component \texttt{AdminReportsCharts} được tải động và dùng \texttt{recharts} để dựng biểu đồ tròn theo số lượng người dùng của từng loại da. \\ \hline

3 & \textbf{Click "Xuất Dữ Liệu (CSV)"} & Chuyển mảng \texttt{skinTypeDistribution} thành nội dung CSV gồm hai cột \texttt{Type, Count}, tạo liên kết tải xuống và lưu tệp trên thiết bị Admin. \\ \hline

\end{tabularx}
\end{center}


\subsubsection{Màn hình Nhập / Xuất Dữ liệu (Import / Export)}

\textbf{a) Giao diện}

Màn hình tại \texttt{/admin/import-export} cung cấp ba nhóm chức năng cho Thành phần, Quy tắc an toàn và Sản phẩm: xuất Excel, nhập Excel và xóa toàn bộ dữ liệu.

\textbf{b) Mô tả các đối tượng trên màn hình}

\begin{center}
\renewcommand{\arraystretch}{1.5}
\rowcolors{2}{gray!10}{white}
\begin{tabularx}{\textwidth}{|c|l|l|l|X|}
\hline
\rowcolor{blue!20}
\textbf{STT} & \textbf{Tên đối tượng} & \textbf{Kiểu} & \textbf{Ràng buộc} & \textbf{Chức năng} \\ \hline

1 & \textbf{Workflow Tip} & Container & - & Hướng dẫn quy trình xuất dữ liệu mẫu, chỉnh sửa Excel và nhập lại hệ thống. \\ \hline

2 & \textbf{Export Cards} & Card / Button & Theo từng loại dữ liệu & Tải xuống toàn bộ dữ liệu Thành phần, Quy tắc hoặc Sản phẩm dưới dạng tệp \texttt{.xlsx}. \\ \hline

3 & \textbf{Import Drop Zone} & File Input & Tệp \texttt{.xlsx} hoặc \texttt{.xls} & Nhận tệp bằng kéo thả hoặc chọn tệp, đồng thời hiển thị gợi ý tên cột yêu cầu. \\ \hline

4 & \textbf{Import Result Panel} & Container & Sau khi nhập & Hiển thị số dòng đã tạo, đã cập nhật, đã bỏ qua và danh sách lỗi có thể mở rộng. \\ \hline

5 & \textbf{Delete All Cards} & Card / Button & Yêu cầu xác nhận & Xóa vĩnh viễn toàn bộ dữ liệu của loại được chọn trong khu vực nguy hiểm. \\ \hline

\end{tabularx}
\end{center}

\textbf{c) Danh sách biến cố và xử lý}

\begin{center}
\renewcommand{\arraystretch}{1.5}
\rowcolors{2}{gray!10}{white}
\begin{tabularx}{\textwidth}{|c|l|X|}
\hline
\rowcolor{blue!20}
\textbf{STT} & \textbf{Biến cố} & \textbf{Xử lý} \\ \hline

1 & \textbf{Click "Tải xuống .xlsx"} & Gọi \texttt{GET /api/v1/admin/export/:entity}, nhận Blob từ server và tạo tệp tải xuống với tên chứa loại dữ liệu và thời điểm xuất. \\ \hline

2 & \textbf{Chọn hoặc thả tệp Excel} & Kiểm tra đuôi tệp; nếu hợp lệ, gửi \texttt{POST /api/v1/admin/import/:entity} dưới dạng \texttt{FormData}. Kết quả xử lý được hiển thị theo từng loại dữ liệu. \\ \hline

3 & \textbf{Tệp không hợp lệ / lỗi nhập} & Nếu tệp không thuộc định dạng Excel hoặc server trả lỗi, giao diện hiển thị thông báo lỗi; Admin có thể chọn tệp khác để thử lại. \\ \hline

4 & \textbf{Click "Xóa tất cả"} & Hiển thị cảnh báo không thể hoàn tác; khi xác nhận, gửi \texttt{DELETE /api/v1/admin/:entity/all} và thông báo kết quả. \\ \hline

\end{tabularx}
\end{center}


\subsubsection{Popup Chi tiết Thành phần}

\textbf{a) Giao diện}

Popup được mở trên màn hình Phân tích khi người dùng chọn một kết quả thành phần, phủ lớp nền mờ và tập trung thông tin chi tiết của thành phần đang chọn.

\textbf{b) Mô tả các đối tượng trên màn hình}

\begin{center}
\renewcommand{\arraystretch}{1.5}
\rowcolors{2}{gray!10}{white}
\begin{tabularx}{\textwidth}{|c|l|l|l|X|}
\hline
\rowcolor{blue!20}
\textbf{STT} & \textbf{Tên đối tượng} & \textbf{Kiểu} & \textbf{Ràng buộc} & \textbf{Chức năng} \\ \hline

1 & \textbf{Ingredient Name} & Label & Có kết quả được chọn & Hiển thị tên chuẩn hóa của thành phần. \\ \hline

2 & \textbf{Safety Badge} & Badge & Theo \texttt{effect} & Trình bày đánh giá Tốt, Xấu hoặc Trung bình bằng nhãn màu tương ứng. \\ \hline

3 & \textbf{Description} & Text & - & Hiển thị mô tả thành phần; dùng thông báo thay thế khi chưa có dữ liệu mô tả. \\ \hline

4 & \textbf{Recommended Action} & Text & Theo \texttt{effect} & Cung cấp lời khuyên sử dụng hoặc tránh thành phần phù hợp với kết quả đánh giá. \\ \hline

5 & \textbf{Report / Close Button} & Button & - & Mở popup Báo cáo phân loại hoặc đóng popup chi tiết. \\ \hline

\end{tabularx}
\end{center}

\textbf{c) Danh sách biến cố và xử lý}

\begin{center}
\renewcommand{\arraystretch}{1.5}
\rowcolors{2}{gray!10}{white}
\begin{tabularx}{\textwidth}{|c|l|X|}
\hline
\rowcolor{blue!20}
\textbf{STT} & \textbf{Biến cố} & \textbf{Xử lý} \\ \hline

1 & \textbf{Chọn thẻ kết quả} & Gán kết quả vào state \texttt{selectedIngredient}; popup hiển thị tên, mô tả, hiệu ứng và khuyến nghị tương ứng. \\ \hline

2 & \textbf{Click "Báo cáo"} & Lấy hiệu ứng hiện tại làm giá trị đề xuất ban đầu, đặt \texttt{showReportModal = true} để mở biểu mẫu báo cáo. \\ \hline

3 & \textbf{Đóng popup} & Click nút đóng hoặc vùng nền ngoài popup sẽ đặt \texttt{selectedIngredient = null} và quay lại kết quả phân tích. \\ \hline

\end{tabularx}
\end{center}


\subsubsection{Popup Xác nhận Thao tác}

\textbf{a) Giao diện}

Popup xác nhận được thể hiện trực tiếp trong màn hình Quản lý Người dùng; các thao tác xóa tại những màn hình quản trị khác sử dụng hộp xác nhận tương đương của trình duyệt.

\textbf{b) Mô tả các đối tượng trên màn hình}

\begin{center}
\renewcommand{\arraystretch}{1.5}
\rowcolors{2}{gray!10}{white}
\begin{tabularx}{\textwidth}{|c|l|l|l|X|}
\hline
\rowcolor{blue!20}
\textbf{STT} & \textbf{Tên đối tượng} & \textbf{Kiểu} & \textbf{Ràng buộc} & \textbf{Chức năng} \\ \hline

1 & \textbf{Confirmation Title} & Label & Theo thao tác & Hiển thị xác nhận xóa hoặc thay đổi trạng thái tài khoản. \\ \hline

2 & \textbf{Warning Message} & Text & - & Nêu tên người dùng bị tác động và cảnh báo khi thao tác xóa không thể hoàn tác. \\ \hline

3 & \textbf{Cancel Button} & Button & - & Đóng popup và không gửi yêu cầu thay đổi dữ liệu. \\ \hline

4 & \textbf{Confirm Button} & Button & Có thao tác đang chờ & Thực thi thao tác Khóa, Mở khóa hoặc Xóa theo lựa chọn trước đó. \\ \hline

\end{tabularx}
\end{center}

\textbf{c) Danh sách biến cố và xử lý}

\begin{center}
\renewcommand{\arraystretch}{1.5}
\rowcolors{2}{gray!10}{white}
\begin{tabularx}{\textwidth}{|c|l|X|}
\hline
\rowcolor{blue!20}
\textbf{STT} & \textbf{Biến cố} & \textbf{Xử lý} \\ \hline

1 & \textbf{Mở xác nhận} & Lưu loại thao tác, ID và username vào state \texttt{confirmModal}, sau đó hiển thị lớp phủ Modal. \\ \hline

2 & \textbf{Xác nhận trạng thái} & Với Khóa/Mở khóa, gọi \texttt{PATCH /api/v1/admin/users/:id/status}; với Xóa, gọi \texttt{DELETE /api/v1/admin/users/:id}. \\ \hline

3 & \textbf{Hoàn tất hoặc Hủy} & Sau khi thao tác thành công, tải lại danh sách người dùng và đóng Modal; nút Hủy chỉ đóng Modal mà không tác động dữ liệu. \\ \hline

\end{tabularx}
\end{center}


\subsubsection{Màn hình Cộng đồng -- Báo cáo}

\textbf{a) Giao diện}

Màn hình tại \texttt{/community/reports} được bảo vệ bởi \texttt{ProtectedRoute}, cho phép người dùng đăng nhập xem và bình chọn các báo cáo đang chờ xử lý.

\textbf{b) Mô tả các đối tượng trên màn hình}

\begin{center}
\renewcommand{\arraystretch}{1.5}
\rowcolors{2}{gray!10}{white}
\begin{tabularx}{\textwidth}{|c|l|l|l|X|}
\hline
\rowcolor{blue!20}
\textbf{STT} & \textbf{Tên đối tượng} & \textbf{Kiểu} & \textbf{Ràng buộc} & \textbf{Chức năng} \\ \hline

1 & \textbf{Sort Buttons} & Button & - & Chuyển thứ tự hiển thị giữa Nổi bật nhất và Mới nhất. \\ \hline

2 & \textbf{Report Card} & Card & Báo cáo trạng thái chờ & Hiển thị thành phần, loại da, hiệu ứng đề xuất, tác giả, ngày gửi, lý do và dẫn chứng. \\ \hline

3 & \textbf{Up / Down Vote} & Button & Người dùng đăng nhập & Bình chọn đồng ý hoặc phản đối và hiển thị điểm số của báo cáo. \\ \hline

4 & \textbf{Evidence Link} & Link & Khi có URL dẫn chứng & Mở tài liệu dẫn chứng được người gửi cung cấp trong tab mới. \\ \hline

5 & \textbf{Empty State} & Container & Không có báo cáo & Thông báo hiện tại không có báo cáo nào đang chờ xử lý. \\ \hline

\end{tabularx}
\end{center}

\textbf{c) Danh sách biến cố và xử lý}

\begin{center}
\renewcommand{\arraystretch}{1.5}
\rowcolors{2}{gray!10}{white}
\begin{tabularx}{\textwidth}{|c|l|X|}
\hline
\rowcolor{blue!20}
\textbf{STT} & \textbf{Biến cố} & \textbf{Xử lý} \\ \hline

1 & \textbf{Tải danh sách / đổi sắp xếp} & Gọi \texttt{GET /api/v1/reports/pending?sortBy=votes} hoặc \texttt{sortBy=newest}; lưu danh sách báo cáo và phiếu hiện tại của người dùng. \\ \hline

2 & \textbf{Click bình chọn} & Cập nhật giao diện tức thời theo hướng lạc quan, sau đó gửi \texttt{POST /api/v1/reports/vote} với \texttt{reportId} và \texttt{voteType}. \\ \hline

3 & \textbf{Bình chọn thất bại} & Nếu request thất bại, khôi phục trạng thái phiếu và điểm báo cáo trước khi người dùng thao tác. \\ \hline

\end{tabularx}
\end{center}


\subsubsection{Màn hình Admin -- Duyệt Báo cáo Cộng đồng}

\textbf{a) Giao diện}

Màn hình tại \texttt{/admin/community-reports} nằm trong layout quản trị, trình bày bảng báo cáo cộng đồng chờ duyệt và Modal duyệt hoặc từ chối.

\textbf{b) Mô tả các đối tượng trên màn hình}

\begin{center}
\renewcommand{\arraystretch}{1.5}
\rowcolors{2}{gray!10}{white}
\begin{tabularx}{\textwidth}{|c|l|l|l|X|}
\hline
\rowcolor{blue!20}
\textbf{STT} & \textbf{Tên đối tượng} & \textbf{Kiểu} & \textbf{Ràng buộc} & \textbf{Chức năng} \\ \hline

1 & \textbf{Pending Reports Table} & Table & Báo cáo chờ xử lý & Hiển thị thành phần, người báo cáo, đề xuất, lý do, điểm bình chọn và liên kết dẫn chứng. \\ \hline

2 & \textbf{Approve Button} & Button & Theo từng báo cáo & Mở xác nhận duyệt; khi hoàn tất sẽ cập nhật quy tắc phân loại theo đề xuất. \\ \hline

3 & \textbf{Reject Button} & Button & Theo từng báo cáo & Mở xác nhận từ chối báo cáo mà không thay đổi quy tắc hiện tại. \\ \hline

4 & \textbf{Admin Note Textarea} & Text Area & Tùy chọn & Nhập lý do duyệt hoặc từ chối để phản hồi cho người gửi. \\ \hline

5 & \textbf{Resolve Modal} & Popup & Chọn Duyệt/Từ chối & Hiển thị tác động của quyết định và yêu cầu xác nhận trước khi xử lý. \\ \hline

\end{tabularx}
\end{center}

\textbf{c) Danh sách biến cố và xử lý}

\begin{center}
\renewcommand{\arraystretch}{1.5}
\rowcolors{2}{gray!10}{white}
\begin{tabularx}{\textwidth}{|c|l|X|}
\hline
\rowcolor{blue!20}
\textbf{STT} & \textbf{Biến cố} & \textbf{Xử lý} \\ \hline

1 & \textbf{Truy cập màn hình} & Gọi \texttt{GET /api/v1/reports/pending?sortBy=votes} kèm Token và hiển thị danh sách báo cáo ưu tiên theo điểm cộng đồng. \\ \hline

2 & \textbf{Click "Duyệt" / "Từ chối"} & Lưu ID báo cáo và trạng thái \texttt{APPROVED} hoặc \texttt{REJECTED}, sau đó mở Modal nhập ghi chú. \\ \hline

3 & \textbf{Xác nhận xử lý} & Gửi \texttt{POST /api/v1/reports/resolve} gồm \texttt{reportId}, \texttt{status}, \texttt{adminNote}. Báo cáo đã xử lý được loại khỏi bảng; trường hợp duyệt sẽ tự động cập nhật quy tắc và tạo thông báo cho tác giả. \\ \hline

\end{tabularx}
\end{center}


\subsubsection{Popup Thông báo}

\textbf{a) Giao diện}

Dropdown Thông báo được gắn tại Navbar cho người dùng đã đăng nhập, hiển thị biểu tượng chuông, số thông báo chưa đọc và danh sách thông báo gần đây.

\textbf{b) Mô tả các đối tượng trên màn hình}

\begin{center}
\renewcommand{\arraystretch}{1.5}
\rowcolors{2}{gray!10}{white}
\begin{tabularx}{\textwidth}{|c|l|l|l|X|}
\hline
\rowcolor{blue!20}
\textbf{STT} & \textbf{Tên đối tượng} & \textbf{Kiểu} & \textbf{Ràng buộc} & \textbf{Chức năng} \\ \hline

1 & \textbf{Bell Button / Badge} & Button / Badge & Người dùng đăng nhập & Mở dropdown và hiển thị số thông báo chưa đọc, tối đa thể hiện là \texttt{9+}. \\ \hline

2 & \textbf{Notification List} & List & Tối đa 15 mục trên giao diện & Hiển thị tiêu đề, nội dung, thời gian tương đối và dấu hiệu chưa đọc của từng thông báo. \\ \hline

3 & \textbf{Mark All Read Button} & Button & Có thông báo chưa đọc & Đánh dấu toàn bộ thông báo của người dùng là đã đọc. \\ \hline

4 & \textbf{Notification Item} & Button & - & Đánh dấu mục được chọn là đã đọc và chuyển đến đường dẫn liên quan nếu thông báo có liên kết. \\ \hline

5 & \textbf{Empty State} & Container & Danh sách rỗng & Thông báo người dùng chưa có thông báo nào. \\ \hline

\end{tabularx}
\end{center}

\textbf{c) Danh sách biến cố và xử lý}

\begin{center}
\renewcommand{\arraystretch}{1.5}
\rowcolors{2}{gray!10}{white}
\begin{tabularx}{\textwidth}{|c|l|X|}
\hline
\rowcolor{blue!20}
\textbf{STT} & \textbf{Biến cố} & \textbf{Xử lý} \\ \hline

1 & \textbf{Tải thông báo} & Khi người dùng đăng nhập, gọi \texttt{GET /api/v1/notifications}; hệ thống tự tải lại mỗi 30 giây khi dropdown đang đóng và tải lại khi mở. \\ \hline

2 & \textbf{Click một thông báo} & Cập nhật trạng thái đã đọc trên giao diện, gửi \texttt{PATCH /api/v1/notifications/:id/read}; nếu có \texttt{link}, đóng dropdown và điều hướng tới liên kết. \\ \hline

3 & \textbf{Click "Đánh dấu đã đọc"} & Cập nhật tất cả mục sang trạng thái đã đọc và gửi \texttt{PATCH /api/v1/notifications/read-all}. \\ \hline

4 & \textbf{Click ngoài dropdown} & Trình lắng nghe sự kiện chuột phát hiện click ngoài vùng tham chiếu và đóng bảng thông báo. \\ \hline

\end{tabularx}
\end{center}


\subsubsection{Popup Báo cáo Phân loại}

\textbf{a) Giao diện}

Popup được mở từ nút Báo cáo trong Chi tiết Thành phần của màn hình Phân tích, cho phép người dùng đóng góp điều chỉnh kết quả phân loại.

\textbf{b) Mô tả các đối tượng trên màn hình}

\begin{center}
\renewcommand{\arraystretch}{1.5}
\rowcolors{2}{gray!10}{white}
\begin{tabularx}{\textwidth}{|c|l|l|l|X|}
\hline
\rowcolor{blue!20}
\textbf{STT} & \textbf{Tên đối tượng} & \textbf{Kiểu} & \textbf{Ràng buộc} & \textbf{Chức năng} \\ \hline

1 & \textbf{Effect Options} & Button Group & Chọn một giá trị & Chọn đề xuất đánh giá mới cho thành phần: Tốt, Trung bình hoặc Xấu theo loại da hiện tại. \\ \hline

2 & \textbf{Reason Textarea} & Text Area & Ít nhất 20 ký tự & Nhập lý do cho rằng đánh giá hiện tại chưa chính xác. \\ \hline

3 & \textbf{Evidence URL Input} & URL Input & Tùy chọn & Nhập liên kết tới bài nghiên cứu hoặc tài liệu dẫn chứng. \\ \hline

4 & \textbf{Submit Button} & Button & Lý do hợp lệ & Gửi báo cáo lên hệ thống; hiển thị trạng thái đang gửi trong khi xử lý. \\ \hline

5 & \textbf{Cancel / Close Button} & Button & - & Đóng popup mà không gửi báo cáo. \\ \hline

\end{tabularx}
\end{center}

\textbf{c) Danh sách biến cố và xử lý}

\begin{center}
\renewcommand{\arraystretch}{1.5}
\rowcolors{2}{gray!10}{white}
\begin{tabularx}{\textwidth}{|c|l|X|}
\hline
\rowcolor{blue!20}
\textbf{STT} & \textbf{Biến cố} & \textbf{Xử lý} \\ \hline

1 & \textbf{Mở popup} & Hệ thống lấy tên và hiệu ứng của \texttt{selectedIngredient}, mặc định chọn hiệu ứng hiện tại và hiển thị loại da của người dùng. \\ \hline

2 & \textbf{Click "Gửi báo cáo"} & Kiểm tra lý do tối thiểu 20 ký tự; gọi \texttt{GET /api/v1/ingredients/search?name=...} để lấy ID thành phần, sau đó gửi \texttt{POST /api/v1/reports} với hiệu ứng đề xuất, lý do và dẫn chứng. \\ \hline

3 & \textbf{Gửi thành công / thất bại} & Khi thành công, thông báo cảm ơn, đóng Modal và làm sạch trường nhập; nếu thất bại, hiển thị cảnh báo lỗi do server trả về. \\ \hline

\end{tabularx}
\end{center}

% \subsubsection{Trang chủ}

% a) Giao diện 
% (paste ảnh ở đây)

% b) Mô tả các đối tượng trên màn hình 
% 1 bảng gồm các cột STT, Tên, Kiểu, Ràng buộc, Chức năng

% c) Danh sách biến cố và xử lý

% 1 bảng gồm các cột STT, Biến cố, Xử lý
\end{document}




